\section{A uniform lower bound near slope one for the three-arrow quiver}\label{sec:central-lower}
\begin{definition}\label{def:central-euler-characteristic}
For positive coprime integers $a$ and $r$, let
\[
 E_r(a)=E_3(a,a+r).
\]
\end{definition}
From here to the end of the manuscript, let
\[
 \beta=\frac{256}{27}.
\]

\subsection{Graphs and quivers}\label{subsec:graphs-quivers}
\begin{definition}[Graphs]\label{def:graphs}
For a set $V$, define
\[
 \operatorname{Sym}^{2}(V)=(V\times V)/{\sim},\qquad
 (u,v)\sim(u',v')\ \Longleftrightarrow\
 (u',v')=(u,v)\ \text{or}\ (u',v')=(v,u).
\]
Let $\llbracket u,v\rrbracket$ denote the class of $(u,v)$.
For $U\subseteq V$, we identify $\operatorname{Sym}^{2}(U)$ with its image in $\operatorname{Sym}^{2}(V)$.
A \textit{graph} $G$ consists of a vertex set $\operatorname{Vert}(G)$,
an edge set $\operatorname{Edge}(G)$, and an endpoint map
\[
 \partial_G:\operatorname{Edge}(G)\longrightarrow
 \operatorname{Sym}^{2}(\operatorname{Vert}(G)).
\]
We call $G$ \textit{finite} if $\operatorname{Vert}(G)$ and $\operatorname{Edge}(G)$ are finite.
\begin{enumerate}[label=\arabic*.,ref=\arabic*,font=\itshape]
\item An edge $e$ is a \textit{loop at $v$} if $\partial_G(e)=\llbracket v,v\rrbracket$.
\item Edges $e$ and $e'$ are \textit{parallel} if $\partial_G(e)=\partial_G(e')$.
\item\label{d:graph-path-tuple} For $n\in\Z_{\ge0}$, consider a tuple
\[
 \gamma=(v_0,e_1,v_1,\ldots,e_n,v_n)
\]
such that $v_i\in\operatorname{Vert}(G)$ for $i\in[0,n]_{\Z}$ and
\[
 e_i\in\operatorname{Edge}(G),\qquad
 \partial_G(e_i)=\llbracket v_{i-1},v_i\rrbracket
 \qquad(i\in[n]_{\Z}).
\]
The tuple $\gamma$ starts at $v_0$ and ends at $v_n$.
We call $\gamma$ a \textit{simple path of length $n$} if $v_0,\ldots,v_n$ are pairwise distinct.
\item A \textit{cycle of length $n\in\Z_{\ge1}$} is a tuple
$(v_0,e_1,v_1,\ldots,e_n,v_n)$ such that $v_n=v_0$,
the vertices $v_0,\ldots,v_{n-1}\in\operatorname{Vert}(G)$ are pairwise distinct,
the edges $e_1,\ldots,e_n\in\operatorname{Edge}(G)$ are pairwise distinct, and
\[
 \partial_G(e_i)=\llbracket v_{i-1},v_i\rrbracket\qquad(i\in[n]_{\Z}).
\]
\end{enumerate}
\end{definition}

\begin{definition}[Graph notions]\label{def:graph-notions}
Let $G$ be a graph.
\begin{enumerate}[font=\itshape]
\item\label{d:graph-connected} We call $G$ \textit{connected} if, for every $\varnothing\ne U\subsetneq\operatorname{Vert}(G)$, there exist
$u\in U$, $v\in\operatorname{Vert}(G)\setminus U$, and $e\in\operatorname{Edge}(G)$ such that
\[
 \partial_G(e)=\llbracket u,v\rrbracket.
\]
\item We call $G$ \textit{simple} if it has no loops and $\partial_G$ is injective.
For a simple graph, we identify each edge with its image under $\partial_G$ and let
\[
 G=(\operatorname{Vert}(G),\operatorname{Edge}(G)),\qquad
 \operatorname{Edge}(G)\subseteq
 \{\llbracket u,v\rrbracket:u,v\in\operatorname{Vert}(G),\ u\ne v\}.
\]
We call a simple graph without cycles a \textit{forest} and a connected forest with at least one vertex a \textit{tree}.
\item Let $G'$ be a graph. A \textit{homomorphism} $\psi:G\to G'$ is a pair of maps
\[
 \psi_0:\operatorname{Vert}(G)\longrightarrow\operatorname{Vert}(G'),\qquad
 \psi_1:\operatorname{Edge}(G)\longrightarrow\operatorname{Edge}(G')
\]
such that, for every $e\in\operatorname{Edge}(G)$, $u\in\operatorname{Vert}(G)$, and $v\in\operatorname{Vert}(G)$,
\[
 \partial_G(e)=\llbracket u,v\rrbracket\quad\Longrightarrow\quad
 \partial_{G'}(\psi_1(e))=\llbracket \psi_0(u),\psi_0(v)\rrbracket.
\]
We call $\psi$ an \textit{isomorphism} if $\psi_0$ and $\psi_1$ are bijective.
We compose homomorphisms componentwise and let
$1_G=(1_{\operatorname{Vert}(G)},1_{\operatorname{Edge}(G)})$.
If $G$ and $G'$ are simple, we denote a homomorphism $\psi:G\to G'$ by its vertex map $\psi_0$.
\item Let $v\in\operatorname{Vert}(G)$. We define
\[
 \operatorname{AdjVert}_G(v)=\{u\in\operatorname{Vert}(G):
 \partial_G(e)=\llbracket v,u\rrbracket\text{ for some }e\in\operatorname{Edge}(G)\}.
\]
We call $u$ and $v$ \textit{adjacent} if $u\in\operatorname{AdjVert}_G(v)$.
We define the set of \textit{edges incident to $v$} by
\[
 \operatorname{IncEdge}_G(v)
 =\{e\in\operatorname{Edge}(G):\partial_G(e)=\llbracket v,u\rrbracket
 \text{ for some }u\in\operatorname{Vert}(G)\}.
\]
We define the \textit{degree of $v$} by
\[
\begin{aligned}
 \deg_G(v)
 &=\#\{e\in\operatorname{Edge}(G):\partial_G(e)=\llbracket v,u\rrbracket
       \text{ for some }u\ne v\}\\
 &\quad+2\#\{e\in\operatorname{Edge}(G):\partial_G(e)=\llbracket v,v\rrbracket\}.
\end{aligned}
\]
We define the set of \textit{leaves} by
\[
 \operatorname{Leaf}(G)=\{v\in\operatorname{Vert}(G):\deg_G(v)=1\}.
\]
\item\label{d:graph-subgraph} A graph $H$ is a \textit{subgraph} of $G$ if
\[
 \operatorname{Vert}(H)\subseteq\operatorname{Vert}(G),\qquad
 \operatorname{Edge}(H)\subseteq\operatorname{Edge}(G)
\]
and
\[
 \partial_H(e)=\partial_G(e)\qquad(e\in\operatorname{Edge}(H)).
\]
\item\label{d:graph-induced-subgraph} For $U\subseteq\operatorname{Vert}(G)$, the \textit{induced subgraph} $\operatorname{Subgraph}(G,U)$
has vertex set $U$ and edge set
\[
 \operatorname{Edge}(\operatorname{Subgraph}(G,U))
 =\{e\in\operatorname{Edge}(G):\partial_G(e)\in\operatorname{Sym}^{2}(U)\}.
\]
\item\label{d:graph-connected-components} Let $\operatorname{CtdComp}(G)$ denote the set of induced subgraphs
$\operatorname{Subgraph}(G,U)$ for nonempty $U\subseteq\operatorname{Vert}(G)$
that are maximal under inclusion among connected induced subgraphs of $G$.
We call an element of $\operatorname{CtdComp}(G)$ a \textit{connected component} of $G$.
\item\label{d:graph-edge-coloring} Let $C$ be a set.
An \textit{edge coloring} is a map $\kappa:\operatorname{Edge}(G)\to C$.
We call $(G,\kappa)$ a \textit{colored graph}.
\item A \textit{color-preserving homomorphism}
$(G,\kappa)\to(G',\kappa')$ is a graph homomorphism $\psi=(\psi_0,\psi_1):G\to G'$ such that
\[
 \kappa'\circ\psi_1=\kappa.
\]
We call it a \textit{color-preserving isomorphism} if $\psi_0$ and $\psi_1$ are bijective.
\item If $G$ is simple, we call $\kappa$ a \textit{proper edge coloring} if
\[
\begin{gathered}
 \llbracket u,v\rrbracket,\llbracket u,w\rrbracket\in\operatorname{Edge}(G),\qquad v\ne w\\
 \Longrightarrow\quad
 \kappa(\llbracket u,v\rrbracket)\ne\kappa(\llbracket u,w\rrbracket)\qquad(u,v,w\in\operatorname{Vert}(G)).
\end{gathered}
\]
\end{enumerate}
\end{definition}

\begin{definition}[Bipartite graphs]\label{def:bipartite-graphs}
Let $G$ be a graph. Let $I\subseteq\operatorname{Vert}(G)$ and $J\subseteq\operatorname{Vert}(G)$.
\begin{enumerate}[font=\itshape]
\item We call the pair $\mathbb G=(G,(I,J))$ a \textit{bipartite graph} if
\[
 \operatorname{Vert}(G)=I\sqcup J,\qquad
 \partial_G(\operatorname{Edge}(G))\subseteq
 \{\llbracket u,v\rrbracket:u\in I,\ v\in J\}.
\]
\item Let $\mathbb G=(G,(I,J))$ be a bipartite graph.
We call $I$ the \textit{source set} and $J$ the \textit{sink set} of $\mathbb G$.
Let $\operatorname{Src}_{\mathbb G}=I$ and $\operatorname{Snk}_{\mathbb G}=J$.
For $v\in\operatorname{Vert}(G)$, define
\[
\begin{aligned}
 \operatorname{SrcAdj}_{\mathbb G}(v)&=\operatorname{AdjVert}_G(v)\cap I,\\
 \operatorname{SnkAdj}_{\mathbb G}(v)&=\operatorname{AdjVert}_G(v)\cap J.
\end{aligned}
\]
\item Let $\mathbb G=(G,(I,J))$ be a bipartite graph. For a subgraph $H$ of $G$, define
\[
 \operatorname{SubSrc}_{\mathbb G}(H)=I\cap\operatorname{Vert}(H),\qquad
 \operatorname{SubSnk}_{\mathbb G}(H)=J\cap\operatorname{Vert}(H).
\]
\end{enumerate}
\end{definition}

\begin{lemma}[Rigidity of properly colored trees]\label{lem:colored-tree-rigidity}
Let $\mathbb T=(T,(I,J))$ be a bipartite graph such that $T$ is a finite tree.
Let $C$ be a set and let $\kappa:\operatorname{Edge}(T)\to C$ be a proper edge coloring.
Let $\phi=(\phi_0,\phi_1):T\to T$ be an automorphism satisfying
\[
 \kappa\circ\phi_1=\kappa,\qquad \phi_0(I)=I,\qquad \phi_0(J)=J.
\]
Then, $\phi=1_T$.
\end{lemma}
\begin{proof}
Let us use induction on $n=\#\operatorname{Vert}(T)$.
First, suppose $n=1$. The only vertex is fixed and there are no edges. Thus, $\phi=1_T$.

Second, suppose $n=2$. Let $I=\{u\}$ and $J=\{v\}$.
Since $\phi_0(I)=I$ and $\phi_0(J)=J$, we have $\phi_0(u)=u$ and $\phi_0(v)=v$.
The unique edge $\llbracket u,v\rrbracket$ is also fixed.

Third, suppose $n\ge3$. Let us remove the leaves and apply induction. Define
\[
 L=\operatorname{Leaf}(T),\qquad
 V'=\operatorname{Vert}(T)\setminus L,\qquad
 T'=\operatorname{Subgraph}(T,V').
\]
Let $(v_0,e_1,v_1,\ldots,e_\ell,v_\ell)$ be a longest simple path in $T$.
Then,
\[
 \ell\ge2,\qquad v_0,v_\ell\in L,\qquad
 \{v_1,\ldots,v_{\ell-1}\}\subseteq V'.
\]

Let us prove that $T'$ is connected.
For $u\in V'$ and $w\in V'$, let $P=(x_0,f_1,x_1,\ldots,f_r,x_r)$ be a simple path in $T$ from $x_0=u$ to $x_r=w$.
For $i\in[1,r-1]_{\Z}$, the distinct edges $f_i$ and $f_{i+1}$ are incident to $x_i$.
Hence, $\deg_T(x_i)\ge2$ and every vertex of $P$ belongs to $V'$.
Thus, $P$ lies in $T'$.
The graph $T'$ has no cycles. Therefore, $T'$ is a tree and
\[
 1\le\#\operatorname{Vert}(T')\le n-2.
\]

Let us apply the induction hypothesis to $T'$.
The automorphism $\phi$ preserves degrees. Hence, $\phi_0(L)=L$ and $\phi_0(V')=V'$.
Its restriction
\[
 \phi'=\bigl(\phi_0|_{V'},\phi_1|_{\operatorname{Edge}(T')}\bigr):T'\to T'
\]
preserves the coloring $\kappa|_{\operatorname{Edge}(T')}$ and the $(I\cap V',J\cap V')$-bipartition.
The induction hypothesis gives
\[
 \phi_0(v)=v\qquad(v\in V').
\]

Let us verify that $\phi_0(u)=u$ for each $u\in L$.
Its unique neighbor $v$ lies in $V'$ because $n\ge3$.
Let $e=\llbracket u,v\rrbracket$. Since $\phi_0(v)=v$,
\[
 e,\phi_1(e)\in\operatorname{IncEdge}_T(v),\qquad
 \kappa(\phi_1(e))=\kappa(e).
\]
Properness of $\kappa$ gives $\phi_1(e)=e$. Hence, $\phi_0(u)=u$.
Thus, $\phi_0=1_{\operatorname{Vert}(T)}$.
The injectivity of $\partial_T$ then gives $\phi_1=1_{\operatorname{Edge}(T)}$.
\end{proof}

\begin{definition}[Edge deletion and cut components]\label{def:cut-components}
Let $G$ be a graph.
\begin{enumerate}[font=\itshape]
\item\label{d:edge-deletion} For $e\in\operatorname{Edge}(G)$, we define a graph
\[
 \operatorname{DelEdge}(G,e)
 =\left(\operatorname{Vert}(G),\operatorname{Edge}(G)\setminus\{e\},
 \partial_G\big|_{\operatorname{Edge}(G)\setminus\{e\}}\right).
\]
\item\label{d:vertex-component} For $v\in\operatorname{Vert}(G)$, let $\operatorname{CtdComp}(G;v)$ be the connected component of $G$ containing $v$.
\item\label{d:sink-cut-component} Let $\mathbb S=(S,(I_S,J_S))$ be a bipartite graph such that $S$ is a tree.
Let $e=\llbracket x,b\rrbracket\in\operatorname{Edge}(S)$ for $x\in I_S$ and $b\in J_S$.
We define the \textit{sink-side cut component} by
\[
 \operatorname{SnkCtdComp}(\mathbb S;e)=\operatorname{CtdComp}(\operatorname{DelEdge}(S,e);b).
\]
\end{enumerate}
\end{definition}

\begin{definition}[Quivers]\label{def:quivers}
A \textit{quiver} is a triple
\[
 Q=(\operatorname{Vert}(Q),\operatorname{Edge}(Q),(s_Q,t_Q)).
\]
The sets $\operatorname{Vert}(Q)$ and $\operatorname{Edge}(Q)$ are its vertex and arrow sets.
The \textit{tail} and \textit{head} maps are
\[
 s_Q,t_Q:\operatorname{Edge}(Q)\longrightarrow\operatorname{Vert}(Q).
\]
We call $Q$ \textit{finite} if $\operatorname{Vert}(Q)$ and $\operatorname{Edge}(Q)$ are finite.
Let $\alpha:u\to v$ mean that $s_Q(\alpha)=u$ and $t_Q(\alpha)=v$.
If the map $\alpha\mapsto(s_Q(\alpha),t_Q(\alpha))$ is injective, we identify each arrow $\alpha$ with
$(s_Q(\alpha),t_Q(\alpha))$.
\end{definition}

\Needspace{6\baselineskip}
\begin{definition}[Directed paths and acyclicity]\label{def:directed-paths}
Let $Q$ be a quiver.
\begin{enumerate}[font=\itshape]
\item\label{d:quiver-directed-path} A \textit{directed path of length $n\in\Z_{\ge0}$} is a tuple
$(v_0,\alpha_1,v_1,\ldots,\alpha_n,v_n)$ such that
\[
 s_Q(\alpha_i)=v_{i-1},\qquad t_Q(\alpha_i)=v_i\qquad(i\in[n]_{\Z}).
\]
Here $v_i\in\operatorname{Vert}(Q)$ for $i\in[0,n]_{\Z}$ and
$\alpha_i\in\operatorname{Edge}(Q)$ for $i\in[n]_{\Z}$.
\item\label{d:quiver-directed-cycle} A directed path $(v_0,\alpha_1,v_1,\ldots,\alpha_n,v_n)$ is a \textit{directed cycle} if $n\ge1$, $v_n=v_0$, and $v_0,\ldots,v_{n-1}$ are pairwise distinct.
\item\label{d:quiver-acyclic} We call $Q$ \textit{acyclic} if it has no directed cycle.
\end{enumerate}
\end{definition}

\Needspace{6\baselineskip}
\begin{definition}[Quiver notions]\label{def:quiver-notions}
Let $Q$ be a quiver.
\begin{enumerate}[font=\itshape]
\item The \textit{underlying graph} $\operatorname{UnderG}(Q)$ of $Q$ has
\[
 \operatorname{Vert}(\operatorname{UnderG}(Q))=\operatorname{Vert}(Q),\qquad
 \operatorname{Edge}(\operatorname{UnderG}(Q))=\{\bar\alpha:\alpha\in\operatorname{Edge}(Q)\}.
\]
Here $\alpha\mapsto\bar\alpha$ is a bijection onto a copy of $\operatorname{Edge}(Q)$. The endpoint map of $\operatorname{UnderG}(Q)$ is
\[
 \partial_{\operatorname{UnderG}(Q)}(\bar\alpha)=\llbracket s_Q(\alpha),t_Q(\alpha)\rrbracket
 \qquad(\alpha\in\operatorname{Edge}(Q)).
\]
Let $\deg_Q(v)=\deg_{\operatorname{UnderG}(Q)}(v)$.
For $v\in\operatorname{Vert}(Q)$, we define the set of \textit{arrows incident to $v$} by
\[
\begin{aligned}
 \operatorname{IncEdgeQuiv}_Q(v)&:=\{\alpha\in\operatorname{Edge}(Q):\bar\alpha\in\operatorname{IncEdge}_{\operatorname{UnderG}(Q)}(v)\}\\
 &=\{\alpha\in\operatorname{Edge}(Q):s_Q(\alpha)=v\ \text{or}\ t_Q(\alpha)=v\}.
\end{aligned}
\]
We define $\operatorname{AdjVert}_Q(v)=\operatorname{AdjVert}_{\operatorname{UnderG}(Q)}(v)$.
We use $\operatorname{UnderG}(Q)$ to discuss trees and connected components of $Q$.
\item Let $Q$ and $Q'$ be quivers. A \textit{homomorphism} $\phi:Q\to Q'$ is a pair of maps
\[
 \phi_0:\operatorname{Vert}(Q)\longrightarrow\operatorname{Vert}(Q'),\qquad
 \phi_1:\operatorname{Edge}(Q)\longrightarrow\operatorname{Edge}(Q')
\]
satisfying
\[
 s_{Q'}\circ\phi_1=\phi_0\circ s_Q,\qquad
 t_{Q'}\circ\phi_1=\phi_0\circ t_Q.
\]
We call $\phi$ an \textit{isomorphism} if $\phi_0$ and $\phi_1$ are bijective.
We compose homomorphisms componentwise and let
$1_Q=(1_{\operatorname{Vert}(Q)},1_{\operatorname{Edge}(Q)})$.
The underlying graph homomorphism is
\[
 \operatorname{UnderG}(\phi)=(\phi_0,\bar\phi_1):
 \operatorname{UnderG}(Q)\longrightarrow\operatorname{UnderG}(Q'),\qquad
 \bar\phi_1(\bar\alpha)=\overline{\phi_1(\alpha)}.
\]
\item A quiver $Q'$ is a \textit{subquiver} of $Q$ if
\[
 \operatorname{Vert}(Q')\subseteq\operatorname{Vert}(Q),\qquad
 \operatorname{Edge}(Q')\subseteq\operatorname{Edge}(Q)
\]
and
\[
 s_{Q'}(\alpha)=s_Q(\alpha),\qquad t_{Q'}(\alpha)=t_Q(\alpha)
 \qquad(\alpha\in\operatorname{Edge}(Q')).
\]
For $U\subseteq\operatorname{Vert}(Q)$, the \textit{full subquiver} $\operatorname{SubQuiv}(Q,U)$
has vertex set $U$ and arrow set
\[
 \operatorname{Edge}(\operatorname{SubQuiv}(Q,U))
 =\{\alpha\in\operatorname{Edge}(Q):s_Q(\alpha),t_Q(\alpha)\in U\}.
\]
\end{enumerate}
\end{definition}

\begin{definition}[Bipartite quivers]\label{def:bipartite-quivers}
Let $Q$ be a quiver. Let $I\subseteq\operatorname{Vert}(Q)$ and $J\subseteq\operatorname{Vert}(Q)$.
\begin{enumerate}[font=\itshape]
\item We call the pair $\mathbb Q=(Q,(I,J))$ a \textit{bipartite quiver} if
\[
 \operatorname{Vert}(Q)=I\sqcup J,\qquad
 (s_Q,t_Q)(\operatorname{Edge}(Q))\subseteq I\times J.
\]
The underlying bipartite graph of $\mathbb Q$ is $(\operatorname{UnderG}(Q),(I,J))$.
\item Let $\mathbb Q=(Q,(I,J))$ and $\mathbb Q'=(Q',(I',J'))$ be bipartite quivers.
A quiver isomorphism $(\phi_0,\phi_1):Q\to Q'$ \textit{preserves the bipartitions} if
\[
 \phi_0(I)=I',\qquad \phi_0(J)=J'.
\]
\item Let $\mathbb Q=(Q,(I,J))$ be a bipartite quiver. We call $\mathbb Q$ \textit{simple} if the map
\[
 (s_Q,t_Q):\operatorname{Edge}(Q)\longrightarrow I\times J
\]
is injective.
\end{enumerate}
\end{definition}

See Remark~4.7(1) of~\cite{Weist} for arrow colors of bipartite quivers.

\begin{definition}[Colored quivers]\label{def:colored-quivers}
Let $Q$ be a quiver. Let $C$ be a set.
\begin{enumerate}[font=\itshape]
\item A \textit{colored quiver} with color set $C$ is a pair
\[
 \widetilde Q=(Q,c),\qquad c:\operatorname{Edge}(Q)\longrightarrow C.
\]
We define $\operatorname{QCol}(\widetilde Q)=c$ and use the notation
\[
\begin{gathered}
 \operatorname{Vert}(\widetilde Q)=\operatorname{Vert}(Q),\qquad
 \operatorname{Edge}(\widetilde Q)=\operatorname{Edge}(Q),\\
 s_{\widetilde Q}=s_Q,\qquad t_{\widetilde Q}=t_Q,\qquad
 \operatorname{UnderG}(\widetilde Q)=\operatorname{UnderG}(Q).
\end{gathered}
\]
Let
\[
 \deg_{\widetilde Q}(v)=\deg_Q(v),\qquad
 \operatorname{IncEdgeQuiv}_{\widetilde Q}(v)=\operatorname{IncEdgeQuiv}_Q(v)
 \qquad(v\in\operatorname{Vert}(Q)).
\]
For $v\in\operatorname{Vert}(Q)$, define
\[
 \operatorname{AdjVert}_{\widetilde Q}(v)=\operatorname{AdjVert}_Q(v).
\]
\item\label{d:colored-arrow} Let $\widetilde Q=(Q,c)$ be a colored quiver. Let $v\in\operatorname{Vert}(Q)$ and $d\in C$.
Define the set of \textit{$(v,d)$-colored arrows} by
\[
 \operatorname{ColArrow}_{\widetilde Q}(v,d)
 :=\{\alpha\in\operatorname{Edge}(Q):t_Q(\alpha)=v,\ c(\alpha)=d\}.
\]
\item\label{d:arrow-operators} Let $\widetilde Q=(Q,c)$ be a colored quiver.
\begin{enumerate}[label=(\alph*),ref=\alph*,font=\itshape]
\item\label{d:delete-arrow} For $\alpha\in\operatorname{Edge}(Q)$, we define
$\operatorname{DelArrow}(\widetilde Q,\alpha)$ by keeping $\operatorname{Vert}(Q)$,
replacing the arrow set by $\operatorname{Edge}(Q)\setminus\{\alpha\}$,
and restricting $s_Q$, $t_Q$, and $c$ to this set.
\item\label{d:add-arrow} For $x\in\operatorname{Vert}(Q)$, $v\in\operatorname{Vert}(Q)$, and $d\in C$, we define
$\operatorname{AddArrow}(\widetilde Q;x,v,d)$ by keeping $\operatorname{Vert}(Q)$
and adjoining a new arrow $\alpha_*$ to $\operatorname{Edge}(Q)$.
The source, target, and coloring maps extend those of $\widetilde Q$ by
\[
 s(\alpha_*)=x,\qquad t(\alpha_*)=v,\qquad c(\alpha_*)=d.
\]
\end{enumerate}
\item Let $\widetilde Q=(Q,c)$ and $\widetilde Q'=(Q',c')$ be colored quivers with color set $C$.
A \textit{color-preserving homomorphism} $\widetilde Q\to\widetilde Q'$
is a quiver homomorphism $(\phi_0,\phi_1):Q\to Q'$ such that
\[
 \operatorname{QCol}(\widetilde Q')\circ\phi_1
 =\operatorname{QCol}(\widetilde Q).
\]
We call this homomorphism a \textit{color-preserving isomorphism}
if $\phi_0$ and $\phi_1$ are bijective.
\item Let $\widetilde Q=(Q,c)$ and $\widetilde Q'=(Q',c')$ be colored quivers with color set $C$.
We call $\widetilde Q'$ a \textit{colored subquiver} of $\widetilde Q$ if $Q'$ is a subquiver of $Q$ and
\[
 c'=c\big|_{\operatorname{Edge}(Q')}.
\]
For $U\subseteq\operatorname{Vert}(Q)$, we define the \textit{full colored subquiver} on $U$ by
\[
 \operatorname{SubColQuiv}(\widetilde Q,U)=\left(\operatorname{SubQuiv}(Q,U),c\big|_{\operatorname{Edge}(\operatorname{SubQuiv}(Q,U))}\right).
\]
\item Let $\widetilde Q=(Q,c)$ be a colored quiver. Let $I\subseteq\operatorname{Vert}(Q)$ and $J\subseteq\operatorname{Vert}(Q)$.
We call the pair $\widetilde{\mathbb Q}=(\widetilde Q,(I,J))$ a \textit{colored bipartite quiver} if $\mathbb Q=(Q,(I,J))$ is a bipartite quiver.
We use
\[
\begin{gathered}
 \operatorname{Vert}(\widetilde{\mathbb Q})=\operatorname{Vert}(Q),\qquad
 \operatorname{Edge}(\widetilde{\mathbb Q})=\operatorname{Edge}(Q),\\
 s_{\widetilde{\mathbb Q}}=s_Q,\qquad t_{\widetilde{\mathbb Q}}=t_Q,\qquad
 \operatorname{QCol}(\widetilde{\mathbb Q})=c,\\
 \operatorname{UnderG}(\widetilde{\mathbb Q})=\operatorname{UnderG}(Q),\qquad
 \deg_{\widetilde{\mathbb Q}}(v)=\deg_Q(v)\quad(v\in\operatorname{Vert}(Q)).
\end{gathered}
\]
\item Let $\widetilde{\mathbb Q}=(\widetilde Q,(I,J))$ be a colored bipartite quiver.
We call $\widetilde{\mathbb Q}$ \textit{properly colored} if the restriction
\[
 c\big|_{\operatorname{IncEdgeQuiv}_Q(v)}:\operatorname{IncEdgeQuiv}_Q(v)\longrightarrow C
\]
is injective for every $v\in\operatorname{Vert}(Q)$.
\item\label{d:vacant-slot} Let $\widetilde{\mathbb Q}=(\widetilde Q,(I,J))$ be a colored bipartite quiver.
We define its set of \textit{colored-vacant slots} by
\[
 \operatorname{ColVacSlot}(\widetilde{\mathbb Q})
 :=\{(v,d)\in J\times C:
 \operatorname{ColArrow}_{\widetilde Q}(v,d)=\varnothing\}.
\]
We call $d$ the color of the colored-vacant slot $(v,d)$.
\end{enumerate}
\end{definition}

\subsection{Representations}\label{subsec:quiver-representations}
\begin{definition}[Representations of quivers]\label{def:quiver-representations}
Let $Q$ be a quiver.
\begin{enumerate}[font=\itshape]
\item A \textit{representation of $Q$ over $\C$} is a family
\[
 R=\bigl((R_v)_{v\in\operatorname{Vert}(Q)},
 (R_\alpha)_{\alpha\in\operatorname{Edge}(Q)}\bigr)
\]
such that each $R_v$ is a finite-dimensional complex vector space and each map
\[
 R_\alpha:R_{s_Q(\alpha)}\longrightarrow R_{t_Q(\alpha)}
 \qquad(\alpha\in\operatorname{Edge}(Q))
\]
is linear.
\item\label{d:representation-hom} Let $Q'$ be a quiver. Let $R$ and $R'$ be representations of $Q$ and $Q'$, respectively.
For a quiver homomorphism $\phi=(\phi_0,\phi_1):Q\to Q'$, a \textit{homomorphism}
$f:R\to R'$ \textit{over $\phi$} is a tuple $f=(f_v)_{v\in\operatorname{Vert}(Q)}$ of complex linear maps
$f_v:R_v\to R'_{\phi_0(v)}$ such that
\[
 f_{t_Q(\alpha)}\circ R_\alpha
 =R'_{\phi_1(\alpha)}\circ f_{s_Q(\alpha)}
 \qquad(\alpha\in\operatorname{Edge}(Q)).
\]
\item For a representation $R''$ of a quiver $Q''$ and a homomorphism $g:R'\to R''$
over $\psi:Q'\to Q''$, we define $g\circ f$ over $\psi\circ\phi$ by
\[
 (g\circ f)_v=g_{\phi_0(v)}\circ f_v\qquad(v\in\operatorname{Vert}(Q)).
\]
\item Let $\operatorname{Rep}_{\C}(Q)$ be the category of representations of $Q$ over $\C$ whose morphisms are homomorphisms over $1_Q$.
\item For $R\in\operatorname{Rep}_{\C}(Q)$, we define its \textit{dimension vector} and \textit{support} by
\[
\begin{aligned}
 \dim_{\C}(R)&=\bigl(\dim_{\C}(R_i)\bigr)_{i\in\operatorname{Vert}(Q)},\\
 \operatorname{RepSupp}(R)&=\{i\in\operatorname{Vert}(Q):R_i\ne0\}.
\end{aligned}
\]
The \textit{support quiver} is $\operatorname{SubQuiv}(Q,\operatorname{RepSupp}(R))$.
\item Let $R\in\operatorname{Rep}_{\C}(Q)$ and $R'\in\operatorname{Rep}_{\C}(Q)$.
An \textit{isomorphism} $\phi:R\to R'$ in $\operatorname{Rep}_{\C}(Q)$ is a tuple
\[
 \phi=(\phi_v)_{v\in\operatorname{Vert}(Q)}
\]
such that each $\phi_v:R_v\to R'_v$ is a complex linear isomorphism and
\[
 \phi_{t_Q(\alpha)}\circ R_\alpha
 =R'_\alpha\circ\phi_{s_Q(\alpha)}
 \qquad(\alpha\in\operatorname{Edge}(Q)).
\]
\item Let $R\in\operatorname{Rep}_{\C}(Q)$.
A \textit{subrepresentation} $W$ of $R$ consists of complex linear subspaces
$W_v\subseteq R_v$ for $v\in\operatorname{Vert}(Q)$ such that
\[
 R_\alpha(W_{s_Q(\alpha)})\subseteq W_{t_Q(\alpha)}
 \qquad(\alpha\in\operatorname{Edge}(Q)).
\]
Its arrow maps are
\[
 W_\alpha=R_\alpha\big|_{W_{s_Q(\alpha)}}:
 W_{s_Q(\alpha)}\longrightarrow W_{t_Q(\alpha)}.
\]
We denote the set of these subrepresentations by $\operatorname{SubRep}(R)$.
The zero subrepresentation has $W_v=0$ for every $v\in\operatorname{Vert}(Q)$.
We call $W$ \textit{proper} if $W\ne R$.
\end{enumerate}
\end{definition}

\Needspace{9\baselineskip}
\begin{definition}[Restriction and extension by zero]\label{def:rep-restriction-zero}
Let $Q$ be a quiver, let $U\subseteq\operatorname{Vert}(Q)$. Let $Q_U=\operatorname{SubQuiv}(Q,U)$.
\begin{enumerate}[font=\itshape]
\item Define
$\operatorname{Res}_{Q,U}:\operatorname{Rep}_{\C}(Q)\to\operatorname{Rep}_{\C}(Q_U)$
by keeping the spaces $R_v$ for $v\in U$, the maps $R_\alpha$ for $\alpha\in\operatorname{Edge}(Q_U)$, and the components $f_v$ for $v\in U$ of each homomorphism $f$.
\item Define
$\operatorname{ExtZero}_{Q,U}:\operatorname{Rep}_{\C}(Q_U)\to\operatorname{Rep}_{\C}(Q)$
by
\[
 (\operatorname{ExtZero}_{Q,U}R)_v=
 \begin{cases}R_v&v\in U,\\0&v\notin U,\end{cases}
 \qquad
 (\operatorname{ExtZero}_{Q,U}R)_\alpha=
 \begin{cases}R_\alpha&s_Q(\alpha),t_Q(\alpha)\in U,\\0&\text{otherwise}.\end{cases}
\]
For a homomorphism $f$, its component is $f_v$ on $U$ and the zero map outside $U$.
\end{enumerate}
\end{definition}

\begin{lemma}\label{lem:rep-restriction-zero}
Let $Q$ be a quiver, let $U\subseteq\operatorname{Vert}(Q)$. Let $Q_U=\operatorname{SubQuiv}(Q,U)$.
Then, we have the following.

\begin{enumerate}
\item\label{c:restriction-zero-left} Let $R\in\operatorname{Rep}_{\C}(Q_U)$.
The maps
\[
 (\varepsilon_R)_v=1_{R_v}\qquad(v\in U)
\]
define an isomorphism
\[
 \varepsilon_R:\operatorname{Res}_{Q,U}(\operatorname{ExtZero}_{Q,U}R)\longrightarrow R.
\]
\item\label{c:restriction-zero-right} Let $S\in\operatorname{Rep}_{\C}(Q)$ satisfy $\operatorname{RepSupp}(S)\subseteq U$.
The maps
\[
 (\eta_S)_v=
 \begin{cases}1_{S_v}&v\in U,\\0&v\notin U\end{cases}
\]
define an isomorphism
\[
 \eta_S:\operatorname{ExtZero}_{Q,U}(\operatorname{Res}_{Q,U}S)\longrightarrow S.
\]
\item\label{c:restriction-zero-subrep} Let $S$ satisfy the hypotheses of Claim~\ref{c:restriction-zero-right}.
The map
\[
 \rho_S:\operatorname{SubRep}(S)\longrightarrow
 \operatorname{SubRep}(\operatorname{Res}_{Q,U}S),\qquad
 W\longmapsto\operatorname{Res}_{Q,U}W
\]
is a bijection. Its inverse is
\[
 \sigma_S:\operatorname{SubRep}(\operatorname{Res}_{Q,U}S)\longrightarrow
 \operatorname{SubRep}(S),\qquad Z\longmapsto\widehat Z.
\]
For $Z\in\operatorname{SubRep}(\operatorname{Res}_{Q,U}S)$, define
\[
 \widehat Z_v=
 \begin{cases}Z_v&v\in U,\\0&v\notin U\end{cases}
 \qquad(v\in\operatorname{Vert}(Q)).
\]
For $W\in\operatorname{SubRep}(S)$ and $v\in U$,
\[
 \dim_{\C}(\rho_S(W)_v)=\dim_{\C}(W_v).
\]
\end{enumerate}
\end{lemma}
\begin{proof}
\prf{Claim~\ref{c:restriction-zero-left}} Let us compare the vertex spaces and arrow maps after restriction.
By Definition~\ref{def:rep-restriction-zero},
\[
\begin{aligned}
 (\operatorname{Res}_{Q,U}(\operatorname{ExtZero}_{Q,U}R))_v&=R_v
 &&(v\in U),\\
 (\operatorname{Res}_{Q,U}(\operatorname{ExtZero}_{Q,U}R))_\alpha&=R_\alpha
 &&(\alpha\in\operatorname{Edge}(Q_U)).
\end{aligned}
\]
Thus, the identity maps define $\varepsilon_R$ and its inverse.

\prf{Claim~\ref{c:restriction-zero-right}} Let us check the subrepresentation at the vertices outside $U$.
For $v\notin U$, we have $S_v=0$.
The map $(\eta_S)_v$ is therefore the unique isomorphism between zero vector spaces.
For $v\in U$, it is the identity on $S_v$.
Let $\widehat S=\operatorname{ExtZero}_{Q,U}(\operatorname{Res}_{Q,U}S)$.
For each $\alpha\in\operatorname{Edge}(Q)$, the equation
\[
 (\eta_S)_{t_Q(\alpha)}\circ\widehat S_\alpha
 =S_\alpha\circ(\eta_S)_{s_Q(\alpha)}
\]
holds as an equality of maps $\widehat S_{s_Q(\alpha)}\to S_{t_Q(\alpha)}$.
If $s_Q(\alpha)\in U$ and $t_Q(\alpha)\in U$, both sides equal $S_\alpha$. Otherwise, both sides are zero.
Hence, $\eta_S$ is an isomorphism of representations.

\prf{Claim~\ref{c:restriction-zero-subrep}} Let us verify $\rho_S(\widehat Z)=Z$ and $\widehat{\rho_S(W)}=W$.
Let $Z\in\operatorname{SubRep}(\operatorname{Res}_{Q,U}S)$ and $\alpha\in\operatorname{Edge}(Q)$.
If both endpoints of $\alpha$ lie in $U$, the inclusion
$S_\alpha(\widehat Z_{s_Q(\alpha)})\subseteq\widehat Z_{t_Q(\alpha)}$
is the subrepresentation condition for $Z$.
Otherwise, $S_{s_Q(\alpha)}=0$ or $S_{t_Q(\alpha)}=0$.
Then, $S_\alpha(\widehat Z_{s_Q(\alpha)})=0\subseteq\widehat Z_{t_Q(\alpha)}$.
For every $W\in\operatorname{SubRep}(S)$, we have $W_v=0$ for $v\notin U$.
Consequently,
\[
 \sigma_S\circ\rho_S=1_{\operatorname{SubRep}(S)},\qquad
 \rho_S\circ\sigma_S=1_{\operatorname{SubRep}(\operatorname{Res}_{Q,U}S)}.
\]
The dimension equality follows from $\rho_S(W)_v=W_v$ for $v\in U$.
\end{proof}

\begin{definition}[Euler and Tits forms]\label{def:euler-tits-forms}
Let $H$ be a finite acyclic quiver. For
\[
 d=(d_v)_{v\in\operatorname{Vert}(H)}\in\Z^{\operatorname{Vert}(H)},\qquad
 \delta=(\delta_v)_{v\in\operatorname{Vert}(H)}\in\Z^{\operatorname{Vert}(H)},
\]
we define the \textit{Euler form} and \textit{Tits quadratic form}
\[
\langle d,\delta\rangle_H=\sum_{v\in\operatorname{Vert}(H)}d_v\delta_v
 -\sum_{\alpha\in\operatorname{Edge}(H)}d_{s_H(\alpha)}\delta_{t_H(\alpha)},\qquad Q_H(d)=\langle d,d\rangle_H.
\]
\end{definition}

\begin{definition}[Representations of colored quivers]\label{def:colored-quiver-representations}
Let $C$ be a set, let $Q$ be a quiver, and let $\widetilde Q=(Q,c)$ be a colored quiver with color set $C$.
\begin{enumerate}[font=\itshape]
\item We define a \textit{representation $R$ of $\widetilde Q$ over $\C$} as a representation of $Q$ over $\C$.
\item Let $R$ be a representation of $\widetilde Q$ over $\C$. We define its \textit{colored support quiver} by
\[
 \operatorname{SuppColQuiv}(R):=\operatorname{SubColQuiv}(\widetilde Q,\operatorname{RepSupp}(R)).
\]
\item\begin{enumerate}[label=(\alph*),ref=\alph*,font=\itshape]
\item Let $\widetilde Q'=(Q',c')$ be a colored quiver with color set $C$.
Let $R$ and $R'$ be representations of $\widetilde Q$ and $\widetilde Q'$, respectively.
A \textit{color-preserving homomorphism} $\Phi:R\to R'$ is a tuple
\[
 \Phi=\bigl(\phi_0,\phi_1,(f_v)_{v\in\operatorname{Vert}(Q)}\bigr)
\]
such that $(\phi_0,\phi_1):\widetilde Q\to\widetilde Q'$ is a color-preserving quiver homomorphism
and the tuple $(f_v)_{v\in\operatorname{Vert}(Q)}$ is a homomorphism of representations over $(\phi_0,\phi_1)$.
\item For a composable color-preserving homomorphism
$\Psi=(\psi_0,\psi_1,(g_w)_{w\in\operatorname{Vert}(Q')})$, we define
\[
 \Psi\circ\Phi=
 \bigl(\psi_0\circ\phi_0,\psi_1\circ\phi_1,
 (g_{\phi_0(v)}\circ f_v)_{v\in\operatorname{Vert}(Q)}\bigr).
\]
The identity of $R$ is $(1_{\operatorname{Vert}(Q)},1_{\operatorname{Edge}(Q)},(1_{R_v})_{v\in\operatorname{Vert}(Q)})$.
\item Let $\operatorname{ColRep}_{\C}(\widetilde Q)$ denote the category of representations of $\widetilde Q$ and color-preserving homomorphisms.
For $R\in\operatorname{ColRep}_{\C}(\widetilde Q)$ and $R'\in\operatorname{ColRep}_{\C}(\widetilde Q)$, let
\[
 \operatorname{Hom}_{\operatorname{ColRep}_{\C}(\widetilde Q)}(R,R')
\]
denote the set of color-preserving homomorphisms $R\to R'$.
\item We call $\Phi$ a \textit{color-preserving isomorphism} if $\phi_0$ and $\phi_1$ are bijective
and every $f_v$ is a complex linear isomorphism.
\end{enumerate}
\item Let $R\in\operatorname{ColRep}_{\C}(\widetilde Q)$.
A \textit{colored subrepresentation} $W$ of $R$ is a representation of $\widetilde Q$
such that $W\in\operatorname{SubRep}(R)$.
We denote the set of these colored subrepresentations by $\operatorname{SubColRep}(R)$.
\item\label{d:total-dimensions} Let $(Q,(I,J))$ be a bipartite quiver. Let $R\in\operatorname{Rep}_{\C}(Q)$ have finite support $\operatorname{RepSupp}(R)$.
We define the \textit{total source dimension} and \textit{total sink dimension} of $R$ by
\[
 \operatorname{SrcDim}_{\C,(I,J)}(R)=\sum_{v\in I}\dim_{\C}R_v,\qquad
 \operatorname{SnkDim}_{\C,(I,J)}(R)=\sum_{v\in J}\dim_{\C}R_v.
\]
We define its \textit{total dimension pair} by
\[
 \operatorname{TotDim}_{\C,(I,J)}(R):=\bigl(\operatorname{SrcDim}_{\C,(I,J)}(R),\operatorname{SnkDim}_{\C,(I,J)}(R)\bigr).
\]
\end{enumerate}
\end{definition}

\begin{definition}[Bipartite stability]\label{def:bipartite-stability}
Let $\mathbb Q=(Q,(I,J))$ be a bipartite quiver.
Let $R\in\operatorname{Rep}_{\C}(Q)$ be nonzero and have finite support.
Let $(a,b)=\operatorname{TotDim}_{\C,(I,J)}(R)$.
We call $R$ \textit{stable} if every nonzero proper $W\in\operatorname{SubRep}(R)$ satisfies
\[
 a\operatorname{SnkDim}_{\C,(I,J)}(W)-b\operatorname{SrcDim}_{\C,(I,J)}(W)>0.
\]
\end{definition}

\begin{definition}[Arrow operators on representations]\label{def:representation-arrow-operators}
Let $\widetilde Q=(Q,c)$ be a colored quiver with color set $C$. Let $R\in\operatorname{Rep}_{\C}(Q)$.
\begin{enumerate}[font=\itshape]
\item\label{d:representation-delete-arrow} For $\alpha\in\operatorname{Edge}(Q)$, we define
$\operatorname{DelArrow}(R,\alpha)$ on $\operatorname{DelArrow}(\widetilde Q,\alpha)$
by keeping all vertex spaces and all arrow maps except $R_\alpha$.
\item\label{d:representation-add-arrow} For $x\in\operatorname{Vert}(Q)$, $v\in\operatorname{Vert}(Q)$, $d\in C$, and a linear map $f:R_x\to R_v$,
we define
\[
 \operatorname{AddArrow}(R;x,v,d,f)
\]
on $\operatorname{AddArrow}(\widetilde Q;x,v,d)$ by keeping the vertex spaces and
existing arrow maps and setting
\[
 (\operatorname{AddArrow}(R;x,v,d,f))_{\alpha_*}=f.
\]
\end{enumerate}
\end{definition}




\begin{definition}[Isomorphisms and markings of colored representations]\label{def:colored-representation-markings}
Let $C$ be a set. Let $\widetilde Q=(Q,c)$ and $\widetilde Q'=(Q',c')$ be colored quivers with color set $C$.
Let $\widetilde{\mathbb Q}=(\widetilde Q,(I,J))$ and $\widetilde{\mathbb Q}'=(\widetilde Q',(I',J'))$ be colored bipartite quivers such that $Q$ and $Q'$ are finite.
Let $R\in\operatorname{ColRep}_{\C}(\widetilde Q)$ and $R'\in\operatorname{ColRep}_{\C}(\widetilde Q')$.
\begin{enumerate}[font=\itshape]
\item\label{d:unmarked-colored-rep-iso} An isomorphism of $(\widetilde{\mathbb Q},R)$ and $(\widetilde{\mathbb Q}',R')$ is a color-preserving isomorphism $R\to R'$ that takes $(I,J)$ to $(I',J')$.
We denote the isomorphism class of $(\widetilde{\mathbb Q},R)$ by
$\operatorname{CIsomClass}_{\C}(\widetilde{\mathbb Q},R)$.
\item\label{d:ordered-colored-marking} Let $k\in\Z_{\ge0}$. An \textit{ordered marking of colored-vacant slots} of $R$ is a tuple
\[
 \mathbf s=((v_i,d_i))_{i\in[k]_{\Z}},\qquad
 (v_i,d_i)\in\operatorname{ColVacSlot}(\widetilde{\mathbb Q}).
\]
The marked representation is the pair $(R;\mathbf s)$.
Distinct marked colored-vacant slots means $(v_i,d_i)\ne(v_j,d_j)$ for $i\ne j$.
\item\label{d:marked-rep-category} Fix $\widetilde{\mathbb Q}$ and $\mathbf s$. We define the category
$\operatorname{MarkedColRep}_{\C}(\widetilde{\mathbb Q};\mathbf s)$ as follows.
Its objects are the pairs $(R;\mathbf s)$ with $R\in\operatorname{ColRep}_{\C}(\widetilde Q)$.
Its homomorphisms $(R;\mathbf s)\to(R';\mathbf s)$ are the homomorphisms
$\Phi=(\phi_0,\phi_1,(f_v)_v)\in\operatorname{Hom}_{\operatorname{ColRep}_{\C}(\widetilde Q)}(R,R')$
such that
\[
\begin{gathered}
 \phi_0(I)\subseteq I,\qquad
 \phi_0(J)\subseteq J,\\
 \phi_0(v_i)=v_i\qquad(i\in[k]_{\Z}).
\end{gathered}
\]
We use the identities and composition of $\operatorname{ColRep}_{\C}(\widetilde Q)$.
\item\label{d:marked-rep-iso} Let $\mathbf s'=((v'_i,d'_i))_{i\in[k]_{\Z}}$ be an ordered marking of $R'$.
An isomorphism $(\widetilde{\mathbb Q},R;\mathbf s)\to(\widetilde{\mathbb Q}',R';\mathbf s')$ is an isomorphism
$\Phi=(\phi_0,\phi_1,(f_v)_{v\in\operatorname{Vert}(Q)})$
of $(\widetilde{\mathbb Q},R)$ and $(\widetilde{\mathbb Q}',R')$ satisfying
\[
 (\phi_0(v_i),d_i)=(v'_i,d'_i)\qquad(i\in[k]_{\Z}).
\]
We denote its \textit{marked color-preserving isomorphism class} by
$\operatorname{MarkedCIsom}_{\C}(\widetilde{\mathbb Q},R;\mathbf s)$.
\end{enumerate}
\end{definition}

\subsection{Stable zero-sum representations and gluing}
From here to the end of this section, all colored quivers have color set $[0,2]_{\Z}$.

\begin{definition}[Zero-sum representations]\label{def:reflected-block}
Let $\mathbb U=(Q,(I_Q,J_Q))$ be a bipartite quiver such that $Q$ is finite.
We define the \textit{zero-sum representation} $\operatorname{ZS}(\mathbb U)$ of $Q$ as follows.
\begin{enumerate}[font=\itshape]
\item\label{d:reflected-kernel} For $u\in J_{Q}$, let $\operatorname{ZS}(\mathbb U)_u=\C$.
For $v\in I_Q$, put $A_v=\operatorname{IncEdgeQuiv}_Q(v)$.
On the space $\C^{A_v}$ of tuples $z=(z_\alpha)_{\alpha\in A_v}$, define
\begin{equation}\label{eq:zero-sum-map}
 \Sigma_v\colon\C^{A_v}\longrightarrow\C,\qquad
 \Sigma_v(z)=\sum_{\alpha\in A_v}z_\alpha,
\end{equation}
and let $\operatorname{ZS}(\mathbb U)_v=\ker\Sigma_v$.
\item\label{d:reflected-projections} For each arrow $\alpha:v\to u$ of $Q$, define
\begin{equation}\label{eq:zero-sum-arrow}
 \operatorname{ZS}(\mathbb U)_\alpha\colon
 \operatorname{ZS}(\mathbb U)_v\longrightarrow\operatorname{ZS}(\mathbb U)_u,
 \qquad z\longmapsto z_\alpha.
\end{equation}
\item Let $\widetilde Q=(Q,c)$ be a colored quiver.
Let $\widetilde{\mathbb Q}=(\widetilde Q,(I_Q,J_Q))$ be a colored bipartite quiver.
Let $\operatorname{ZS}(\widetilde{\mathbb Q})=\operatorname{ZS}(\mathbb U)$.
\end{enumerate}
\end{definition}

\begin{definition}[Colored source-trivalent trees and their quivers]\label{def:source-trivalent-tree}
Let $\mathbb T=(T,(I_T,J_T))$ be a bipartite graph such that $T$ is a finite tree.
\begin{enumerate}[font=\itshape]
\item We call $\mathbb T$ \textit{source-trivalent} if
\[
 \deg_T(u)=3\qquad(u\in I_T).
\]
\item A \textit{colored source-trivalent tree} is a pair $(\mathbb T,\kappa)$ such that $\mathbb T$ is source-trivalent and
\[
 \kappa\colon\operatorname{Edge}(T)\longrightarrow[0,2]_{\Z}
\]
is a proper edge coloring.
\item\label{d:reflected-quiver} For a colored source-trivalent tree $(\mathbb T,\kappa)$, define the colored bipartite quiver $Q(\mathbb T,\kappa)$ by
\[
\begin{aligned}
 I_{Q(\mathbb T,\kappa)}
 &=\{v\in J_{T}:\deg_T(v)\ge2\},\\
 J_{Q(\mathbb T,\kappa)}&=I_{T}.
\end{aligned}
\]
Its vertex set is $I_{Q(\mathbb T,\kappa)}\sqcup J_{Q(\mathbb T,\kappa)}$.
We define its arrows, tail and head maps, and colors by
\[
\begin{aligned}
 \operatorname{Edge}(Q(\mathbb T,\kappa))
 &=\{(v,u)\in I_{Q(\mathbb T,\kappa)}\times J_{Q(\mathbb T,\kappa)}:
 \llbracket u,v\rrbracket\in\operatorname{Edge}(T)\},\\
 s_{Q(\mathbb T,\kappa)}(v,u)&=v,\qquad t_{Q(\mathbb T,\kappa)}(v,u)=u,\\
 \operatorname{QCol}(Q(\mathbb T,\kappa))(v,u)&=\kappa(\llbracket u,v\rrbracket).
\end{aligned}
\]
\end{enumerate}
\end{definition}

\begin{lemma}[Source dimension of a subrepresentation]\label{lem:zero-sum-source-bound}
Let $(\mathbb T,\kappa)$ be a colored source-trivalent tree.
Let $\mathbb T=(T,(I_T,J_T))$.
Let $W\in\operatorname{SubColRep}(\operatorname{ZS}(Q(\mathbb T,\kappa)))$ satisfy $W\ne0$.
Let
\[
\begin{gathered}
 S=\{u\in I_{T}:\dim W_u=1\},\qquad n=\#S,\\
 u_W=\operatorname{SrcDim}_{\C,(I_{Q(\mathbb T,\kappa)},J_{Q(\mathbb T,\kappa)})}(W).
\end{gathered}
\]
Then, $n\ge1$ and $u_W\le n-1$.
\end{lemma}
\begin{proof}
Let us bound $\dim_{\C}W_v$ for $v\in I_{Q(\mathbb T,\kappa)}$.
Let
\[
 N=\{v\in J_{T}:\llbracket u,v\rrbracket\in\operatorname{Edge}(T)
       \text{ for some }u\in S\}.
\]
For $v\in J_{T}$, let
\[
 q_v=\#(S\cap\operatorname{SrcAdj}_{\mathbb T}(v)).
\]
Fix $v\in I_{Q(\mathbb T,\kappa)}$.
The arrows starting at $v$ are $(v,u)$ for $u\in\operatorname{SrcAdj}_{\mathbb T}(v)$.
For $z\in W_v$, equations~\eqref{eq:zero-sum-map} and~\eqref{eq:zero-sum-arrow} and the subrepresentation condition give
\begin{equation}\label{eq:zero-sum-subspace}
 \sum_{u\in\operatorname{SrcAdj}_{\mathbb T}(v)}z_{(v,u)}=0,
 \qquad z_{(v,u)}\in W_u
 \quad(u\in\operatorname{SrcAdj}_{\mathbb T}(v)).
\end{equation}
If $u\notin S$, then $W_u=0$ and $z_{(v,u)}=0$.
Thus, only the $q_v$ coordinates indexed by $u\in S$ can be nonzero. Their sum is zero by~\eqref{eq:zero-sum-subspace}.
If $q_v>0$, then $\dim_{\C}W_v\le q_v-1$. If $q_v=0$, then $W_v=0$.
If $n=0$, then $W_u=0$ for each $u\in J_{Q(\mathbb T,\kappa)}$ by the definition of $S$.
The bound on $\dim_{\C}W_v$ gives $W_v=0$ for each $v\in I_{Q(\mathbb T,\kappa)}$.
Then, $W=0$, contrary to the hypothesis. Hence, $n\ge1$.

Let us count the edges of $F=\operatorname{Subgraph}(T,S\sqcup N)$.
By Definition~\ref{def:graph-notions}(\ref{d:graph-induced-subgraph}),
\[
 \operatorname{Edge}(F)=
 \{\llbracket u,v\rrbracket\in\operatorname{Edge}(T):u\in S,\ v\in N\}.
\]
The graph $F$ is a forest because $T$ has no cycles.
Let $C=\#\operatorname{CtdComp}(F)\ge1$.
Each $u\in S$ has degree $3$ in $F$.
Each $v\in N\cap\operatorname{Leaf}(T)$ satisfies $q_v=1$.
For $v\in I_{Q(\mathbb T,\kappa)}\setminus N$, we have $W_v=0$.
Therefore,
\[
 3n=\#\operatorname{Edge}(F)=n+\#N-C,\qquad
 u_W\le\sum_{v\in N}(q_v-1)=3n-\#N=n-C\le n-1.
\]
\end{proof}

\begin{lemma}\label{lem:reflected-blocks}
Let $(\mathbb T,\kappa)$ be a colored source-trivalent tree.
Let $\mathbb T=(T,(I_T,J_T))$.
Assume $\#I_{T}=\nu\ge2$.
Let $H$ be the quiver underlying $Q(\mathbb T,\kappa)$ and let $d=\dim_{\C}(\operatorname{ZS}(Q(\mathbb T,\kappa)))$.
Then, we have the following.

\begin{enumerate}
\item\label{c:reflected-blocks-dimension} The representation $\operatorname{ZS}(Q(\mathbb T,\kappa))$ satisfies the following.
\begin{enumerate}[label=(\alph*),ref=\theenumi(\alph*)]
\item\label{c:reflected-blocks-dimension-vector} We have
\[
 \operatorname{TotDim}_{\C,(I_{Q(\mathbb T,\kappa)},J_{Q(\mathbb T,\kappa)})}(\operatorname{ZS}(Q(\mathbb T,\kappa)))=(\nu-1,\nu).
\]
\item\label{c:reflected-blocks-tits} We have $Q_H(d)=1$.
\end{enumerate}
\item\label{c:reflected-blocks-sinks} Each $u\in J_{Q(\mathbb T,\kappa)}$ satisfies
$\dim_{\C}(\operatorname{ZS}(Q(\mathbb T,\kappa))_u)=1$.
\item\label{c:reflected-blocks-sources} Each $v\in I_{Q(\mathbb T,\kappa)}$ satisfies
\[
 \bigl(\dim_{\C}(\operatorname{ZS}(Q(\mathbb T,\kappa))_v),\deg_{Q(\mathbb T,\kappa)}(v)\bigr)
 \in\{(1,2),(2,3)\}.
\]
\item\label{c:reflected-blocks-stable} The representation $\operatorname{ZS}(Q(\mathbb T,\kappa))$ is stable.
\item\label{c:reflected-blocks-generation} For each $u\in J_{Q(\mathbb T,\kappa)}$, we have
\[
 \sum_{\substack{\alpha\in\operatorname{Edge}(Q(\mathbb T,\kappa))\\t_{Q(\mathbb T,\kappa)}(\alpha)=u}}
 \operatorname{im}\bigl(\operatorname{ZS}(Q(\mathbb T,\kappa))_\alpha\bigr)
 =\operatorname{ZS}(Q(\mathbb T,\kappa))_u.
\]
\end{enumerate}
\end{lemma}
\begin{proof}
\prf{Claim~\ref{c:reflected-blocks-dimension-vector}} Let us compute the total dimension from the tree degrees.
Let $s=\#J_{T}$.
The identity $\#\operatorname{Edge}(T)=\#\operatorname{Vert}(T)-1$ gives
\[
 3\nu=\nu+s-1,\qquad s=2\nu+1.
\]
For $v\in J_{T}$, let $n_v=\deg_T(v)$.
By Definition~\ref{def:reflected-block}(\ref{d:reflected-kernel}), we obtain
\[
\begin{aligned}
 \operatorname{SrcDim}_{\C,(I_{Q(\mathbb T,\kappa)},J_{Q(\mathbb T,\kappa)})}(\operatorname{ZS}(Q(\mathbb T,\kappa)))
 &=\sum_{v\in J_{T}}(n_v-1)=3\nu-(2\nu+1)=\nu-1,\\
 \operatorname{SnkDim}_{\C,(I_{Q(\mathbb T,\kappa)},J_{Q(\mathbb T,\kappa)})}(\operatorname{ZS}(Q(\mathbb T,\kappa)))&=\#I_{T}=\nu.
\end{aligned}
\]
\prf{Claim~\ref{c:reflected-blocks-tits}} Claim~\ref{c:reflected-blocks-dimension-vector} gives
\[
Q_H(d)=\nu+\sum_{v\in J_{T}}\bigl((n_v-1)^2-n_v(n_v-1)\bigr)
 =\nu-\sum_{v\in J_{T}}(n_v-1)=1.
\]

\prf{Claim~\ref{c:reflected-blocks-sinks}} Let us count the sinks after deleting the leaves.
The sinks of $Q(\mathbb T,\kappa)$ belong to
$I_{T}$. By Definition~\ref{def:reflected-block}(\ref{d:reflected-kernel}),
$\operatorname{ZS}(Q(\mathbb T,\kappa))_u=\C$ for every such vertex $u$.

\prf{Claim~\ref{c:reflected-blocks-sources}} Let us count the sources after deleting the leaves.
By Definition~\ref{def:source-trivalent-tree}(2), the incident edges at each vertex have distinct colors in $[0,2]_{\Z}$. Hence, $n_v\le3$ for $v\in J_T$.
The source vertices of $Q(\mathbb T,\kappa)$ are
the vertices $v\in J_{T}$ satisfying $n_v\in[2,3]_{\Z}$.
For such a vertex, $\dim_{\C}(\operatorname{ZS}(Q(\mathbb T,\kappa))_v)=n_v-1$ and $\deg_{Q(\mathbb T,\kappa)}(v)=n_v$.

\prf{Claim~\ref{c:reflected-blocks-stable}} Let us prove stability by bounding each subrepresentation weight.
Let $W\in\operatorname{SubColRep}(\operatorname{ZS}(Q(\mathbb T,\kappa)))$ satisfy $0\ne W\ne\operatorname{ZS}(Q(\mathbb T,\kappa))$.
Claim~\ref{c:reflected-blocks-dimension-vector} gives the dimension pair $(\nu-1,\nu)$.
Define
\[
 S=\{u\in I_{T}:\dim W_u=1\},\qquad n=\#S,
\]
and let $u_W=\operatorname{SrcDim}_{\C,(I_{Q(\mathbb T,\kappa)},J_{Q(\mathbb T,\kappa)})}(W)$.
Lemma~\ref{lem:zero-sum-source-bound} gives $n\ge1$ and $u_W\le n-1$.
For $n<\nu$, we obtain
\[
 (\nu-1)n-\nu u_W\ge\nu-n>0.
\]
If $n=\nu$, properness gives $u_W\le\nu-2$. Hence,
\[
 (\nu-1)n-\nu u_W=\nu(\nu-1-u_W)\ge\nu>0.
\]

\prf{Claim~\ref{c:reflected-blocks-generation}} Let us verify surjectivity of the arrow maps in Definition~\ref{def:reflected-block}(\ref{d:reflected-projections}).
Since $T$ is connected and $\nu\ge2$, every $u\in I_{T}$
has a vertex $v\in\operatorname{SnkAdj}_{\mathbb T}(u)$ of degree $d\ge2$.
Let $e_1,\ldots,e_d$ be the standard basis of $\C^d$.
If $i$ and $j$ are distinct coordinate indices, then
\[
 e_i-e_j\in\ker\!\left(\C^d\xrightarrow{\;\sum\;}\C\right),
 \qquad (e_i-e_j)_i=1.
\]
Thus, each arrow map from this zero-sum space to $\C$ is surjective.
\end{proof}

\begin{lemma}\label{lem:vacant-slots}

Let $(\mathbb T,\kappa)$ be a colored source-trivalent tree.
Let $\mathbb T=(T,(I_T,J_T))$.
The map sending $(u,d)\in\operatorname{ColVacSlot}(Q(\mathbb T,\kappa))$ to the unique $v\in J_T$ satisfying
\[
 \deg_T(v)=1,\qquad \llbracket u,v\rrbracket\in\operatorname{Edge}(T),\qquad
 \kappa(\llbracket u,v\rrbracket)=d
\]
is a bijection onto $J_T\cap\operatorname{Leaf}(T)$.
\end{lemma}
\begin{proof}
Let us construct the inverse map from sink leaves to colored-vacant slots.
For $u\in I_{T}$, the three edges in $\operatorname{IncEdge}_T(u)$
have distinct colors in $[0,2]_{\Z}$.
The quiver $Q(\mathbb T,\kappa)$ keeps exactly the edges from $u$ to vertices of degree at least two.
Thus, $(u,d)\in\operatorname{ColVacSlot}(Q(\mathbb T,\kappa))$ determines the unique leaf
$v\in J_{T}$ such that $\llbracket u,v\rrbracket\in\operatorname{Edge}(T)$
and $\kappa(\llbracket u,v\rrbracket)=d$.
Conversely, let $v\in J_T$ be a leaf and let $u\in I_T$ satisfy $\operatorname{SrcAdj}_{\mathbb T}(v)=\{u\}$.
The inverse map is
\[
 J_T\cap\operatorname{Leaf}(T)\longrightarrow\operatorname{ColVacSlot}(Q(\mathbb T,\kappa)),
 \qquad v\longmapsto\bigl(u,\kappa(\llbracket u,v\rrbracket)\bigr).
\]
Uniqueness of the neighbor of $v$ and of the edge of color $d$ at $u$ shows that the two maps are inverse.
\end{proof}

\begin{definition}[Gluing zero-sum representations]\label{def:attachments}
Let $\mathbb S=(S,(I_S,J_S))$ be a source-trivalent tree.
For each $b\in J_S$, choose a colored source-trivalent tree $(\mathbb T_b,\kappa_b)$.
Let $\mathbb T_b=(T_b,(I_{T_b},J_{T_b}))$ for $b\in J_S$.
For each $e=\llbracket x,b\rrbracket\in\operatorname{Edge}(S)$ with
$x\in I_{S}$ and $b\in J_{S}$, choose
\[
 (v_e,d_e)\in\operatorname{ColVacSlot}(Q(\mathbb T_b,\kappa_b)),\qquad
 0\ne f_e:\C\longrightarrow\operatorname{ZS}(Q(\mathbb T_b,\kappa_b))_{v_e}.
\]
Let
\[
 \mathbf T=((\mathbb T_b,\kappa_b))_{b\in J_{S}},\qquad
 \mathbf p=((v_e,d_e))_{e\in\operatorname{Edge}(S)},\qquad
 \mathbf f=(f_e)_{e\in\operatorname{Edge}(S)}.
\]
We define the colored representation
\[
 R=\operatorname{Glue}_{\C}(\mathbb S;\mathbf T,\mathbf p,\mathbf f)
\]
by the following conditions.
\begin{enumerate}[font=\itshape]
\item Its quiver consists of disjoint copies of the $Q(\mathbb T_b,\kappa_b)$, a new vertex
$x$ for each $x\in I_{S}$, and one new arrow
\[
 \alpha_e:x\longrightarrow v_e,\qquad c(\alpha_e)=d_e
 \qquad(e=\llbracket x,b\rrbracket\in\operatorname{Edge}(S)).
\]
Each $v_e$ belongs to the copy of $Q(\mathbb T_b,\kappa_b)$ indexed by $b$.
The sources are the new vertices $x\in I_S$ and the copied vertices in $I_{Q(\mathbb T_b,\kappa_b)}$; the sinks are the copied vertices in $J_{Q(\mathbb T_b,\kappa_b)}$.
\item Its restriction to each $Q(\mathbb T_b,\kappa_b)$ is $\operatorname{ZS}(Q(\mathbb T_b,\kappa_b))$.
\item For each new vertex $x$, let $R_x=\C$.
For each new arrow $\alpha_e$, let $R_{\alpha_e}=f_e$.
\end{enumerate}
\end{definition}
\subsection{Flow sum equations}\label{subsec:flows-vertex-sums}
\begin{definition}[Flow sum equations]\label{def:flows-vertex-sums}
Let $\mathbb S=(S,(I_S,J_S))$ be a bipartite graph with a finite graph $S$. Let
\[
 \alpha=(\alpha_v\in\R)_{v\in I_{S}},\qquad
 \xi=(\xi_i\in\R)_{i\in J_{S}}.
\]
\begin{enumerate}[font=\itshape]
\item We define $\operatorname{Flow}(S)=\R^{\operatorname{Edge}(S)}$.
For $f\in\operatorname{Flow}(S)$, let $f_w=f(w)$.
We call
\[
\begin{aligned}
 \sum_{w\in\operatorname{IncEdge}_S(v)}f_w&=\alpha_v
 &&(v\in I_{S}),\\
 \sum_{w\in\operatorname{IncEdge}_S(i)}f_w&=\xi_i
 &&(i\in J_{S})
\end{aligned}
\]
the \textit{$(\alpha,\xi)$-flow sum equations on $S$}.
\item Define the set of \textit{$(\alpha,\xi)$-flows on $S$} by
\[
 \operatorname{Flow}(S,(\alpha,\xi))=
 \{f\in\operatorname{Flow}(S):f\text{ satisfies the }(\alpha,\xi)\text{-flow sum equations}\}.
\]
\item We call an $(\alpha,\xi)$-flow $f$ \textit{positive} if $f_w>0$ for every $w\in\operatorname{Edge}(S)$.
\end{enumerate}
\end{definition}

\begin{lemma}[Flows on a bipartite tree]\label{lem:tree-flow-cut}
Let $\mathbb S=(S,(I_S,J_S))$ be a bipartite graph such that $S$ is a finite tree.
Let $\alpha=(\alpha_v)_{v\in I_S}\in\R^{I_S}$ and
$\xi=(\xi_i)_{i\in J_S}\in\R^{J_S}$.
Assume
\begin{equation}\label{eq:tree-vertex-sum-balance}
 \sum_{v\in I_{S}}\alpha_v
 =\sum_{i\in J_{S}}\xi_i.
\end{equation}
Then, we have the following.

\begin{enumerate}
\item\label{c:tree-flow-existence} The set $\operatorname{Flow}(S,(\alpha,\xi))$ has exactly one element $f$.
\item\label{c:tree-flow-cut} Let $w=\llbracket v,i\rrbracket\in\operatorname{Edge}(S)$
with $i\in J_{S}$.
Let $S_w=\operatorname{SnkCtdComp}(\mathbb S;w)$.
Then,
\[
 f_w=\sum_{j\in \operatorname{SubSnk}_{\mathbb S}(S_w)}\xi_j
      -\sum_{u\in \operatorname{SubSrc}_{\mathbb S}(S_w)}\alpha_u.
\]
\item\label{c:tree-flow-positive-boundary} Assume $f_w>0$ for every $w\in\operatorname{Edge}(S)$.
Let $I\subseteq J_{S}$ and $J\subseteq I_{S}$.
Assume $\varnothing\ne I\sqcup J\subsetneq\operatorname{Vert}(S)$.
Assume $\operatorname{SnkAdj}_{\mathbb S}(v)\subseteq I$ for every $v\in J$.
Then,
\[
 \sum_{i\in I}\xi_i-\sum_{v\in J}\alpha_v>0.
\]
\end{enumerate}
\end{lemma}
\begin{proof}
\prf{Claim~\ref{c:tree-flow-existence}} Let us prove existence and uniqueness by induction on the number of vertices.

First, suppose $S$ has one vertex. If the vertex is a source $v$, equation~\eqref{eq:tree-vertex-sum-balance} gives $\alpha_v=0$.
If the vertex is a sink $i$, that equation gives $\xi_i=0$.
The empty map is the unique $(\alpha,\xi)$-flow.

Second, suppose $\#\operatorname{Vert}(S)\ge2$. Choose $z\in\operatorname{Leaf}(S)$ and let
\[
 S'=\operatorname{Subgraph}(S,\operatorname{Vert}(S)\setminus\{z\}).
\]
The graph $S'$ is a tree with one fewer vertex.

First, let us remove a source leaf $z=v\in I_S$. Let $\operatorname{SnkAdj}_{\mathbb S}(v)=\{i\}$ and $w=\llbracket v,i\rrbracket$.
The flow equation at $v$ forces $f_w=\alpha_v$.
On $S'$, define
\[
 \alpha'_u=\alpha_u\quad(u\in I_S\setminus\{v\}),\qquad
 \xi'_j=\begin{cases}\xi_i-\alpha_v,&j=i,\\\xi_j,&j\in J_S\setminus\{i\}.\end{cases}
\]
Both sums in~\eqref{eq:tree-vertex-sum-balance} decrease by $\alpha_v$.

Second, let us remove a sink leaf $z=i\in J_S$. Let $\operatorname{SrcAdj}_{\mathbb S}(i)=\{v\}$ and $w=\llbracket v,i\rrbracket$.
The flow equation at $i$ forces $f_w=\xi_i$.
On $S'$, define
\[
 \xi'_j=\xi_j\quad(j\in J_S\setminus\{i\}),\qquad
 \alpha'_u=\begin{cases}\alpha_v-\xi_i,&u=v,\\\alpha_u,&u\in I_S\setminus\{v\}.\end{cases}
\]
Both sums in~\eqref{eq:tree-vertex-sum-balance} decrease by $\xi_i$.

Let us extend the flow from $S'$ to $S$.
In either case, $(\alpha',\xi')$ satisfies~\eqref{eq:tree-vertex-sum-balance} on $S'$.
By induction, there is a unique $(\alpha',\xi')$-flow on $S'$.
The flow extends uniquely to $S$ by the forced value of $f_w$.

\prf{Claim~\ref{c:tree-flow-cut}} Let us compute the weight of the cut edge. Sum the sink equations over $S_w$ and subtract the source equations over $S_w$.
For every $e\in\operatorname{Edge}(S_w)$, the terms $f_e$ cancel. The edge $w$ contributes $f_w$.

\prf{Claim~\ref{c:tree-flow-positive-boundary}} Let us cancel the terms indexed by edges $\llbracket v,i\rrbracket$ with $v\in J$ and $i\in I$.
The assumption $\operatorname{SnkAdj}_{\mathbb S}(v)\subseteq I$ for every $v\in J$ gives
\[
 \sum_{i\in I}\xi_i-\sum_{v\in J}\alpha_v
 =\sum_{\substack{\llbracket v,i\rrbracket\in\operatorname{Edge}(S)\\
                  i\in I,\ v\notin J}}f_{\llbracket v,i\rrbracket}.
\]
Since $S$ is connected and $\varnothing\ne I\sqcup J\subsetneq\operatorname{Vert}(S)$,
there exist $i\in I$ and $v\notin J$ such that $\llbracket v,i\rrbracket\in\operatorname{Edge}(S)$.
Every summand is positive.
\end{proof}

\begin{lemma}[Weights of subrepresentations of a zero-sum representation]\label{lem:zero-sum-weight}
Let $(\mathbb T,\kappa)$ be a colored source-trivalent tree.
Let $\mathbb T=(T,(I_T,J_T))$, $Q=Q(\mathbb T,\kappa)$, and $Z=\operatorname{ZS}(Q)$.
Let $a_0\in\Z_{\ge0}$ satisfy
\[
 \operatorname{TotDim}_{\C,(I_Q,J_Q)}(Z)=(a_0,a_0+1).
\]
Let $0\ne W\in\operatorname{SubColRep}(Z)$. Let
\[
 (u,t)=\operatorname{TotDim}_{\C,(I_Q,J_Q)}(W).
\]
Let $y\in\R_{\ge0}$.
Then, the following statements hold.
\begin{enumerate}
\item\label{c:zero-sum-weight-bound} We have
\[
 y(t-u)-u\ge y-a_0.
\]
\item\label{c:zero-sum-weight-equality} We have
\[
 y(t-u)-u=y-a_0\quad\Longleftrightarrow\quad W=Z.
\]
\end{enumerate}
\end{lemma}
\begin{proof}
\prf{Claim~\ref{c:zero-sum-weight-bound}} Let us prove $t-u\ge1$ and $u\le a_0$.
Lemma~\ref{lem:zero-sum-source-bound} gives $t-u\ge1$.
For every $z\in I_Q$, the inclusion $W_z\subseteq Z_z$ gives
$\dim_\C W_z\le\dim_\C Z_z$.
Summing over $I_Q$ gives $u\le a_0$.
Therefore,
\[
 y(t-u)-u-(y-a_0)=y(t-u-1)+(a_0-u)\ge0.
\]

\prf{Claim~\ref{c:zero-sum-weight-equality}} Let us prove that a nonzero proper subrepresentation satisfies $u<a_0$.
Every $z\in J_Q=I_T$ has $Z_z=\C$.
Hence, $\#I_T=\#J_Q=a_0+1$.
If $a_0=0$, then $J_Q=\{z_0\}$ for one vertex $z_0$,
$Z_{z_0}=\C$. Every source space is zero.
The condition $W\ne0$ gives $W_{z_0}=\C$. Thus, $W=Z$.

Let us treat the case in which the subrepresentation has full source dimension.
Suppose that $a_0\ge1$ and $u=a_0$.
Then,
\[
 0=a_0-u=\sum_{z\in I_Q}
       \bigl(\dim_\C Z_z-\dim_\C W_z\bigr).
\]
Each summand is nonnegative.
Thus, $W_z=Z_z$ for every $z\in I_Q$.
Since $\#I_T=a_0+1\ge2$, Claim~\ref{c:reflected-blocks-generation}
of Lemma~\ref{lem:reflected-blocks} applies.
For every $z\in J_Q$, the subrepresentation condition gives
\[
 Z_z
 =\sum_{\substack{\beta\in\operatorname{Edge}(Q)\\t_Q(\beta)=z}}
       Z_\beta(Z_{s_Q(\beta)})
 =\sum_{\substack{\beta\in\operatorname{Edge}(Q)\\t_Q(\beta)=z}}
       Z_\beta(W_{s_Q(\beta)})
 \subseteq W_z\subseteq Z_z.
\]
Consequently, $W=Z$.
For $W\ne Z$, we therefore have $a_0-u>0$, and
\[
 y(t-u)-u-(y-a_0)=y(t-u-1)+(a_0-u)>0.
\]
For $W=Z$, substitution of $(u,t)=(a_0,a_0+1)$ gives equality.
\end{proof}

\begin{lemma}[Weights of subrepresentations of a gluing]\label{lem:gluing-subrepresentation-weights}
Let $\mathbb S=(S,(I_S,J_S))$ be a source-trivalent tree.
For $i\in J_S$, let $(\mathbb T_i,\kappa_i)$ be a colored source-trivalent tree, and let
\[
 Q_i=Q(\mathbb T_i,\kappa_i),\qquad Z_i=\operatorname{ZS}(Q_i).
\]
Let $a_i\in\Z_{\ge0}$ satisfy
\[
 \operatorname{TotDim}_{\C,(I_{Q_i},J_{Q_i})}(Z_i)=(a_i,a_i+1)
 \qquad(i\in J_S).
\]
For every $e=\llbracket v,i\rrbracket\in\operatorname{Edge}(S)$ with
$v\in I_S$ and $i\in J_S$, fix a colored-vacant slot and a nonzero complex linear map
\[
 p_e=(z_e,d_e)\in\operatorname{ColVacSlot}(Q_i),\qquad
 0\ne\phi_e:\C\longrightarrow(Z_i)_{z_e}.
\]
Let
\[
 \mathbf T=((\mathbb T_i,\kappa_i))_{i\in J_S},\qquad
 \mathbf p=(p_e)_{e\in\operatorname{Edge}(S)},\qquad
 \mathbf f=(\phi_e)_{e\in\operatorname{Edge}(S)},
\]
and let $R=\operatorname{Glue}_\C(\mathbb S;\mathbf T,\mathbf p,\mathbf f)$.
Use the vertex copies in Definition~\ref{def:attachments}. Let
\[
 I_R=I_S\sqcup\bigsqcup_{i\in J_S}I_{Q_i},\qquad
 J_R=\bigsqcup_{i\in J_S}J_{Q_i}.
\]
Let
\[
 c=\#I_S,\qquad k=c+1,\qquad
 a=\sum_{i\in J_S}a_i+c,\qquad x=\frac ak.
\]
Let $\alpha_v=x+1$ for $v\in I_S$ and $\xi_i=x-a_i$ for $i\in J_S$.
Fix $f\in\operatorname{Flow}(S,(\alpha,\xi))$.
Then, the following statements hold.
\begin{enumerate}
\item\label{c:gluing-weight-decomposition} Let $W\in\operatorname{SubColRep}(R)$.
For $i\in J_S$, let $W_i$ be the subrepresentation of $Z_i$ defined by
\[
 (W_i)_z=W_z\quad(z\in\operatorname{Vert}(Q_i)),\qquad
 (W_i)_\beta=R_\beta|_{W_{s_{Q_i}(\beta)}}
 \quad(\beta\in\operatorname{Edge}(Q_i)).
\]
Define
\[
 I=\{i\in J_S:W_i\ne0\},\qquad
 J=\{v\in I_S:W_v\ne0\}.
\]
\Needspace{5\baselineskip}
Let
\[
 (u_i,t_i)=\operatorname{TotDim}_{\C,(I_{Q_i},J_{Q_i})}(W_i)
 \quad(i\in I),\qquad
 (u,t)=\operatorname{TotDim}_{\C,(I_R,J_R)}(W).
\]
Then, $\operatorname{SnkAdj}_{\mathbb S}(v)\subseteq I$ for every $v\in J$. We have
\begin{equation}\label{eq:gluing-boundary-flow}
 \sum_{i\in I}(x-a_i)-\#J(x+1)
 =
 \sum_{\substack{e=\llbracket v,i\rrbracket\in\operatorname{Edge}(S)\\
                  v\in I_S\setminus J,\ i\in I}}f_e.
\end{equation}
Moreover,
\begin{equation}\label{eq:gluing-subobject-weight}
\begin{aligned}
 \frac{at-(a+k)u}{k}
 &=
 \sum_{\substack{e=\llbracket v,i\rrbracket\in\operatorname{Edge}(S)\\
                  v\in I_S\setminus J,\ i\in I}}f_e\\
 &\quad+\sum_{i\in I}
       \bigl(x(t_i-u_i-1)+(a_i-u_i)\bigr).
\end{aligned}
\end{equation}
\item\label{c:gluing-cut-subrepresentation} Let
$w=\llbracket v,i\rrbracket\in\operatorname{Edge}(S)$ with $v\in I_S$ and $i\in J_S$.
Let
\[
 S_w=\operatorname{SnkCtdComp}(\mathbb S;w),\qquad
 A_w=\operatorname{SubSnk}_{\mathbb S}(S_w),\qquad
 B_w=\operatorname{SubSrc}_{\mathbb S}(S_w).
\]
For each vertex $z$ of the quiver of $R$, define the subspace
\[
 W^w_z=
 \begin{cases}
 R_z,&z\in B_w\sqcup\bigsqcup_{j\in A_w}\operatorname{Vert}(Q_j),\\
 0,&\text{otherwise}.
 \end{cases}
\]
These subspaces and the restrictions of the arrow maps of $R$ define
a nonzero proper subrepresentation $W^w$.
For $(u_w,t_w)=\operatorname{TotDim}_{\C,(I_R,J_R)}(W^w)$, we have
\[
 u_w=\sum_{j\in A_w}a_j+\#B_w,\qquad
 t_w=\sum_{j\in A_w}(a_j+1),\qquad
 at_w-(a+k)u_w=kf_w.
\]
\end{enumerate}
\end{lemma}
\begin{proof}
\prf{Claim~\ref{c:gluing-weight-decomposition}} Let us determine which restrictions of $W$ are nonzero.
For $v\in J$, the inclusion $0\ne W_v\subseteq R_v=\C$ gives $W_v=\C$.
For $e=\llbracket v,i\rrbracket\in\operatorname{Edge}(S)$ with $v\in J$, we have
\[
 0\ne\phi_e(\C)=\phi_e(W_v)\subseteq W_{z_e}.
\]
Thus, $W_i\ne0$ and $i\in I$.
This proves $\operatorname{SnkAdj}_{\mathbb S}(v)\subseteq I$ for $v\in J$.

Let us sum the flow equations over $I$ and $J$.
They give
\[
 \sum_{i\in I}(x-a_i)-\#J(x+1)
 =
 \sum_{i\in I}\sum_{e\in\operatorname{IncEdge}_S(i)}f_e
 -
 \sum_{v\in J}\sum_{e\in\operatorname{IncEdge}_S(v)}f_e.
\]
An edge between $J$ and $I$ occurs in both sums.
There is no edge from $J$ to $J_S\setminus I$.
Subtracting the two finite sums proves~\eqref{eq:gluing-boundary-flow}.

Let us compute the stability weight from the vertex dimensions.
The disjoint vertex sets in the gluing give
\[
 u=\sum_{i\in I}u_i+\#J,\qquad t=\sum_{i\in I}t_i.
\]
Since $a=kx$ and $k>0$, we obtain
\[
\begin{aligned}
 \frac{at-(a+k)u}{k}
 &=x(t-u)-u\\
 &=\sum_{i\in I}\bigl(x(t_i-u_i)-u_i\bigr)-\#J(x+1)\\
 &=\sum_{i\in I}(x-a_i)-\#J(x+1)
   +\sum_{i\in I}\bigl(x(t_i-u_i-1)+(a_i-u_i)\bigr).
\end{aligned}
\]
Substitution of~\eqref{eq:gluing-boundary-flow} proves~\eqref{eq:gluing-subobject-weight}.

\prf{Claim~\ref{c:gluing-cut-subrepresentation}} Let us verify the subrepresentation condition at each arrow.
For $\beta\in\operatorname{Edge}(Q_j)$, both endpoint spaces of $W^w$ are full if $j\in A_w$ and zero if $j\notin A_w$.
For an added arrow $\alpha_e:y\to z_e$ indexed by $e=\llbracket y,j\rrbracket$, suppose $W^w_y\ne0$.
Then, $y\in B_w$.
The source endpoint $v$ of $w$ does not belong to $B_w$. Thus, $e\ne w$.
The edge $e$ therefore belongs to the component $S_w$ whenever $y\in B_w$.
Hence, $j\in A_w$ and
\[
 R_{\alpha_e}(W^w_y)=\phi_e(R_y)\subseteq R_{z_e}=W^w_{z_e}.
\]
If $W^w_y=0$, the subrepresentation condition holds because $\phi_e(0)=0$.

Let us check that the cut subrepresentation is nonzero and proper.
The index $i$ belongs to $A_w$. The representation $Z_i$ has sink dimension $a_i+1>0$.
Thus, $W^w\ne0$.
Also, $v\notin B_w$ gives $W^w_v=0\ne R_v=\C$. Thus, $W^w\ne R$.

Let us compute the dimension pair and its weight.
The full constituent spaces over $A_w$ and the full new source spaces over $B_w$ give
\[
 u_w=\sum_{j\in A_w}a_j+\#B_w,\qquad
 t_w=\sum_{j\in A_w}(a_j+1).
\]
Therefore,
\[
\begin{aligned}
 \frac{at_w-(a+k)u_w}{k}
 &=x(t_w-u_w)-u_w\\
 &=\sum_{j\in A_w}(x-a_j)-\sum_{y\in B_w}(x+1)=f_w.
\end{aligned}
\]
Summing the source and sink equations for $f$ gives~\eqref{eq:tree-vertex-sum-balance}.
The equality $\sum_{j\in A_w}(x-a_j)-\sum_{y\in B_w}(x+1)=f_w$ follows from Claim~\ref{c:tree-flow-cut} of Lemma~\ref{lem:tree-flow-cut}.
Multiplying by $k$ gives $at_w-(a+k)u_w=kf_w$.
\end{proof}

\begin{lemma}[Stability under gluing]\label{lem:flow-stability}
Let $\mathbb S=(S,(I_S,J_S))$ be a source-trivalent tree.
Let $c=\#I_{S}$.
For $i\in J_S$, let $(\mathbb T_i,\kappa_i)$ be a colored source-trivalent tree.
Let $\mathbf T=((\mathbb T_i,\kappa_i))_{i\in J_S}$.
Let $R=\operatorname{Glue}_{\C}(\mathbb S;\mathbf T,\mathbf p,\mathbf f)$ be a gluing of their zero-sum representations using colored-vacant slots $\mathbf p$ and nonzero complex linear maps $\mathbf f$.
For $i\in J_S$, let
\[
 \operatorname{TotDim}_{\C,(I_{Q(\mathbb T_i,\kappa_i)},J_{Q(\mathbb T_i,\kappa_i)})}
 (\operatorname{ZS}(Q(\mathbb T_i,\kappa_i)))=(a_i,a_i+1).
\]
Let
\[
 \widetilde Q_R=\operatorname{SuppColQuiv}(R),\qquad
 G_R=\operatorname{UnderG}(\widetilde Q_R).
\]
Let $\widetilde{\mathbb Q}_R=(\widetilde Q_R,(I_R,J_R))$ for
\[
\begin{aligned}
 I_R&=\operatorname{RepSupp}(R)\cap
 \left(I_S\sqcup\bigsqcup_{i\in J_S}I_{Q(\mathbb T_i,\kappa_i)}\right),\\
 J_R&=\operatorname{RepSupp}(R)\cap
 \bigsqcup_{i\in J_S}J_{Q(\mathbb T_i,\kappa_i)}.
\end{aligned}
\]
The disjoint unions use the vertex copies in $R$.
Assume the following conditions on the added arrows.
\begin{hypotheses}
\item\label{hyp:gluing-attachments} Each edge $\llbracket x,i\rrbracket$ of $S$ with $x\in I_{S}$ and $i\in J_{S}$ specifies one added $(v,d)$-colored arrow $x\to v$ for some $(v,d)\in\operatorname{ColVacSlot}(Q(\mathbb T_i,\kappa_i))$.
\item\label{hyp:gluing-slots} Each $(v,d)\in\operatorname{ColVacSlot}(Q(\mathbb T_i,\kappa_i))$ receives at most one added $(v,d)$-colored arrow. The three arrows starting at each $x\in I_{S}$ have distinct colors.
\end{hypotheses}
Let
\[
 k=c+1,\qquad a=\sum_{i\in J_{S}}a_i+c,\qquad x=\frac ak.
\]
For $\alpha_v=x+1$ and $\xi_i=x-a_i$, the $(\alpha,\xi)$-flow sum equations on $S$ are
\begin{equation}
\sum_{w\in\operatorname{IncEdge}_S(v)}f_w=x+1\quad(v\in I_{S}),\qquad
\sum_{w\in\operatorname{IncEdge}_S(i)}f_w=x-a_i\quad(i\in J_{S}).\label{flow}
\end{equation}
Then, we have the following.

\begin{enumerate}
\Needspace{6\baselineskip}
\item\label{c:flow-stability-tree} The following statements hold.
\begin{enumerate}[label=(\alph*),ref=\theenumi(\alph*)]
\item\label{c:flow-stability-support-tree} The graph $G_R$ is a tree.
\item\label{c:flow-stability-support-colors} The coloring $\operatorname{QCol}(\widetilde Q_R)$ induces a proper edge coloring of $G_R$.
\item\label{c:flow-stability-block-count} $\#J_{S}=2c+1$.
\item\label{c:flow-stability-total-dimension} We have $\operatorname{TotDim}_{\C,(I_R,J_R)}(R)=(a,a+k)$.
\end{enumerate}
\item\label{c:flow-stability-unique} The $(\alpha,\xi)$-flow sum equations~\eqref{flow} have a unique solution.
\item\label{c:flow-stability-positive} Let $f$ be the unique element of $\operatorname{Flow}(S,(\alpha,\xi))$ from Claim~\ref{c:flow-stability-unique}.
The representation $R$ is stable if and only if $f_w>0$ for every $w\in\operatorname{Edge}(S)$.
\end{enumerate}
\end{lemma}
\begin{proof}
\prf{Claim~\ref{c:flow-stability-support-tree}} Let us verify that $G_R$ is a tree.
Let $B=J_{S}$ and $b=\#B$.
Counting the edges of $S$ gives $3c=c+b-1$. Thus, $b=2c+1$.
For $i\in B$, let
\[
 G_i=\operatorname{UnderG}(\operatorname{SuppColQuiv}(\operatorname{ZS}(Q(\mathbb T_i,\kappa_i)))),\qquad
 v_i=\#\operatorname{Vert}(G_i),\qquad e_i=\#\operatorname{Edge}(G_i).
\]
Definition~\ref{def:reflected-block}(\ref{d:reflected-kernel}) gives $\#I_{T_i}=a_i+1\ge1$.
By Definitions~\ref{def:reflected-block} and~\ref{def:source-trivalent-tree}, $G_i$ is obtained from $T_i$ by deleting its sink leaves.
Every path between retained vertices avoids the deleted leaves. Thus, $G_i$ is a nonempty tree and $e_i=v_i-1$.
Replacing each $i\in B$ by $G_i$ and adding the arrows in assumption~\ref{hyp:gluing-attachments} gives the connected graph $G_R$.
Then,
\[
\begin{aligned}
 \#\operatorname{Vert}(G_R)&=c+\sum_{i\in B}v_i,\\
 \#\operatorname{Edge}(G_R)&=3c+\sum_{i\in B}e_i
 =3c+\sum_{i\in B}(v_i-1)
 =\#\operatorname{Vert}(G_R)-1.
\end{aligned}
\]
Hence, $G_R$ is a tree. Two colored-vacant slots sharing a vertex contribute two distinct added arrows to $3c$.
\prf{Claim~\ref{c:flow-stability-support-colors}} Assumption~\ref{hyp:gluing-slots} gives distinct colors at each new source.
At each old vertex, the old incident arrows have distinct colors because $\kappa_i$ is proper.
Each added arrow uses a colored-vacant slot by assumption~\ref{hyp:gluing-attachments}.
Its color is therefore unused by the old arrows at its target.
Assumption~\ref{hyp:gluing-slots} gives distinct colors among the added arrows at each old vertex.

\prf{Claim~\ref{c:flow-stability-block-count}} Let us count the sinks of the gluing tree.
Let $b=\#J_{S}$.
The identity $3c=c+b-1$ in the proof of Claim~\ref{c:flow-stability-support-tree} gives $b=2c+1$.

\prf{Claim~\ref{c:flow-stability-total-dimension}}
By Claim~\ref{c:flow-stability-block-count}, we obtain
\[
\begin{aligned}
 \operatorname{SrcDim}_{\C,(I_R,J_R)}(R)
 &=\sum_{i\in J_{S}}a_i+c=a,\\
 \operatorname{SnkDim}_{\C,(I_R,J_R)}(R)
 &=\sum_{i\in J_{S}}(a_i+1)=a-c+2c+1=a+k.
\end{aligned}
\]

\prf{Claim~\ref{c:flow-stability-unique}} Let us verify equality of the two total vertex weights.
Claim~\ref{c:flow-stability-block-count} gives
\[
\begin{aligned}
 \sum_{i\in J_{S}}\xi_i
 &=\sum_{i\in J_{S}}(x-a_i)\\
 &=(2c+1)x-(a-c)=c(x+1)
 =\sum_{v\in I_{S}}\alpha_v.
\end{aligned}
\]
Claim~\ref{c:tree-flow-existence} of Lemma~\ref{lem:tree-flow-cut} gives the unique $(\alpha,\xi)$-flow.

\prf{Claim~\ref{c:flow-stability-positive}} Let us obtain stability from positive flows.
Assume $f_e>0$ for every $e\in\operatorname{Edge}(S)$.
Let $0\ne W\ne R$ be a subrepresentation of $R$.
For $i\in J_S$, let
\[
 Q_i=Q(\mathbb T_i,\kappa_i),\qquad Z_i=\operatorname{ZS}(Q_i),
\]
and let $W_i$ be the restriction of $W$ to the copy of $Q_i$.
Define
\[
 I=\{i\in J_S:W_i\ne0\},\qquad
 J=\{v\in I_S:W_v\ne0\}.
\]
\Needspace{5\baselineskip}
Let
\[
 (u_i,t_i)=\operatorname{TotDim}_{\C,(I_{Q_i},J_{Q_i})}(W_i)\quad(i\in I),
 \qquad
 (u,t)=\operatorname{TotDim}_{\C,(I_R,J_R)}(W).
\]
We have $a_i\ge0$, $a\ge0$, $k>0$, and $x\ge0$.
Claim~\ref{c:zero-sum-weight-bound} of Lemma~\ref{lem:zero-sum-weight} gives
\[
 x(t_i-u_i-1)+(a_i-u_i)\ge0\qquad(i\in I).
\]
Claim~\ref{c:gluing-weight-decomposition}
of Lemma~\ref{lem:gluing-subrepresentation-weights} gives
$\operatorname{SnkAdj}_{\mathbb S}(v)\subseteq I$ for $v\in J$ and
the weight formula~\eqref{eq:gluing-subobject-weight}.
The dimension pair of $R$ is $(a,a+k)$ by Claim~\ref{c:flow-stability-total-dimension}.

First, suppose $I\sqcup J\subsetneq\operatorname{Vert}(S)$.
Since $W\ne0$, we have $I\sqcup J\ne\varnothing$.
Equation~\eqref{eq:gluing-boundary-flow} and
Claim~\ref{c:tree-flow-positive-boundary} of Lemma~\ref{lem:tree-flow-cut} give
\[
 \sum_{\substack{e=\llbracket v,i\rrbracket\in\operatorname{Edge}(S)\\
                  v\in I_S\setminus J,\ i\in I}}f_e
 =\sum_{i\in I}(x-a_i)-\#J(x+1)>0.
\]
Hence,~\eqref{eq:gluing-subobject-weight} gives $at-(a+k)u>0$.

Second, suppose $I\sqcup J=\operatorname{Vert}(S)$.
Then, $I=J_S$, $J=I_S$, and $W_v=R_v=\C$ for every $v\in I_S$.
Since $W\ne R$, there is $i_0\in J_S$ such that $0\ne W_{i_0}\ne Z_{i_0}$.
Claim~\ref{c:zero-sum-weight-equality} of Lemma~\ref{lem:zero-sum-weight}
and Claim~\ref{c:zero-sum-weight-bound} of Lemma~\ref{lem:zero-sum-weight} give
\[
 x(t_{i_0}-u_{i_0}-1)+(a_{i_0}-u_{i_0})>0.
\]
For $i\in I\setminus\{i_0\}$, Claim~\ref{c:zero-sum-weight-bound} of Lemma~\ref{lem:zero-sum-weight} gives $x(t_i-u_i-1)+(a_i-u_i)\ge0$.
The boundary-edge sum in~\eqref{eq:gluing-subobject-weight} is empty because $J=I_S$.
Thus, $at-(a+k)u>0$ in this case as well.
By Definition~\ref{def:bipartite-stability}, $R$ is stable.

Let us prove the converse. Suppose $R$ is stable.
Fix $w=\llbracket v,i\rrbracket\in\operatorname{Edge}(S)$ with $v\in I_S$ and $i\in J_S$.
Claim~\ref{c:gluing-cut-subrepresentation}
of Lemma~\ref{lem:gluing-subrepresentation-weights}
gives a nonzero proper subrepresentation $W^w$ of $R$.
For
\[
 (u_w,t_w)=\operatorname{TotDim}_{\C,(I_R,J_R)}(W^w),
\]
Claim~\ref{c:gluing-cut-subrepresentation} of Lemma~\ref{lem:gluing-subrepresentation-weights} gives
\[
 at_w-(a+k)u_w=kf_w.
\]
Stability gives $kf_w>0$.
Since $k>0$, we obtain $f_w>0$.
\end{proof}

\subsection{Exact counts for fixed colored-vacant-slot colors}
\begin{definition}\label{def:marked-block-counts}
Let $j\in\Z_{\ge1}$. Let $d\in[0,2]_{\Z}$ and $\epsilon\in[0,2]_{\Z}\setminus\{d\}$.
We consider all colored source-trivalent trees $(\mathbb T,\kappa)$ such that
\[
 \operatorname{TotDim}_{\C,(I_{Q(\mathbb T,\kappa)},J_{Q(\mathbb T,\kappa)})}(\operatorname{ZS}(Q(\mathbb T,\kappa)))=(j,j+1).
\]
\begin{enumerate}[font=\itshape]
\item We define $P_j$ to be the number of classes
\[
 \operatorname{MarkedCIsom}_{\C}(Q(\mathbb T,\kappa),\operatorname{ZS}(Q(\mathbb T,\kappa));((v,d))),
 \qquad (v,d)\in\operatorname{ColVacSlot}(Q(\mathbb T,\kappa)).
\]
\item We define $Q_j^{=}$ to be the number of classes
\[
 \operatorname{MarkedCIsom}_{\C}(Q(\mathbb T,\kappa),\operatorname{ZS}(Q(\mathbb T,\kappa));((v_1,d),(v_2,d)))
\]
for $(v_i,d)\in\operatorname{ColVacSlot}(Q(\mathbb T,\kappa))$ with $i\in[2]_{\Z}$ and $v_1\ne v_2$.
\item We define $Q_j^{\ne}$ to be the number of classes
\[
 \operatorname{MarkedCIsom}_{\C}(Q(\mathbb T,\kappa),\operatorname{ZS}(Q(\mathbb T,\kappa));((v_1,d),(v_2,\epsilon)))
\]
for $(v_1,d)\in\operatorname{ColVacSlot}(Q(\mathbb T,\kappa))$ and $(v_2,\epsilon)\in\operatorname{ColVacSlot}(Q(\mathbb T,\kappa))$.
\end{enumerate}
\end{definition}
\begin{definition}[Colored branches and their series]\label{def:colored-branches}
Let $d\in[0,2]_{\Z}$. Let $t$, $y_0$, $y_1$, and $y_2$ be commuting indeterminates.
\begin{enumerate}[font=\itshape]
\item\label{d:branch-set} Let $\mathbb B=(B,(I_B,J_B))$ be a bipartite graph such that $B$ is a finite tree.
We define $\operatorname{Branch}(\mathbb B;d)$ to be the set of pairs $(\kappa,r)$ satisfying the following conditions.
\begin{enumerate}[label=(\alph*),ref=\theenumi(\alph*),font=\itshape]
\item $\kappa:\operatorname{Edge}(B)\to[0,2]_{\Z}$ is a proper edge coloring.
\item $r\in\operatorname{Vert}(B)$ and $\kappa(e)\ne d$ for every $e\in\operatorname{IncEdge}_B(r)$.
\item For every $v\in I_{B}$,
\[
 \deg_B(v)=\begin{cases}2,&v=r,\\3,&v\ne r.\end{cases}
\]
\end{enumerate}
A \textit{branch with parent color $d$} is a tuple $\mathfrak b=(\mathbb B,\kappa,r)$ such that $(\kappa,r)\in\operatorname{Branch}(\mathbb B;d)$.
We call $r$ its \textit{root}.
\item\label{d:branch-types} We call $\mathfrak b$ a \textit{source branch} if $r\in I_{B}$ and a \textit{sink branch} if $r\in J_{B}$.
\item\label{d:branch-parent} For $\mathfrak b=(\mathbb B,\kappa,r)$, form $B^+$ by adjoining a new vertex $p$ in the part opposite to $r$ and the edge $\llbracket p,r\rrbracket$.
We call $p$ the \textit{parent vertex} and $\llbracket p,r\rrbracket$ the \textit{parent edge}.
Extend $\kappa$ to $\kappa^+:\operatorname{Edge}(B^+)\to[0,2]_{\Z}$ by setting
\[
 \kappa^+(e)=\kappa(e)\quad(e\in\operatorname{Edge}(B)),\qquad
 \kappa^+(\llbracket p,r\rrbracket)=d.
\]
\item\label{d:branch-isomorphism} Let $\mathfrak b=(\mathbb B,\kappa,r)$ and $\mathfrak b'=(\mathbb B',\kappa',r')$ be branches with parent color $d$.
Let $\mathbb B=(B,(I_B,J_B))$ and $\mathbb B'=(B',(I_{B'},J_{B'}))$.
An isomorphism $\mathfrak b\to\mathfrak b'$ is a graph isomorphism
$\psi=(\psi_0,\psi_1):B\to B'$ such that $\psi_0(I_B)=I_{B'}$, $\psi_0(J_B)=J_{B'}$, and
\[
 \psi_0(r)=r',\qquad \kappa'(\psi_1(e))=\kappa(e)\quad(e\in\operatorname{Edge}(B)).
\]
Let $\mathcal B_d^{\mathrm{src}}$ and $\mathcal B_d^{\mathrm{snk}}$ be the sets of isomorphism classes of source and sink branches with parent color $d$.
\item\label{d:branch-statistics} For $\mathfrak b=(\mathbb B,\kappa,r)$, define
\[
 n(\mathfrak b)=\#I_{B}.
\]
For $e\in[0,2]_{\Z}$, let $\ell_e(\mathfrak b)$ be the number of vertices
$v\in J_{B}$ such that $\deg_{B^+}(v)=1$ and the unique edge incident to $v$ in $B^+$ has color $e$ under $\kappa^+$.
\item\label{d:branch-weight} We define the weight of $\mathfrak b$ by
\[
 w(\mathfrak b)=t^{n(\mathfrak b)}\prod_{e=0}^2 y_e^{\ell_e(\mathfrak b)}.
\]
\item\label{d:tree-weight} Let $\mathcal T$ be the set of color-preserving, bipartition-preserving isomorphism classes of colored source-trivalent trees $(\mathbb T,\kappa)$.
For $\mathbb T=(T,(I_T,J_T))$, define $n(\mathbb T)=\#I_T$.
For $e\in[0,2]_{\Z}$, let $\ell_e(\mathbb T,\kappa)$ count the sink leaves of $T$ whose unique incident edge has color $e$.
Define
\[
 w(\mathbb T,\kappa)=t^{n(\mathbb T)}\prod_{e=0}^2 y_e^{\ell_e(\mathbb T,\kappa)}.
\]
\item\label{d:branch-series} We define
\[
 T_d=\sum_{[\mathfrak b]\in\mathcal B_d^{\mathrm{src}}}w(\mathfrak b),\qquad
 D_d=\sum_{[\mathfrak b]\in\mathcal B_d^{\mathrm{snk}}}w(\mathfrak b),\qquad
 U=\sum_{[(\mathbb T,\kappa)]\in\mathcal T}w(\mathbb T,\kappa).
\]
\end{enumerate}
\end{definition}

\Needspace{7\baselineskip}
\begin{lemma}[Branch-series identities]\label{lem:branch-series}
Let $\{d,\epsilon,f\}=[0,2]_{\Z}$.
Then, we have the following.
\begin{enumerate}
\item\label{c:slots-finite} For every $d\in[0,2]_{\Z}$ and $n\in\Z_{\ge0}$, there are finitely many classes in each of
$\mathcal B_d^{\mathrm{src}}$, $\mathcal B_d^{\mathrm{snk}}$, and $\mathcal T$ with source count $n$.
The weights $w(\mathfrak b)$ and $w(\mathbb T,\kappa)$ are invariant under branch isomorphisms and color-preserving, bipartition-preserving tree isomorphisms, respectively.
Thus, $T_d$, $D_d$, and $U$ belong to $\Z[y_0,y_1,y_2][[t]]$.
\item\label{c:branch-series-recursions} We have
\begin{equation}\label{eq:marked-branch-series}
D_d=y_d+T_\epsilon+T_f+T_\epsilon T_f,\qquad T_d=tD_\epsilon D_f.
\end{equation}
\item\label{c:branch-series-leaf} We have
\[
 \frac{\partial U}{\partial y_d}=T_d.
\]
\end{enumerate}
\end{lemma}
\begin{proof}
\prf{Claim~\ref{c:slots-finite}} Let us count the vertices for each source count $n$.
For $\mathfrak b=(\mathbb B,\kappa,r)$ with $\mathbb B=(B,(I_B,J_B))$ and $\#I_B=n$, we have
\[
 \#\operatorname{Edge}(B)=\sum_{v\in I_B}\deg_B(v)
 =\begin{cases}
 3n-1,&r\in I_B,\\
 3n,&r\in J_B.
 \end{cases}
\]
Since $B$ is a tree,
\[
 \#\operatorname{Vert}(B)=\#\operatorname{Edge}(B)+1
 =\begin{cases}
 3n,&r\in I_B,\\
 3n+1,&r\in J_B.
 \end{cases}
\]
For a source-trivalent tree $\mathbb T=(T,(I_T,J_T))$ with $\#I_T=n$, we obtain
\[
 \#\operatorname{Edge}(T)=\sum_{v\in I_T}3=3n,\qquad
 \#\operatorname{Vert}(T)=3n+1.
\]
For each of these vertex counts $N$, relabel the vertices by $[N]_{\Z}$.
There are finitely many choices of
\[
 I\subseteq[N]_{\Z},\qquad J=[N]_{\Z}\setminus I,\qquad
 E\subseteq\{\llbracket u,v\rrbracket:u,v\in[N]_{\Z},\ u\ne v\},\qquad
 \kappa:E\to[0,2]_{\Z}.
\]
For a branch, there are also finitely many choices of $r\in[N]_{\Z}$.
Every class at source count $n$ has such a representative.
Thus, each set of classes at source count $n$ is finite.

Let us verify that isomorphisms preserve the weights.
Let $\psi=(\psi_0,\psi_1):\mathfrak b\to\mathfrak b'$ be an isomorphism of branches with parent color $d$.
Let $\mathfrak b'=(\mathbb B',\kappa',r')$ and $\mathbb B'=(B',(I_{B'},J_{B'}))$.
Let $p$ and $p'$ be their parent vertices.
Extend $\psi$ to $\psi^+:B^+\to B'^+$ by
\[
\begin{aligned}
 \psi_0^+(v)&=\psi_0(v)\quad(v\in\operatorname{Vert}(B)),&
 \psi_0^+(p)&=p',\\
 \psi_1^+(g)&=\psi_1(g)\quad(g\in\operatorname{Edge}(B)),&
 \psi_1^+(\llbracket p,r\rrbracket)&=\llbracket p',r'\rrbracket.
\end{aligned}
\]
Since $\psi_0(r)=r'$ and both parent edges have color $d$, we have
\[
\begin{aligned}
 \deg_{B'^+}(\psi_0^+(v))&=\deg_{B^+}(v)
 &&(v\in\operatorname{Vert}(B^+)),\\
 \kappa'^+(\psi_1^+(g))&=\kappa^+(g)
 &&(g\in\operatorname{Edge}(B^+)).
\end{aligned}
\]
Also, $\psi_0(I_B)=I_{B'}$ and $\psi_0(J_B)=J_{B'}$.
Consequently,
\[
 n(\mathfrak b)=n(\mathfrak b'),\qquad
 \ell_e(\mathfrak b)=\ell_e(\mathfrak b')\quad(e\in[0,2]_{\Z}),
\]
and
\[
 w(\mathfrak b)
 =t^{n(\mathfrak b)}\prod_{e=0}^2y_e^{\ell_e(\mathfrak b)}
 =t^{n(\mathfrak b')}\prod_{e=0}^2y_e^{\ell_e(\mathfrak b')}
 =w(\mathfrak b').
\]

Let us verify that isomorphisms of colored bipartite trees preserve their weights.
For a color-preserving, bipartition-preserving tree isomorphism
$(\mathbb T,\kappa)\to(\mathbb T',\kappa')$, preservation of degrees and colors gives
\[
 n(\mathbb T)=n(\mathbb T'),\qquad
 \ell_e(\mathbb T,\kappa)=\ell_e(\mathbb T',\kappa')
 \quad(e\in[0,2]_{\Z}).
\]
Hence, $w(\mathbb T,\kappa)=w(\mathbb T',\kappa')$.
For each $n$, the coefficient of $t^n$ in $T_d$, $D_d$, and $U$ is therefore a finite sum of well-defined monomials in $\Z[y_0,y_1,y_2]$.

\prf{Claim~\ref{c:branch-series-recursions}} Let us specify the branches obtained by deleting the root.
Let $\mathfrak b=(\mathbb B,\kappa,r)$ have parent color $d$ and let $\mathbb B=(B,(I_B,J_B))$.
Let
\[
 A=\{\kappa(g):g\in\operatorname{IncEdge}_B(r)\}
 \subseteq\{\epsilon,f\}.
\]
For each $e\in A$, proper coloring gives a unique vertex $u_e$ satisfying
\[
 \llbracket r,u_e\rrbracket\in\operatorname{Edge}(B),\qquad
 \kappa(\llbracket r,u_e\rrbracket)=e.
\]
Let $B_e$ be the connected component of
$\operatorname{Subgraph}(B,\operatorname{Vert}(B)\setminus\{r\})$ containing $u_e$.
Define
\[
\begin{aligned}
 I_e&=I_B\cap\operatorname{Vert}(B_e),&
 J_e&=J_B\cap\operatorname{Vert}(B_e),\\
 \kappa_e&=\kappa|_{\operatorname{Edge}(B_e)},&
 \mathfrak b_e&=((B_e,(I_e,J_e)),\kappa_e,u_e).
\end{aligned}
\]
Since $B$ is a tree, its components after deleting $r$ are exactly the graphs $B_e$ for $e\in A$.
Each $\mathfrak b_e$ has parent color $e$.

First, let us treat a sink root. Let $r\in J_B$.
Then, $u_e\in I_e$ and
\[
 \deg_{B_e}(v)=
 \begin{cases}
 2,&v=u_e,\\
 3,&v\in I_e\setminus\{u_e\}.
 \end{cases}
\]
Thus, $\mathfrak b_e$ is a source branch.
The map
\[
 [\mathfrak b]\longmapsto(A,([\mathfrak b_e])_{e\in A})
\]
gives a bijection
\[
 \mathcal B_d^{\mathrm{snk}}
 \longrightarrow
 \bigsqcup_{A\subseteq\{\epsilon,f\}}
 \prod_{e\in A}\mathcal B_e^{\mathrm{src}}.
\]
To construct its inverse, choose disjoint representatives
$\mathfrak b_e=((B_e,(I_e,J_e)),\kappa_e,u_e)$ for $e\in A$ and a new sink $r$.
Let
\[
\begin{aligned}
 I_B&=\bigsqcup_{e\in A}I_e,&
 J_B&=\{r\}\sqcup\bigsqcup_{e\in A}J_e,\\
 \operatorname{Edge}(B)&=
 \bigsqcup_{e\in A}\operatorname{Edge}(B_e)
 \sqcup\{\llbracket r,u_e\rrbracket:e\in A\}.&&
\end{aligned}
\]
Let $\kappa|_{\operatorname{Edge}(B_e)}=\kappa_e$ and
$\kappa(\llbracket r,u_e\rrbracket)=e$ for $e\in A$.
These constructions are inverse on isomorphism classes.
For $A=\varnothing$, the branch consists of the sink $r$ and has weight $y_d$.
For $A\ne\varnothing$, the root $r$ is not a sink leaf in $B^+$. Thus,
\[
 n(\mathfrak b)=\sum_{e\in A}n(\mathfrak b_e),\qquad
 \ell_h(\mathfrak b)=\sum_{e\in A}\ell_h(\mathfrak b_e)
 \quad(h\in[0,2]_{\Z}).
\]
Hence,
\[
\begin{aligned}
 D_d
 &=y_d+\sum_{\varnothing\ne A\subseteq\{\epsilon,f\}}
   \prod_{e\in A}T_e\\
 &=y_d+T_\epsilon+T_f+T_\epsilon T_f.
\end{aligned}
\]

Second, let us treat a source root. Let $r\in I_B$.
Then, $A=\{\epsilon,f\}$, $u_e\in J_e$, and every vertex of $I_e$ has degree three in $B_e$.
Thus, $\mathfrak b_\epsilon$ and $\mathfrak b_f$ are sink branches.
The map
\[
 [\mathfrak b]\longmapsto([\mathfrak b_\epsilon],[\mathfrak b_f])
\]
gives a bijection
\[
 \mathcal B_d^{\mathrm{src}}
 \longrightarrow\mathcal B_\epsilon^{\mathrm{snk}}\times\mathcal B_f^{\mathrm{snk}}.
\]
Its inverse adjoins a new source $r$ to disjoint representatives of the two sink branches and adds the edges
$\llbracket r,u_\epsilon\rrbracket$ and $\llbracket r,u_f\rrbracket$ of colors $\epsilon$ and $f$.
The vertex parts are
\[
 I_B=\{r\}\sqcup I_\epsilon\sqcup I_f,\qquad
 J_B=J_\epsilon\sqcup J_f.
\]
We have
\[
\begin{aligned}
 n(\mathfrak b)&=1+n(\mathfrak b_\epsilon)+n(\mathfrak b_f),\\
 \ell_h(\mathfrak b)&=\ell_h(\mathfrak b_\epsilon)+\ell_h(\mathfrak b_f)
 \quad(h\in[0,2]_{\Z}).
\end{aligned}
\]
Therefore,
\[
 T_d=\sum_{[\mathfrak b_\epsilon]\in\mathcal B_\epsilon^{\mathrm{snk}}}
     \sum_{[\mathfrak b_f]\in\mathcal B_f^{\mathrm{snk}}}
     t\,w(\mathfrak b_\epsilon)w(\mathfrak b_f)
 =tD_\epsilon D_f.
\]

\prf{Claim~\ref{c:branch-series-leaf}} Let us specify the deletion of a marked sink leaf.
Let $\mathbb T=(T,(I_T,J_T))$ and let $v\in J_T$ be a sink leaf.
Let $u\in I_T$ be its unique neighbor and assume
\[
 \kappa(\llbracket u,v\rrbracket)=d.
\]
Let
\[
\begin{aligned}
 B&=\operatorname{Subgraph}(T,\operatorname{Vert}(T)\setminus\{v\}),\\
 \mathbb B&=(B,(I_T,J_T\setminus\{v\})),\qquad
 \kappa_B=\kappa|_{\operatorname{Edge}(B)}.
\end{aligned}
\]
Then, $\mathfrak b=(\mathbb B,\kappa_B,u)$ is a source branch with parent color $d$, since
\[
 \deg_B(u)=2,\qquad
 \deg_B(x)=3\quad(x\in I_T\setminus\{u\}),\qquad
 \kappa_B(g)\ne d\quad(g\in\operatorname{IncEdge}_B(u)).
\]
Thus, deletion defines the map on marked isomorphism classes
\[
 [((\mathbb T,\kappa);v)]\longmapsto[(\mathbb B,\kappa_B,u)].
\]

Let us construct the inverse map.
Let $(\mathbb B,\kappa_B,u)$ be a source branch with $\mathbb B=(B,(I_B,J_B))$.
Adjoin a new sink $v$ and let
\[
\begin{aligned}
 I_T&=I_B,&J_T&=J_B\sqcup\{v\},\\
 \operatorname{Edge}(T)&=\operatorname{Edge}(B)\sqcup\{\llbracket u,v\rrbracket\},&&\\
 \kappa|_{\operatorname{Edge}(B)}&=\kappa_B,&
 \kappa(\llbracket u,v\rrbracket)&=d.
\end{aligned}
\]
The marked leaf is $v$.
Both operations preserve isomorphisms and their compositions return the original marked tree or branch up to isomorphism.

Let us compute the weight after deleting the marked leaf.
Let $n=n(\mathbb T)$ and $\ell_e=\ell_e(\mathbb T,\kappa)$ for $e\in[0,2]_{\Z}$.
Deleting $v$ gives
\[
 n(\mathfrak b)=n,\qquad
 \ell_e(\mathfrak b)=
 \begin{cases}
 \ell_d-1,&e=d,\\
 \ell_e,&e\ne d.
 \end{cases}
\]
By Lemma~\ref{lem:colored-tree-rigidity}, the $\ell_d$ choices of $v$ give $\ell_d$ distinct marked isomorphism classes.
For $\ell_d\ge1$, their total branch weight is
\[
 \ell_d t^n y_d^{\ell_d-1}
 \prod_{e\in[0,2]_{\Z}\setminus\{d\}}y_e^{\ell_e}
 =\frac{\partial}{\partial y_d}
   \left(t^n\prod_{e=0}^2y_e^{\ell_e}\right).
\]
For $\ell_d=0$, both contributions are zero.
Claim~\ref{c:slots-finite} makes each coefficient a finite sum. Thus,
\[
 \frac{\partial U}{\partial y_d}
 =\sum_{[\mathfrak b]\in\mathcal B_d^{\mathrm{src}}}w(\mathfrak b)
 =T_d.\qedhere
\]
\end{proof}

\Needspace{7\baselineskip}
\begin{lemma}\label{slots}
Let $j\in\Z_{>0}$.
Then, the following hold.

\begin{subequations}\label{counts}
\begin{enumerate}
\item\label{c:slots-single} We have
\begin{equation}
P_j=\frac1{3j+4}\binom{4j+4}{j+1}.
\end{equation}
\item\label{c:slots-equal} We have
\begin{equation}
Q_j^{=}=2\binom{4j+2}{j-1}.
\end{equation}
\item\label{c:slots-distinct} We have
\begin{equation}
Q_j^{\ne}=\frac{2j+3}{2j}Q_j^{=}\ge Q_j^{=}.
\end{equation}
\end{enumerate}
\end{subequations}
\end{lemma}
\begin{proof}
\prf{Claim~\ref{c:slots-single}} Let us recover a colored tree from its quiver.
Let $(\mathbb T,\kappa)$ be a colored source-trivalent tree.
Let $\mathbb T=(T,(I_T,J_T))$ and assume $\#I_T=j+1$.
Put $Q=Q(\mathbb T,\kappa)$.
For each $(u,e)\in\operatorname{ColVacSlot}(Q)$, choose a new vertex $v_{u,e}$.
Choose these vertices distinct and outside $\operatorname{Vert}(Q)$.
Up to a color-preserving isomorphism, the tree is recovered by
\[
 \begin{aligned}
 I_T&=J_Q,\\
 J_T&=I_Q\sqcup\{v_{u,e}:(u,e)\in\operatorname{ColVacSlot}(Q)\},\\
 \operatorname{Edge}(T)
 &=\{\llbracket t_Q(\alpha),s_Q(\alpha)\rrbracket:
       \alpha\in\operatorname{Edge}(Q)\}\\
 &\quad\sqcup\{\llbracket u,v_{u,e}\rrbracket:
       (u,e)\in\operatorname{ColVacSlot}(Q)\},\\
 \kappa(\llbracket t_Q(\alpha),s_Q(\alpha)\rrbracket)
 &=\operatorname{QCol}(Q)(\alpha),\qquad
 \kappa(\llbracket u,v_{u,e}\rrbracket)=e.
 \end{aligned}
\]
Lemma~\ref{lem:vacant-slots} identifies the marked slot $(u,e)$ with the marked leaf $v_{u,e}$.

Let us construct the induced isomorphism of zero-sum representations.
Let $(\mathbb T',\kappa')$ be another colored source-trivalent tree.
Let $\mathbb T'=(T',(I_{T'},J_{T'}))$ and assume $\#I_{T'}=j+1$.
Put $Q'=Q(\mathbb T',\kappa')$ and let
$\phi=(\phi_0,\phi_1):Q\to Q'$ be a color-preserving isomorphism preserving the bipartitions.
For $v\in I_Q$, define
\[
 F_v:\operatorname{ZS}(Q)_v\longrightarrow
        \operatorname{ZS}(Q')_{\phi_0(v)},\qquad
 (F_vz)_{\phi_1(\alpha)}=z_\alpha
 \quad(\alpha\in\operatorname{IncEdgeQuiv}_Q(v)).
\]
The map $\phi_1$ is a bijection on the incident arrow sets. Hence,
\[
 \sum_{\alpha'\in\operatorname{IncEdgeQuiv}_{Q'}(\phi_0(v))}
 (F_vz)_{\alpha'}
 =\sum_{\alpha\in\operatorname{IncEdgeQuiv}_Q(v)}z_\alpha=0.
\]
For $u\in J_Q$, set $F_u=1_{\C}$.
For every arrow $\alpha:v\to u$, the identity $(F_vz)_{\phi_1(\alpha)}=z_\alpha$ gives
\[
 \operatorname{ZS}(Q')_{\phi_1(\alpha)}\circ F_v
 =F_u\circ\operatorname{ZS}(Q)_\alpha.
\]
The inverse maps are obtained from $\phi^{-1}$. Thus,
$(\phi_0,\phi_1,(F_v)_{v\in\operatorname{Vert}(Q)})$ is an isomorphism of the zero-sum representations.

Let us identify the marked isomorphism classes.
The extension of $\phi_0$ to the deleted leaves is
\[
 v_{u,e}\longmapsto v'_{\phi_0(u),e}.
\]
The map $v_{u,e}\mapsto v'_{\phi_0(u),e}$ preserves the marked leaves exactly when $\phi_0$ preserves the marked slots.
Conversely, an isomorphism of marked representations includes an isomorphism of the marked colored quivers by Definition~\ref{def:colored-representation-markings}(\ref{d:marked-rep-iso}).
The reconstruction therefore gives a bijection between the marked tree classes and the marked representation classes.

Let us express $P_j$ as a coefficient.
Let $d\in[0,2]_{\Z}$ and let $T=T_0|_{y_0=y_1=y_2=1}$.
Permuting the colors gives $T_d|_{y_0=y_1=y_2=1}=T$ for every $d\in[0,2]_{\Z}$.
Claims~\ref{c:branch-series-recursions} and~\ref{c:branch-series-leaf} of Lemma~\ref{lem:branch-series} give
\[
 T=t(1+T)^4,\qquad T(0)=0,\qquad P_j=[t^{j+1}]T.
\]

Let us extract the coefficient of $t^\nu$ for an integer $\nu\ge1$. Let $u$ be a formal indeterminate.
Lagrange inversion gives
\[
 [t^\nu]T
 =\frac1\nu[u^{\nu-1}](1+u)^{4\nu}
 =\frac1\nu\binom{4\nu}{\nu-1}
 =\frac1{3\nu+1}\binom{4\nu}{\nu}.
\]
Taking $\nu=j+1$ proves Claim~\ref{c:slots-single}.

\prf{Claim~\ref{c:slots-equal}} Let us compute $Q_j^{=}$.
Fix $d\in[0,2]_{\Z}$ and $\epsilon\in[0,2]_{\Z}\setminus\{d\}$.
By Claim~\ref{c:branch-series-leaf} of Lemma~\ref{lem:branch-series}, differentiating $T_d$ with respect to $y_d$ counts an ordered pair of distinct sink leaves of color $d$.
Lemma~\ref{lem:vacant-slots} identifies these pairs with the marked colored-vacant slots counted by $Q_j^{=}$.

Let us compute the derivatives in~\eqref{eq:marked-branch-series}.
Write
\[
 A=\left.\frac{\partial T_d}{\partial y_d}\right|_{y_0=y_1=y_2=1},\qquad
 B=\left.\frac{\partial T_\epsilon}{\partial y_d}\right|_{y_0=y_1=y_2=1}
 \quad(\epsilon\ne d).
\]
The two possible colors $\epsilon\ne d$ give the same $B$ by color symmetry.
Differentiating both recursions gives
\[
 A=\frac{2T}{1+T}(A+B),\qquad
 B=\frac{T}{(1+T)^2}+\frac{T}{1+T}(A+3B).
\]
The equation $A=\frac{2T}{1+T}(A+B)$ gives $(1-T)A=2TB$.
Substituting $(1-T)A=2TB$ into $B=\frac{T}{(1+T)^2}+\frac{T}{1+T}(A+3B)$ gives
\[
 B=\frac{T(1-T)}{(1+T)(1-3T)},\qquad
 A=\frac{2T^2}{(1+T)(1-3T)}.
\]
All denominators have constant term one and are invertible as formal power series.

Let us extract the coefficient that counts the two markings.
Differentiating $T=t(1+T)^4$ with respect to $t$ gives
\[
 tT'=\frac{T(1+T)}{1-3T},\qquad A=tT'\frac{2T}{(1+T)^2}.
\]
The derivative form of Lagrange inversion gives, for $\nu\ge2$,
\[
 [t^\nu]A=[u^{\nu-1}]\frac{2u}{(1+u)^2}(1+u)^{4\nu}
 =2\binom{4\nu-2}{\nu-2}.
\]
Taking $\nu=j+1$ gives $Q_j^{=}=2\binom{4j+2}{j-1}$.

\prf{Claim~\ref{c:slots-distinct}} Let us compute $Q_j^{\ne}$.
Fix $d\in[0,2]_{\Z}$ and $\epsilon\in[0,2]_{\Z}\setminus\{d\}$.
By Claim~\ref{c:branch-series-leaf} of Lemma~\ref{lem:branch-series}, differentiating $T_d$ with respect to $y_\epsilon$ counts an ordered pair of sink leaves of colors $d$ and $\epsilon$.
Lemma~\ref{lem:vacant-slots} identifies these pairs with the marked colored-vacant slots counted by $Q_j^{\ne}$.
The derivative calculation in Claim~\ref{c:slots-equal} and color symmetry give
\[
 \left.\frac{\partial T_d}{\partial y_\epsilon}\right|_{y_0=y_1=y_2=1}
 =\frac{T(1-T)}{(1+T)(1-3T)}
 =tT'\frac{1-T}{(1+T)^2}.
\]

Let us extract the coefficient for an integer $\nu\ge2$.
Lagrange inversion gives
\[
 [t^\nu]\left(tT'\frac{1-T}{(1+T)^2}\right)
 =[u^{\nu-1}](1-u)(1+u)^{4\nu-2}
 =\binom{4\nu-2}{\nu-1}-\binom{4\nu-2}{\nu-2}.
\]
Taking $\nu=j+1$ and using Claim~\ref{c:slots-equal}, we obtain
\[
 \frac{Q_j^{\ne}}{Q_j^{=}}
 =\frac12\left(\frac{\binom{4j+2}{j}}{\binom{4j+2}{j-1}}-1\right)
 =\frac12\left(\frac{3j+3}{j}-1\right)=\frac{2j+3}{2j}.
\]
\end{proof}

\Needspace{6\baselineskip}
\begin{definition}[Admissible pairs]\label{def:admissible-pairs}
Let $(p_j)_{j\in\Z_{>0}}$ and $(q_j)_{j\in\Z_{>0}}$ be real sequences and let $N\in\Z_{>0}$.
We call a pair $(A,B)\in\R_{>0}^2$ \textit{$(p_j,q_j,N)$-admissible} if
\[
\begin{aligned}
 p_j&\ge A\beta^j j^{-\frac32},\\
 q_j&\ge B\beta^j j^{-\frac12}
\end{aligned}
\qquad(j\in\Z_{\ge N}).
\]
\end{definition}

\begin{lemma}\label{lem:slot-bounds}
There exists a $(P_j,Q_j^{=},256)$-admissible pair $(A,B)$.
\end{lemma}
\begin{proof}
Let us obtain positive lower bounds for the normalized counts.
Claims~\ref{c:slots-single} and~\ref{c:slots-equal} of Lemma~\ref{slots} and Stirling's formula give
\[
 P_j\beta^{-j}j^{\frac{3}{2}}\longrightarrow
 \frac\beta3\sqrt{\frac{2}{3\pi}},\qquad
 Q_j^{=}\beta^{-j}j^{\frac{1}{2}}\longrightarrow
 \frac{32}{27}\sqrt{\frac{2}{3\pi}}.
\]

Let us include the finitely many smaller indices in these bounds.
Put
\[
 L_P=\frac\beta3\sqrt{\frac{2}{3\pi}},\qquad
 L_Q=\frac{32}{27}\sqrt{\frac{2}{3\pi}}.
\]
Choose $J>256$ such that, for every integer $j\ge J$,
\[
 P_j\beta^{-j}j^{\frac32}\ge\frac{L_P}{2},\qquad
 Q_j^{=}\beta^{-j}j^{\frac12}\ge\frac{L_Q}{2}.
\]
The explicit formulas in Lemma~\ref{slots} show that both counts are positive for $j\ge256$.
Hence, the finite minima
\[
\begin{aligned}
 A&=\min\left(\left\{\frac{L_P}{2}\right\}
       \cup\{P_j\beta^{-j}j^{\frac32}:256\le j<J\}\right)>0,\\
 B&=\min\left(\left\{\frac{L_Q}{2}\right\}
       \cup\{Q_j^{=}\beta^{-j}j^{\frac12}:256\le j<J\}\right)>0
\end{aligned}
\]
give a $(P_j,Q_j^{=},256)$-admissible pair.
\end{proof}

\subsection{A family of positive flows}
\begin{definition}[Labeled source-trivalent chains]\label{def:labeled-chain}
Let $c\in\Z_{>0}$. We define the \textit{labeled source-trivalent chain} $\operatorname{LSTrivCh}(c)$ as follows.
The vertices $x_t$ for $t\in[c]_{\Z}$, $z_t$ for $t\in[0,c]_{\Z}$, and $y_t$ for $t\in[c]_{\Z}$ are pairwise distinct.
\begin{enumerate}[font=\itshape]
\item Its bipartition is
\[
\begin{aligned}
 I_{\operatorname{LSTrivCh}(c)}&=\{x_t:t\in[c]_{\Z}\},\\
 J_{\operatorname{LSTrivCh}(c)}&=
 \{z_t:t\in[0,c]_{\Z}\}\sqcup\{y_t:t\in[c]_{\Z}\}.
\end{aligned}
\]
\item For $t\in[c]_{\Z}$, define
\[
\begin{aligned}
 \operatorname{LeftEdge}_c(t)&=\llbracket x_t,z_{t-1}\rrbracket,\\
 \operatorname{MidEdge}_c(t)&=\llbracket x_t,y_t\rrbracket,\\
 \operatorname{RightEdge}_c(t)&=\llbracket x_t,z_t\rrbracket.
\end{aligned}
\]
The edge set of $\operatorname{LSTrivCh}(c)$ consists of $\operatorname{LeftEdge}_c(t)$, $\operatorname{MidEdge}_c(t)$, and $\operatorname{RightEdge}_c(t)$ for $t\in[c]_{\Z}$.
We call them the left, middle, and right edges at $x_t$, respectively.
\end{enumerate}
\end{definition}

\begin{lemma}\label{lem:labeled-chain-tree}
For $c\in\Z_{>0}$, let $S=\operatorname{LSTrivCh}(c)$ and $\mathbb S=(S,(I_S,J_S))$.
The bipartite graph $\mathbb S$ is a source-trivalent tree.
Its sink sets satisfy
\[
 \begin{aligned}
 J_S\cap\operatorname{Leaf}(S)&=\{z_0,z_c\}\sqcup\{y_t:t\in[c]_{\Z}\},\\
 \{v\in J_S:\deg_S(v)=2\}&=\{z_t:t\in[c-1]_{\Z}\}.
 \end{aligned}
\]
\end{lemma}
\begin{proof}
First, suppose $c=1$. The graph is a tree with center $x_1$ and three leaves $z_0,y_1,z_1$.
Its only source has degree three, and it has no degree-two sinks.

Second, suppose $c\ge2$. Let us delete $z_0,z_c,y_1,\ldots,y_c$.
The remaining vertices and edges form the path
\[
 x_1,z_1,x_2,\ldots,z_{c-1},x_c.
\]
Each deleted vertex is joined to exactly one source by one edge.
Adding these vertices to the path preserves connectedness and creates no cycle.
Thus, $S$ is a tree. Its degrees are
\[
\begin{aligned}
 \deg_S(x_t)&=3 &&(1\le t\le c),\\
 \deg_S(z_t)&=2 &&(1\le t<c),\\
 \deg_S(z_0)=\deg_S(z_c)=\deg_S(y_t)&=1 &&(1\le t\le c).
\end{aligned}
\]
These formulas give the stated source degrees, sink leaves, and degree-two sinks.
\end{proof}

\begin{definition}[Chain flow family and source dimensions]\label{def:chain-flow-family}
Let $a\in\Z_{>0}$ and $k\in\Z_{>0}$ satisfy $k\ge2$.
Let $c=k-1$, $x=\frac ak$, and $S=\operatorname{LSTrivCh}(c)$.
For $f\in\operatorname{Flow}(S)$ and $t\in[c]_{\Z}$, define
\[
 l_t(f)=f_{\operatorname{LeftEdge}_c(t)},\qquad
 m_t(f)=f_{\operatorname{MidEdge}_c(t)},\qquad
 r_t(f)=f_{\operatorname{RightEdge}_c(t)}.
\]
Define $\mathcal F_{a,k}$ to be the set of $f\in\operatorname{Flow}(S)$ such that, for every $t\in[c]_{\Z}$,
\[
\begin{aligned}
 l_t(f)&\in[\frac{x}{256},\frac{2x}{256})\cap(tx+\Z),\\
 r_t(f)&\in(0,\frac{x}{256})\cap(-tx+\Z),\\
 m_t(f)&=x+1-l_t(f)-r_t(f).
\end{aligned}
\]
For $f\in\operatorname{Flow}(S)$ and $b\in J_S$, define the \textit{source dimension}
\[
 a_b(f)=x-\sum_{w\in\operatorname{IncEdge}_S(b)}f_w.
\]
\end{definition}

\begin{lemma}[Flow choices]\label{lem:flow-choices}
Let $a\in\Z_{>0}$ and $k\in\Z_{>0}$ satisfy $k\ge2$ and $a\ge100000k$.
Let $c=k-1$, $x=\frac ak$, and $S=\operatorname{LSTrivCh}(c)$.
Let
\[
 N=256,\qquad \gamma=1-\frac Nx,\qquad \lambda=\frac1{100}.
\]
Let $f\in\mathcal F_{a,k}$.
Let $l_t=l_t(f)$, $m_t=m_t(f)$, and $r_t=r_t(f)$ for $t\in[c]_{\Z}$.
Then, we have the following.

\begin{enumerate}
\item\label{c:flow-choices-dimensions} The flow $f$ and the source dimensions $(a_b(f))_{b\in J_{S}}$ satisfy the following.
\begin{enumerate}[label=(\alph*),ref=\theenumi(\alph*)]
\item\label{c:flow-choices-positive} We have $l_t>0$, $m_t>0$, and $r_t>0$ for $t\in[c]_{\Z}$.
\item\label{c:flow-choices-integral} For $b\in J_S$, we have $a_b(f)\in\Z$.
\item\label{c:flow-choices-sum} We have $\sum_{b\in J_S}a_b(f)=a-c$.
\item\label{c:flow-choices-range} For every $b\in J_S$, we have
\[
 256<a_b(f)<x.
\]
\end{enumerate}
\item\label{c:flow-choices-count} For $t\in[c]_{\Z}$, we have
\[
\begin{aligned}
 \#\left([\frac{x}{N},\frac{2x}{N})\cap(tx+\mathbb Z)\right)&\ge\frac{\gamma x}{N},\\
 \#\left((0,\frac{x}{N})\cap(-tx+\mathbb Z)\right)&\ge\frac{\gamma x}{N}.
\end{aligned}
\]
\item\label{c:flow-choices-distinct} For $f,f'\in\mathcal F_{a,k}$, we have
\[
 \bigl(a_b(f)=a_b(f')\text{ for every }b\in J_S\bigr)
 \quad\Longrightarrow\quad f=f'.
\]
\item\label{c:flow-choices-stable} Let $\mathbf T=((\mathbb T_b,\kappa_b))_{b\in J_S}$ be a family of colored source-trivalent trees satisfying
\[
 \operatorname{TotDim}_{\C,(I_{Q(\mathbb T_b,\kappa_b)},J_{Q(\mathbb T_b,\kappa_b)})}
 (\operatorname{ZS}(Q(\mathbb T_b,\kappa_b)))=(a_b(f),a_b(f)+1)
 \qquad(b\in J_S).
\]
Let $\mathbf p=((v_e,d_e))_{e\in\operatorname{Edge}(S)}$ be a family of colored-vacant slots and $\mathbf g=(g_e)_{e\in\operatorname{Edge}(S)}$ a family of nonzero complex linear maps such that
\[
 (v_e,d_e)\in\operatorname{ColVacSlot}(Q(\mathbb T_b,\kappa_b)),\qquad
 g_e:\C\longrightarrow\operatorname{ZS}(Q(\mathbb T_b,\kappa_b))_{v_e}
\]
for $e=\llbracket x,b\rrbracket\in\operatorname{Edge}(S)$, $x\in I_S$, and $b\in J_S$.
Let
\[
 R=\operatorname{Glue}_{\C}((S,(I_S,J_S));\mathbf T,\mathbf p,\mathbf g)
\]
satisfy hypotheses~\ref{hyp:gluing-attachments} and~\ref{hyp:gluing-slots} of Lemma~\ref{lem:flow-stability}.
Then, $R$ is stable.
\item\label{c:flow-choices-weight} For $t\in[c-1]_{\Z}$, we have
\begin{equation}\label{weight}
 (l_t+r_t-1)^{-\frac32}(x-r_t-l_{t+1})^{-\frac12}
 \ge\lambda N^2x^{-2}.
\end{equation}
The dimensions $a_{z_0}(f)$, $a_{z_c}(f)$, and $a_{y_c}(f)$ satisfy
\[
 (x-l_1)^{-\frac32}(x-r_c)^{-\frac32}(l_c+r_c-1)^{-\frac32}
 \ge x^{-\frac92}.
\]
\end{enumerate}
\end{lemma}
\begin{proof}
\prf{Claim~\ref{c:flow-choices-positive}} Let us prove positivity of all edge weights.
Since $f\in\mathcal F_{a,k}$, $l_t>0$, $r_t>0$, and $l_t+r_t<\frac{3x}{N}$. Hence,
\[
 m_t=x+1-l_t-r_t>x\left(1-\frac3N\right)+1>0.
\]

\prf{Claim~\ref{c:flow-choices-integral}} Let us prove $a_b(f)\in\Z$ for $b\in J_S$.
For $S=\operatorname{LSTrivCh}(c)$, the dimensions $a_b(f)$ satisfy
\[
\begin{aligned}
 a_{z_0}(f)&=x-l_1,& a_{z_c}(f)&=x-r_c,\\
 a_{y_t}(f)&=l_t+r_t-1&&(t\in[c]_{\Z}),\\
 a_{z_t}(f)&=x-r_t-l_{t+1}&&(t\in[c-1]_{\Z}).
\end{aligned}
\]
 The conditions $l_t-tx\in\Z$, $r_t+tx\in\Z$, and
$kx=a\in\Z$ give
\[
\begin{aligned}
 x-l_1&=-(l_1-x)\in\Z,\\
 x-r_c&=a-(r_c+cx)\in\Z,\\
 l_t+r_t-1&=(l_t-tx)+(r_t+tx)-1\in\Z,\\
 x-r_t-l_{t+1}&=-(r_t+tx)-(l_{t+1}-(t+1)x)\in\Z.
\end{aligned}
\]

\prf{Claim~\ref{c:flow-choices-sum}}
The dimensions $a_b(f)$ in Claim~\ref{c:flow-choices-integral} satisfy
\[
\begin{aligned}
 &(x-l_1)+(x-r_c)+\sum_{t=1}^c(l_t+r_t-1)
       +\sum_{t=1}^{c-1}(x-r_t-l_{t+1})\\
 &\qquad=(c+1)x-c=a-c.
\end{aligned}
\]

\prf{Claim~\ref{c:flow-choices-range}} Let us bound each source dimension between $256$ and $x$.
The bound $a\ge100000k$ gives $x\ge100000$.
First, for $t\in[c]_{\Z}$, the bounds $l_t\in[\frac{x}{N},\frac{2x}{N})$ and $r_t\in(0,\frac{x}{N})$ give
\[
 256<\frac{x}{N}-1<a_{y_t}(f)=l_t+r_t-1<\frac{3x}{N}-1<x.
\]
Second, for $t\in[0,c]_{\Z}$, the formulas in Claim~\ref{c:flow-choices-integral} give
\[
 256<x\left(1-\frac3N\right)<a_{z_t}(f)<x.
\]
For $t=0$, use $a_{z_0}(f)=x-l_1$ and $0<l_1<\frac{2x}{N}$.
For $t=c$, use $a_{z_c}(f)=x-r_c$ and $0<r_c<\frac{x}{N}$.
For $t\in[c-1]_{\Z}$, use $a_{z_t}(f)=x-r_t-l_{t+1}$ and $0<r_t+l_{t+1}<\frac{3x}{N}$.

\prf{Claim~\ref{c:flow-choices-count}} Let us count the integers after translating each interval.
For $t\in[c]_{\Z}$, we have
\[
 \begin{aligned}
 \#\left([\frac{x}{N},\frac{2x}{N})\cap(tx+\Z)\right)
 &=\left\lceil\frac{2x}{N}-tx\right\rceil
   -\left\lceil\frac{x}{N}-tx\right\rceil
 \ge\frac{x}{N}-1,\\
 \#\left((0,\frac{x}{N})\cap(-tx+\Z)\right)
 &=\left\lceil\frac{x}{N}+tx\right\rceil-\lfloor tx\rfloor-1
 \ge\frac{x}{N}-1.
 \end{aligned}
\]
Since $\frac{x}{N}-1=\frac{\gamma x}{N}$, both inequalities in Claim~\ref{c:flow-choices-count} follow.

\prf{Claim~\ref{c:flow-choices-distinct}} Let us recover the flow from its source dimensions.
Let $f\in\mathcal F_{a,k}$ and $\widetilde f\in\mathcal F_{a,k}$ satisfy $a_b(f)=a_b(\widetilde f)$ for every $b\in J_S$.
For $h\in\{f,\widetilde f\}$, we have
\[
 \sum_{w\in\operatorname{IncEdge}_S(v)}h_w=x+1\quad(v\in I_S),\qquad
 \sum_{w\in\operatorname{IncEdge}_S(b)}h_w=x-a_b(f)\quad(b\in J_S).
\]
Claim~\ref{c:tree-flow-existence} of Lemma~\ref{lem:tree-flow-cut} gives $f=\widetilde f$.

\prf{Claim~\ref{c:flow-choices-stable}} Let us prove stability.
Claim~\ref{c:flow-choices-sum} gives $\sum_{b\in J_S}a_b(f)+c=a$.
Definition~\ref{def:chain-flow-family} gives equations~\eqref{flow} for $f$ and the source dimensions $a_b(f)$.
Claim~\ref{c:flow-stability-unique} of Lemma~\ref{lem:flow-stability} therefore identifies $f$ as the unique flow used in Claim~\ref{c:flow-stability-positive} of Lemma~\ref{lem:flow-stability}.
Claim~\ref{c:flow-choices-positive} gives $f_w>0$ for every $w\in\operatorname{Edge}(S)$.
Claim~\ref{c:flow-stability-positive} of Lemma~\ref{lem:flow-stability}
gives stability of $R$.

\prf{Claim~\ref{c:flow-choices-weight}} Let us bound the products of inverse powers of the source dimensions.
For $t\in[c-1]_{\Z}$, we have $a_{y_t}(f)<\frac{3x}{N}$ and $a_{z_t}(f)<\frac{(N-1)x}{N}$.
Hence,
\[
(l_t+r_t-1)^{-\frac{3}{2}}
(x-r_t-l_{t+1})^{-\frac{1}{2}}
\ge\frac{N^2x^{-2}}{\sqrt{3^3(N-1)}}
\ge\lambda N^2x^{-2}.
\]
The inequality $\frac{N^2x^{-2}}{\sqrt{3^3(N-1)}}\ge\lambda N^2x^{-2}$ uses $\lambda=\frac1{100}$ and $3^3(N-1)=6885<100^2$.
We also have $a_{z_0}(f)^{-\frac32}a_{z_c}(f)^{-\frac32}a_{y_c}(f)^{-\frac32}\ge x^{-\frac92}$.
\end{proof}

\subsection{Covering quivers and dimension vectors}\label{subsec:cover-definitions}
\Needspace{10\baselineskip}
\begin{definition}[Coverings, embeddings, and deck transformations]\label{def:quiver-covering}
Let $Q$ and $\widehat Q$ be quivers. 
\begin{enumerate}[font=\itshape]
\item An \textit{embedding} $\iota=(\iota_0,\iota_1):Q\to\widehat Q$ consists of injective maps
\[
 \iota_0:\operatorname{Vert}(Q)\hookrightarrow\operatorname{Vert}(\widehat Q),\qquad
 \iota_1:\operatorname{Edge}(Q)\hookrightarrow\operatorname{Edge}(\widehat Q)
\]
such that, for every $\alpha\in\operatorname{Edge}(Q)$,
\[
 s_{\widehat Q}(\iota_1(\alpha))=\iota_0(s_Q(\alpha)),\qquad
 t_{\widehat Q}(\iota_1(\alpha))=\iota_0(t_Q(\alpha)).
\]
\item A quiver homomorphism $p=(p_0,p_1):\widehat Q\to Q$ is a \textit{covering}
if $p_0$ is surjective and, for every $v\in\operatorname{Vert}(\widehat Q)$, the restrictions
\[
\begin{aligned}
 p_1:s_{\widehat Q}^{-1}(\{v\})&\longrightarrow s_Q^{-1}(\{p_0(v)\}),\\
 p_1:t_{\widehat Q}^{-1}(\{v\})&\longrightarrow t_Q^{-1}(\{p_0(v)\})
\end{aligned}
\]
are bijective.
We call $p$ a \textit{universal covering} if $\operatorname{UnderG}(\widehat Q)$ is a tree.
\item A \textit{deck transformation} of $p$ is a quiver isomorphism
$\delta:\widehat Q\to\widehat Q$ such that
\[
 p_0\circ\delta_0=p_0,\qquad p_1\circ\delta_1=p_1.
\]
We denote the group of these transformations under composition by $\operatorname{Deck}(p)$.
\end{enumerate}
\end{definition}

\begin{definition}[Arrow traversals and lifts]\label{def:path-lift}
Let $Q$ be a quiver and $n\in\Z_{\ge0}$. A \textit{traversal of length $n$} is a tuple
$\gamma=(v_0,\alpha_1,\varepsilon_1,v_1,\ldots,\alpha_n,\varepsilon_n,v_n)$ such that
$v_i\in\operatorname{Vert}(Q)$ for $i\in[0,n]_{\Z}$ and
\[
 \alpha_i\in\operatorname{Edge}(Q),\qquad \varepsilon_i\in\{1,-1\},\qquad
 (s_Q(\alpha_i),t_Q(\alpha_i))=
 \begin{cases}(v_{i-1},v_i)&\varepsilon_i=1,\\(v_i,v_{i-1})&\varepsilon_i=-1\end{cases}
 \quad(i\in[n]_{\Z}).
\]
For a covering $p:\widehat Q\to Q$, a \textit{lift} of $\gamma$ is a traversal
$\widehat\gamma=(\widehat v_0,\widehat\alpha_1,\varepsilon_1,\widehat v_1,\ldots,\widehat\alpha_n,\varepsilon_n,\widehat v_n)$
satisfying $p_0(\widehat v_i)=v_i$ for $i\in[0,n]_{\Z}$ and $p_1(\widehat\alpha_i)=\alpha_i$ for $i\in[n]_{\Z}$.
\end{definition}

\begin{lemma}[Unique traversal lifting]\label{lem:path-lift}
Let $p:\widehat Q\to Q$ be a covering.
Let $n\in\Z_{\ge0}$ and let
$\gamma=(v_0,\alpha_1,\varepsilon_1,v_1,\ldots,\alpha_n,\varepsilon_n,v_n)$
be a traversal in $Q$.
Let $\widehat v_0\in\operatorname{Vert}(\widehat Q)$ satisfy $p_0(\widehat v_0)=v_0$.
Then, there is exactly one traversal
$\widehat\gamma=(\widehat v_0,\widehat\alpha_1,\varepsilon_1,\widehat v_1,\ldots,\widehat\alpha_n,\varepsilon_n,\widehat v_n)$ in $\widehat Q$ satisfying
\[
 p_0(\widehat v_i)=v_i\quad(i\in[0,n]_{\Z}),\qquad
 p_1(\widehat\alpha_i)=\alpha_i\quad(i\in[n]_{\Z}).
\]
\end{lemma}
\begin{proof}
Let us construct the lift one arrow at a time.
For a traversal of length zero, the lift consists of the specified vertex $\widehat v_0$.
For $i\ge1$, suppose that the lifted traversal through $\widehat v_{i-1}$ has been determined and
$p_0(\widehat v_{i-1})=v_{i-1}$.

First, suppose $\varepsilon_i=1$. The covering gives a bijection
\[
 p_1:s_{\widehat Q}^{-1}(\{\widehat v_{i-1}\})
       \longrightarrow s_Q^{-1}(\{v_{i-1}\}).
\]
Let $\widehat\alpha_i$ be the unique preimage of $\alpha_i$ and set
\[
 \widehat v_i=t_{\widehat Q}(\widehat\alpha_i).
\]
Then,
\[
 p_0(\widehat v_i)=t_Q(p_1(\widehat\alpha_i))=t_Q(\alpha_i)=v_i.
\]

Second, suppose $\varepsilon_i=-1$. Use the bijection
\[
 p_1:t_{\widehat Q}^{-1}(\{\widehat v_{i-1}\})
       \longrightarrow t_Q^{-1}(\{v_{i-1}\}).
\]
Let $\widehat\alpha_i$ be the unique preimage of $\alpha_i$ and set
\[
 \widehat v_i=s_{\widehat Q}(\widehat\alpha_i).
\]
Then,
\[
 p_0(\widehat v_i)=s_Q(p_1(\widehat\alpha_i))=s_Q(\alpha_i)=v_i.
\]

Let us prove uniqueness of the whole lift.
At each step, the appropriate covering bijection determines $\widehat\alpha_i$, and its other endpoint determines $\widehat v_i$.
Induction on $i$ proves existence and uniqueness through the last vertex.
\end{proof}

\begin{definition}[Lift of a traversal]\label{def:unique-path-lift}
Let $p:\widehat Q\to Q$ be a covering and let $\gamma$ be a traversal in $Q$ beginning at $v_0$.
For $\widehat v_0\in p_0^{-1}(\{v_0\})$, let
\[
 \operatorname{Lift}_p(\gamma;\widehat v_0)
\]
denote the unique lift of $\gamma$ beginning at $\widehat v_0$ from Lemma~\ref{lem:path-lift}.
\end{definition}

\begin{definition}[Supported dimension vectors]\label{def:cover-dimension-vectors}
Let $Q$ be a quiver.
\begin{enumerate}[font=\itshape]
\item A \textit{dimension vector} is a family $d=(d_v)_{v\in\operatorname{Vert}(Q)}$
of nonnegative integers. We define
\[
 \operatorname{VecSupp}(d)=\{v\in\operatorname{Vert}(Q):d_v>0\}.
\]
We call $d$ \textit{finitely supported} if $\operatorname{VecSupp}(d)$ is finite and let
\[
 \operatorname{DimVec}_{\mathrm{fin}}(Q)
 =\{d\in\Z_{\ge0}^{\operatorname{Vert}(Q)}:\#\operatorname{VecSupp}(d)<\infty\}.
\]
Its \textit{positive support quiver} is $\operatorname{SubQuiv}(Q,\operatorname{VecSupp}(d))$.
\item Let $p:\widehat Q\to Q$ be a covering and let $\delta\in\operatorname{Deck}(p)$.
For $d\in\operatorname{DimVec}_{\mathrm{fin}}(\widehat Q)$, define
\[
 (\delta\cdot d)_v=d_{\delta_0^{-1}(v)}
 \qquad(v\in\operatorname{Vert}(\widehat Q)).
\]
We call $d$ and $d'$ \textit{deck-equivalent} if $d'=\delta\cdot d$ for some $\delta\in\operatorname{Deck}(p)$.
We denote the resulting orbit by $[d]_p$.
\end{enumerate}
\end{definition}

\begin{definition}[Covering arrow coloring]\label{def:covering-arrow-coloring}
Fix a bijection $c_3:\operatorname{Edge}(K(3))\to[0,2]_{\Z}$ and a universal covering
$\varpi=(\varpi_0,\varpi_1):\widetilde K_3\to K(3)$ for the rest of this section.
\begin{enumerate}[font=\itshape]
\item Let $I_{K(3)}$ consist of the first vertex and $J_{K(3)}$ of the second vertex of $K(3)$.
We use the bipartition
\[
 I_{\widetilde K_3}=\varpi_0^{-1}(I_{K(3)}),\qquad
 J_{\widetilde K_3}=\varpi_0^{-1}(J_{K(3)}).
\]
\item We define the \textit{covering arrow coloring} by
\[
 \operatorname{CovCol}_{\varpi}:\operatorname{Edge}(\widetilde K_3)\longrightarrow[0,2]_{\Z},\qquad
 \operatorname{CovCol}_{\varpi}(\alpha)=c_3(\varpi_1(\alpha)).
\]
\end{enumerate}
\end{definition}

\begin{definition}[Cover dimension classes and stability weights]\label{def:cover-dimension-classes}
Let $a\in\Z_{>0}$ and $b\in\Z_{>0}$. Let
$I=I_{\widetilde K_3}$ and $J=J_{\widetilde K_3}$.
\begin{enumerate}[font=\itshape]
\item We define
\[
 \operatorname{CovDim}_{\varpi}(a,b)=
 \left\{d\in\operatorname{DimVec}_{\mathrm{fin}}(\widetilde K_3):
 \sum_{i\in I}d_i=a,\ \sum_{j\in J}d_j=b\right\}.
\]
\item For $e\in\operatorname{DimVec}_{\mathrm{fin}}(\widetilde K_3)$, define
\[
 \theta^{\varpi}_{a,b}(e)=a\sum_{j\in J}e_j-b\sum_{i\in I}e_i.
\]
\item For $k\in\Z_{>0}$, we define the sets of \textit{cover dimension classes}
and \textit{stable cover dimension classes} by
\[
\begin{aligned}
 \mathcal D_{a,k}&=\{[d]_{\varpi}:d\in\operatorname{CovDim}_{\varpi}(a,a+k)\},\\
 \mathcal D^{\mathrm{st}}_{a,k}&=\{[d]_{\varpi}\in\mathcal D_{a,k}:
 \cM^s_d(\widetilde K_3)\ne\varnothing\}.
\end{aligned}
\]
The moduli spaces here use the weight $\theta^{\varpi}_{a,a+k}$.
\end{enumerate}
\end{definition}

\subsection{Counting cover representations}
\begin{definition}\label{def:marked-chain-data}
Let $a\in\Z_{>0}$ and $k\in\Z_{>0}$ satisfy $k\ge2$ and $a\ge100000k$.
Let $c=k-1$ and $S=\operatorname{LSTrivCh}(c)$.
Let $\mathbb S=(S,(I_S,J_S))$.
Define
\[
 L=J_{S}\cap\operatorname{Leaf}(S),\qquad
 I=\{i\in J_{S}:\deg_S(i)=2\}.
\]
We call a tuple
\[
 M=(S,f,\tau,(B_i)_{i\in L\sqcup I})
\]
\textit{marked chain data} if the following conditions hold.
\begin{hypotheses}
\item\label{hyp:marked-flows} Let $f\in\mathcal F_{a,k}$ and let $a_i=a_i(f)$ for $i\in J_S$.
\item\label{hyp:marked-colors} For each $t\in[c]_{\Z}$, the map
$\tau_t\colon[3]_{\Z}\longrightarrow[0,2]_{\Z}$ is a bijection.
The indices $1$, $2$, and $3$ denote $\operatorname{LeftEdge}_c(t)$,
$\operatorname{MidEdge}_c(t)$, and $\operatorname{RightEdge}_c(t)$, respectively.
\item\label{hyp:marked-blocks} For each $i\in L\sqcup I$, let $(\mathbb T_i,\kappa_i)$ be a colored source-trivalent tree.
Let $\mathbb T_i=(T_i,(I_{T_i},J_{T_i}))$ for $i\in L\sqcup I$.
Require the following conditions.
\begin{enumerate}[label=(\roman*),font=\itshape]
\item For $i\in L$, let $\mathbf s_i=(s_{i,1})$ satisfy
\[
 s_{i,1}\in\operatorname{ColVacSlot}(Q(\mathbb T_i,\kappa_i)).
\]
For $i\in I$, let $\mathbf s_i=(s_{i,1},s_{i,2})$ satisfy
\[
 s_{i,j}\in\operatorname{ColVacSlot}(Q(\mathbb T_i,\kappa_i))\quad(j\in[2]_{\Z}),
 \qquad s_{i,1}\ne s_{i,2}.
\]
\item For $i\in L\sqcup I$, let
\[
 B_i=\operatorname{MarkedCIsom}_{\C}(Q(\mathbb T_i,\kappa_i),\operatorname{ZS}(Q(\mathbb T_i,\kappa_i));\mathbf s_i)
\]
and require
\[
 \operatorname{TotDim}_{\C,(I_{Q(\mathbb T_i,\kappa_i)},J_{Q(\mathbb T_i,\kappa_i)})}
 (\operatorname{ZS}(Q(\mathbb T_i,\kappa_i)))=(a_i,a_i+1).
\]
\item Define the colored-vacant slots $p_e$ for $e\in\operatorname{Edge}(S)$ by
\[
 p_e=s_{i,1}\quad(i\in L,\ e\in\operatorname{IncEdge}_S(i)),
\]
and, for $t\in[c-1]_{\Z}$, let
\[
 p_{\operatorname{RightEdge}_c(t)}=s_{z_t,1},\qquad
 p_{\operatorname{LeftEdge}_c(t+1)}=s_{z_t,2}.
\]
\item For $e\in\operatorname{Edge}(S)$, let $p_e=(v_e,d_e)$.
For every $t\in[c]_{\Z}$, require
\[
 d_{\operatorname{LeftEdge}_c(t)}=\tau_t(1),\qquad
 d_{\operatorname{MidEdge}_c(t)}=\tau_t(2),\qquad
 d_{\operatorname{RightEdge}_c(t)}=\tau_t(3).
\]
\end{enumerate}
\end{hypotheses}
\begin{enumerate}[font=\itshape]
\item\label{d:marked-chain-set} Let $\mathcal H_{a,k}$ be the set of tuples $M$ satisfying assumptions~\ref{hyp:marked-flows}--\ref{hyp:marked-blocks}.
For a family $\tau=(\tau_t)_{t\in[c]_{\Z}}$ of bijections $\tau_t:[3]_{\Z}\to[0,2]_{\Z}$, define
\[
 I_=(\tau)=\{z_t:t\in[c-1]_{\Z},\ \tau_t(3)=\tau_{t+1}(1)\},\qquad
 I_{\ne}(\tau)=I\setminus I_=(\tau).
\]
\item\label{d:marked-chain-support} For each $M\in\mathcal H_{a,k}$, fix a family
\[
 \rho_M=((\mathbb T_i,\kappa_i;\mathbf s_i))_{i\in J_S}
\]
satisfying assumption~\ref{hyp:marked-blocks} and defining the classes $B_i$ in $M$.
Let $\rho=(\rho_M)_{M\in\mathcal H_{a,k}}$.
For $M\in\mathcal H_{a,k}$, let
$\mathbf T_M^\rho=((\mathbb T_i,\kappa_i))_{i\in J_S}$,
$\mathbf p_M^\rho=(p_e)_{e\in\operatorname{Edge}(S)}$, and
$\mathbf 1=(1_{\C})_{e\in\operatorname{Edge}(S)}$.
Define
\[
 R_M^\rho=\operatorname{Glue}_{\C}(\mathbb S;\mathbf T_M^\rho,\mathbf p_M^\rho,\mathbf 1),\qquad
 \widetilde Q_M^\rho=\operatorname{SuppColQuiv}(R_M^\rho),\qquad
 d_M^\rho(v)=\dim_{\C}(R_M^\rho)_v.
\]
Use the tagged vertices $(i,v)$ for the copy of $Q(\mathbb T_i,\kappa_i)$ and distinct new vertices $x_t$ for $t\in[c]_{\Z}$.
Define
\[
\begin{aligned}
 I_M^\rho&=\operatorname{Vert}(\widetilde Q_M^\rho)\cap
 \left(\{x_t:t\in[c]_{\Z}\}\sqcup
 \bigsqcup_{i\in J_S}\{(i,v):v\in I_{Q(\mathbb T_i,\kappa_i)}\}\right),\\
 J_M^\rho&=\operatorname{Vert}(\widetilde Q_M^\rho)\cap
 \bigsqcup_{i\in J_S}\{(i,v):v\in J_{Q(\mathbb T_i,\kappa_i)}\},\\
 \widetilde{\mathbb Q}_M^\rho&=(\widetilde Q_M^\rho,(I_M^\rho,J_M^\rho)).
\end{aligned}
\]
Let $\iota_i^\rho:Q(\mathbb T_i,\kappa_i)\hookrightarrow\widetilde{\mathbb Q}_M^\rho$ denote the inclusion of the copy indexed by $i$.
\item\label{d:dimensioned-support-classes} An isomorphism between colored bipartite quivers $\widetilde{\mathbb G}=(G,(I_G,J_G))$ and $\widetilde{\mathbb G}'=(G',(I_{G'},J_{G'}))$ equipped with positive vertex dimensions $d_G$ and $d_{G'}$ is a color-preserving quiver isomorphism $\phi=(\phi_0,\phi_1)$ that preserves the bipartitions and satisfies
\[
 d_{G'}(\phi_0(v))=d_G(v)\qquad(v\in\operatorname{Vert}(G)).
\]
Define
\[
 \mathcal S_{a,k}^\rho=\{[(\widetilde{\mathbb Q}_M^\rho,d_M^\rho)]:M\in\mathcal H_{a,k}\},\qquad
 \pi_\rho:\mathcal H_{a,k}\longrightarrow\mathcal S_{a,k}^\rho,\quad
 \pi_\rho(M)=[(\widetilde{\mathbb Q}_M^\rho,d_M^\rho)].
\]

\item\label{d:connector-set} Let $X=(\widetilde{\mathbb G},d_G)$ represent a class in $\mathcal S_{a,k}^\rho$.
Let $\widetilde{\mathbb G}=(G,(I_G,J_G))$.
Define
\[
 C(X)=\{v\in I_G:d_G(v)=1,\ \deg_G(v)=3\}.
\]
\item\label{d:deleted-components} Define
\[
 \mathcal K(X)=\operatorname{CtdComp}\bigl(\operatorname{Subgraph}(\operatorname{UnderG}(G),\operatorname{Vert}(G)\setminus C(X))\bigr).
\]
\item\label{d:recovered-skeleton} Define the bipartite graph $\mathbb S_X=(S_X,(I_{S_X},J_{S_X}))$ with vertex set $C(X)\sqcup\mathcal K(X)$, identifying each part with its canonical image, by
\[
 I_{S_X}=C(X),\qquad
 J_{S_X}=\mathcal K(X).
\]
For $v\in C(X)$ and $K\in\mathcal K(X)$, define its edges by
\[
 \llbracket v,K\rrbracket\in\operatorname{Edge}(S_X)
 \ \Longleftrightarrow\
 \exists u\in\operatorname{Vert}(K),\quad
 \llbracket v,u\rrbracket\in\operatorname{Edge}(\operatorname{UnderG}(G)).
\]
\end{enumerate}
\end{definition}

\begin{lemma}[Recovery of a colored tree and its markings]\label{lem:retained-tree-recovery}
Let $(\mathbb T,\kappa)$ and $(\mathbb T',\kappa')$ be colored source-trivalent trees with $\#I_T\ge1$ and $\#I_{T'}\ge1$.
Let $Q=Q(\mathbb T,\kappa)$ and $Q'=Q(\mathbb T',\kappa')$.
Let $q=(q_0,q_1):Q\to Q'$ be a color-preserving isomorphism preserving the bipartitions.
Then, we have the following.
\begin{enumerate}
\item\label{c:retained-tree-extension} There is a unique color-preserving isomorphism
\[
 \Psi=(\Psi_0,\Psi_1):(\mathbb T,\kappa)\longrightarrow(\mathbb T',\kappa')
\]
preserving the bipartitions and satisfying
$\Psi_0(v)=q_0(v)$ for $v\in\operatorname{Vert}(Q)$.
\item\label{c:retained-marked-transport} Let $\ell\in\Z_{\ge0}$ and let
$\mathbf s=((u_j,d_j))_{j\in[\ell]_{\Z}}$ be an ordered marking of colored-vacant slots of $Q$.
Let $\mathbf s'=((q_0(u_j),d_j))_{j\in[\ell]_{\Z}}$.
Then, $\mathbf s'$ is an ordered marking of $Q'$ and
\[
 \operatorname{MarkedCIsom}_{\C}(Q,\operatorname{ZS}(Q);\mathbf s)
 =\operatorname{MarkedCIsom}_{\C}(Q',\operatorname{ZS}(Q');\mathbf s').
\]
\end{enumerate}
\end{lemma}
\begin{proof}
\prf{Claim~\ref{c:retained-tree-extension}} Let us define the images of the deleted leaves.
Let
\[
 L=\{v\in J_T:\deg_T(v)=1\},\qquad
 L'=\{v'\in J_{T'}:\deg_{T'}(v')=1\}.
\]
Connectedness, $\#I_T\ge1$, and $\#I_{T'}\ge1$ imply
\[
 \operatorname{Vert}(T)=\operatorname{Vert}(Q)\sqcup L,\qquad
 \operatorname{Vert}(T')=\operatorname{Vert}(Q')\sqcup L'.
\]
For $v\in L$, let $u\in I_T$ be its unique neighbor and let
$d=\kappa(\llbracket u,v\rrbracket)$.
Lemma~\ref{lem:vacant-slots} gives $(u,d)\in\operatorname{ColVacSlot}(Q)$.
Since $q$ preserves incidence and colors,
\[
 (u,d)\in\operatorname{ColVacSlot}(Q)
 \quad\Longleftrightarrow\quad
 (q_0(u),d)\in\operatorname{ColVacSlot}(Q').
\]
Lemma~\ref{lem:vacant-slots} gives a unique $v'\in L'$ satisfying
\[
 \llbracket q_0(u),v'\rrbracket\in\operatorname{Edge}(T'),\qquad
 \kappa'(\llbracket q_0(u),v'\rrbracket)=d.
\]
Define $\eta(v)=v'$.
Applying the construction of $\eta:L\to L'$ to $q^{-1}:Q'\to Q$ gives $\eta':L'\to L$ satisfying
\[
 \eta'(\eta(v))=v\quad(v\in L),\qquad
 \eta(\eta'(v'))=v'\quad(v'\in L').
\]
Define
\[
 \Psi_0(v)=
 \begin{cases}
 q_0(v)&v\in\operatorname{Vert}(Q),\\
 \eta(v)&v\in L,
 \end{cases}
 \qquad
 \Psi_1(\llbracket u,v\rrbracket)
 =\llbracket\Psi_0(u),\Psi_0(v)\rrbracket.
\]
Let $e=\llbracket u,v\rrbracket\in\operatorname{Edge}(T)$ with $u\in I_T$ and $v\in J_T\setminus L$.
The arrow $v\to u$ belongs to $Q$.
Since $q$ preserves incidence and colors, we have
\[
 \Psi_1(e)=\llbracket q_0(u),q_0(v)\rrbracket\in\operatorname{Edge}(T'),\qquad
 \kappa'(\Psi_1(e))=\kappa(e).
\]
For $v\in L$, the construction of $\eta(v)$ gives $\Psi_1(e)\in\operatorname{Edge}(T')$ and $\kappa'(\Psi_1(e))=\kappa(e)$.
The maps defined from $q^{-1}$ and $\eta'$ are inverse to $\Psi_0$ and $\Psi_1$.
Thus, $\Psi$ is a color-preserving, bipartition-preserving isomorphism extending $q$.
Any such isomorphism must send a leaf $v$ with neighbor $u$ and color $d$ to the unique leaf of color $d$ adjacent to $q_0(u)$.
Such an isomorphism therefore agrees with $\Psi$ on all vertices and edges.

\prf{Claim~\ref{c:retained-marked-transport}} Let us give the linear maps on the vertex spaces.
The colored-vacant-slot equivalence in Claim~\ref{c:retained-tree-extension} shows that $\mathbf s'$ is a marking of $Q'$.
For $v\in I_Q$, put
\[
 A_v=\operatorname{IncEdgeQuiv}_Q(v),\qquad
 A'_{q_0(v)}=\operatorname{IncEdgeQuiv}_{Q'}(q_0(v)).
\]
The quiver isomorphism $q$ restricts to a bijection
$q_1:A_v\to A'_{q_0(v)}$.
Define $h_u=\operatorname{id}_{\C}$ for $u\in J_Q$.
For $v\in I_Q$ and $z\in\operatorname{ZS}(Q)_v$, define
\begin{equation}\label{eq:zero-sum-transport}
 (h_v(z))_\beta=z_{q_1^{-1}(\beta)}
 \qquad(\beta\in A'_{q_0(v)}).
\end{equation}

Let us prove that $h_v$ is a linear isomorphism of the zero-sum spaces.
By~\eqref{eq:zero-sum-transport} and~\eqref{eq:zero-sum-map}, we have
\[
 \sum_{\beta\in A'_{q_0(v)}}(h_v(z))_\beta
 =\sum_{\alpha\in A_v}z_\alpha=0.
\]
Thus, $h_v(z)\in\operatorname{ZS}(Q')_{q_0(v)}$.
The map $h_v$ is linear. Its inverse satisfies
\[
 (h_v^{-1}(w))_\alpha=w_{q_1(\alpha)}
 \qquad(\alpha\in A_v,\ w\in\operatorname{ZS}(Q')_{q_0(v)}).
\]

Let us check the arrow maps and the markings.
For $\alpha:v\to u$ and $z\in\operatorname{ZS}(Q)_v$, equations~\eqref{eq:zero-sum-arrow} and~\eqref{eq:zero-sum-transport} give
\[
 \operatorname{ZS}(Q')_{q_1(\alpha)}(h_v(z))
 =(h_v(z))_{q_1(\alpha)}
 =z_\alpha
 =h_u(\operatorname{ZS}(Q)_\alpha(z)).
\]
Hence, $(q_0,q_1,(h_v)_{v\in\operatorname{Vert}(Q)})$ is a representation isomorphism.
For each $j\in[\ell]_{\Z}$, its image of the $j$th marking is $(q_0(u_j),d_j)$.
Hence, $(q_0,q_1,(h_v)_{v\in\operatorname{Vert}(Q)})$ is an isomorphism of the marked representations.
\end{proof}

\begin{lemma}[Recovery from a labeled support]\label{lem:marked-chain-recovery}
Let $a\in\Z_{>0}$ and $k\in\Z_{>0}$ satisfy $k\ge2$ and $a\ge100000k$.
Fix a family of representatives $\rho=(\rho_M)_{M\in\mathcal H_{a,k}}$.
Let $c=k-1$ and $S=\operatorname{LSTrivCh}(c)$.
Let $\mathbb S=(S,(I_S,J_S))$.
For $M=(S,f,\tau,(B_i)_i)\in\mathcal H_{a,k}$, let
$\rho_M=((\mathbb T_i,\kappa_i;\mathbf s_i))_{i\in J_S}$.
Let $\iota_i=\iota_i^\rho$ for $i\in J_S$.

Then, we have the following.
\begin{enumerate}
\item\label{c:marked-representative-independence} For any other family of representatives $\rho'$, we have
\[
 \mathcal S_{a,k}^\rho=\mathcal S_{a,k}^{\rho'},\qquad
 \pi_\rho=\pi_{\rho'}.
\]
\item\label{c:marked-recovery-unique} Assume the following.
\begin{hypotheses}
\item\label{hyp:recovery-support} $H\in\mathcal S_{a,k}^\rho$ and $X=(\widetilde{\mathbb G},d_G)$ represents $H$.
Let $\widetilde{\mathbb G}=(G,(I_G,J_G))$.
\item\label{hyp:recovery-labeling} $\lambda=(\lambda_0,\lambda_1):S\to S_X$ is an isomorphism preserving the bipartition.
\end{hypotheses}
Then, there is at most one $M\in\pi_\rho^{-1}(\{H\})$ admitting an isomorphism
$\phi:(\widetilde{\mathbb Q}_M^\rho,d_M^\rho)\to X$ such that
\begin{equation}\label{eq:recovery-matching}
\begin{aligned}
 \phi_0(x_t)&=\lambda_0(x_t)&& (t\in[c]_{\Z}),\\
 \phi_0\bigl(\iota_{i,0}(\operatorname{Vert}(Q(\mathbb T_i,\kappa_i)))\bigr)
 &=\operatorname{Vert}(\lambda_0(i))&& (i\in J_{S}).
\end{aligned}
\end{equation}
\item\label{c:marked-recovery-skeleton} Let $H\in\mathcal S_{a,k}^\rho$ and let $X=(\widetilde{\mathbb G},d_G)$ represent $H$.
For $M\in\pi_\rho^{-1}(\{H\})$, every isomorphism
$\phi:(\widetilde{\mathbb Q}_M^\rho,d_M^\rho)\to X$
induces an isomorphism $\lambda:S\to S_X$ preserving the bipartition and satisfying~\eqref{eq:recovery-matching}.
\end{enumerate}
\end{lemma}
\begin{proof}
\prf{Claim~\ref{c:marked-representative-independence}} Let us compare the two representatives of each $B_i$.
Fix $M=(S,f,\tau,(B_i)_i)\in\mathcal H_{a,k}$.
The equality of the marked classes gives isomorphisms
\[
 \Phi_i=(\phi_{i,0},\phi_{i,1},(h_{i,v})_v)
 \quad(i\in J_S)
\]
between the marked representations in $\rho_M$ and $\rho'_M$.
They preserve the bipartitions, the arrow colors, and the ordered marked colored-vacant slots.
Define maps on the glued vertices and arrows by
\[
\begin{aligned}
 \psi_0(x_t)&=x_t &&(t\in[c]_{\Z}),\\
 \psi_0(i,v)&=(i,\phi_{i,0}(v)) &&(i\in J_S,\ v\in\operatorname{Vert}(Q(\mathbb T_i,\kappa_i))),\\
 \psi_1(\alpha_e)&=\alpha'_e &&(e\in\operatorname{Edge}(S)),\\
 \psi_1(i,\alpha)&=(i,\phi_{i,1}(\alpha)) &&(i\in J_S,\ \alpha\in\operatorname{Edge}(Q(\mathbb T_i,\kappa_i))).
\end{aligned}
\]
For an added arrow $\alpha_e:x_t\to(i,v_e)$, preservation of its marked colored-vacant slot gives
$\alpha'_e:x_t\to(i,\phi_{i,0}(v_e))$ with the same color $d_e$.
Thus, $\psi=(\psi_0,\psi_1)$ is a color-preserving quiver isomorphism that preserves the bipartitions.
The maps $h_{i,v}$ give
$d_M^{\rho'}(\psi_0(i,v))=d_M^\rho(i,v)$.
At each $x_t$, both dimensions are one.
Restricting $\psi$ to the positive supports gives
\[
 (\widetilde{\mathbb Q}_M^\rho,d_M^\rho)
 \cong(\widetilde{\mathbb Q}_M^{\rho'},d_M^{\rho'}).
\]
Hence, $\pi_\rho(M)=\pi_{\rho'}(M)$ for every $M\in\mathcal H_{a,k}$.
Their image sets are therefore equal.

\prf{Claim~\ref{c:marked-recovery-unique}} Let us identify the components determined by $\lambda$.
Let $M\in\pi_\rho^{-1}(\{H\})$ and let $\phi:(\widetilde{\mathbb Q}_M^\rho,d_M^\rho)\to X$ be an isomorphism satisfying~\eqref{eq:recovery-matching}.
For $i\in J_{S}$, let
\[
\begin{gathered}
 K_i=\lambda_0(i),\qquad
 \widetilde Q_i=\operatorname{SubColQuiv}(G,\operatorname{Vert}(K_i)),\\
 I_i=I_G\cap\operatorname{Vert}(K_i),\qquad
 J_i=J_G\cap\operatorname{Vert}(K_i),\\
 \widetilde{\mathbb Q}_i=(\widetilde Q_i,(I_i,J_i)).
\end{gathered}
\]
Assumption~\ref{hyp:recovery-labeling} and~\eqref{eq:recovery-matching} give
$\widetilde{\mathbb Q}_i\cong Q(\mathbb T_i,\kappa_i)$ by a color-preserving isomorphism taking $(I_i,J_i)$ to $(I_{Q(\mathbb T_i,\kappa_i)},J_{Q(\mathbb T_i,\kappa_i)})$.

Let us recover the colored-vacant slots of the added arrows.
For $e=\llbracket x_t,i\rrbracket\in\operatorname{Edge}(S)$, there is a unique arrow
\[
 \alpha_e:\lambda_0(x_t)\longrightarrow u_e,\qquad
 u_e\in\operatorname{Vert}(K_i).
\]
Existence follows from the definition of $S_X$ and assumption~\ref{hyp:recovery-labeling}.
If there were two such arrows, their edges together with the path inside $K_i$ would form a cycle in $\operatorname{UnderG}(G)$.
This contradicts Claim~\ref{c:flow-stability-support-tree} of Lemma~\ref{lem:flow-stability}.
Define
\[
 d_e=\operatorname{QCol}(G)(\alpha_e),\qquad p_e=(u_e,d_e).
\]
Proper coloring of $G$ gives $p_e\in\operatorname{ColVacSlot}(\widetilde{\mathbb Q}_i)$.

Let us recover the color permutations and the ordered markings.
For $t\in[c]_{\Z}$, the color permutation is
\[
 \bigl(\tau_t(1),\tau_t(2),\tau_t(3)\bigr)
 =\bigl(d_{\operatorname{LeftEdge}_c(t)},
 d_{\operatorname{MidEdge}_c(t)},d_{\operatorname{RightEdge}_c(t)}\bigr).
\]
For $i\in L$, let $e_i$ be its unique incident edge. Its marking is $(p_{e_i})$.
For $t\in[c-1]_{\Z}$, its ordered marking is
\[
 \mathbf s_{z_t}
 =\bigl(p_{\operatorname{RightEdge}_c(t)},p_{\operatorname{LeftEdge}_c(t+1)}\bigr).
\]
These are exactly the orders required by assumption~\ref{hyp:marked-blocks} of Definition~\ref{def:marked-chain-data}.

Let us recover the marked zero-sum representations.
Suppose $M'=(S,f',\tau',(B'_i)_i)\in\pi_\rho^{-1}(\{H\})$ admits an isomorphism $\phi'$ satisfying~\eqref{eq:recovery-matching}.
Let $((\mathbb T'_i,\kappa'_i;\mathbf s'_i))_i=\rho_{M'}$ and let $\iota'_i$ be its inclusion maps.
For $i\in J_S$, let
\[
 Q_i=Q(\mathbb T_i,\kappa_i),\qquad Q'_i=Q(\mathbb T'_i,\kappa'_i),\qquad
 q_i=(\phi'\circ\iota'_i)^{-1}\circ(\phi\circ\iota_i):Q_i\longrightarrow Q'_i.
\]
Both maps in this composition are isomorphisms onto $\widetilde{\mathbb Q}_i$.
For each edge $e$ incident to $i$, let $(v_{i,e},d_e)$ and $(v'_{i,e},d_e)$ be the colored-vacant slots assigned to $e$ in the two representatives.
Uniqueness of $\alpha_e$ gives
\[
 (q_{i,0}(v_{i,e}),d_e)=(v'_{i,e},d_e).
\]
The ordered markings agree under $q_i$ because their orders are determined by $S$.

Let us recover each class $B_i$ for $i\in J_S$.
Assumption~\ref{hyp:marked-blocks} of Definition~\ref{def:marked-chain-data} gives
\[
 \#I_{T_i}=a_i(f)+1\ge1,\qquad
 \#I_{T'_i}=a_i(f')+1\ge1.
\]
Claim~\ref{c:retained-tree-extension} of Lemma~\ref{lem:retained-tree-recovery} recovers an isomorphism of these colored trees.
Claim~\ref{c:retained-marked-transport} of Lemma~\ref{lem:retained-tree-recovery} gives
\[
 B_i=B'_i\qquad(i\in J_S).
\]
Thus, $G$, $d_G$, and $\lambda$ determine every $B_i$.

Let us recover the flow coordinates.
For each $i\in J_{S}$, its source dimension is
\[
 a_i=\sum_{v\in I_i}d_G(v).
\]
For $e\in\operatorname{Edge}(S)$, put $S_e=\operatorname{SnkCtdComp}(\mathbb S;e)$.
Claim~\ref{c:tree-flow-cut} of Lemma~\ref{lem:tree-flow-cut} gives
\[
 f_e=\sum_{i\in \operatorname{SubSnk}_{\mathbb S}(S_e)}\left(\frac ak-a_i\right)
      -\left(\frac ak+1\right)\#\operatorname{SubSrc}_{\mathbb S}(S_e).
\]
In particular,
\[
 l_t=f_{\operatorname{LeftEdge}_c(t)},\qquad
 r_t=f_{\operatorname{RightEdge}_c(t)}\qquad(t\in[c]_{\Z}).
\]
Therefore, $G$, $d_G$, and $\lambda$ determine every entry of $M=(S,f,\tau,(B_i)_{i\in L\sqcup I})$.

\prf{Claim~\ref{c:marked-recovery-skeleton}} Let us identify the sources in $C(X)$.
Claim~\ref{c:flow-choices-range} of Lemma~\ref{lem:flow-choices} and assumption~\ref{hyp:marked-blocks} of Definition~\ref{def:marked-chain-data} give $\#I_{T_i}=a_i(f)+1\ge2$ for $i\in J_S$.
For each $t\in[c]_{\Z}$, the new source satisfies
\[
 d_M^\rho(x_t)=1,\qquad \deg_{\widetilde{\mathbb Q}_M^\rho}(x_t)=3.
\]
For $i\in J_S$ and $v\in I_{Q(\mathbb T_i,\kappa_i)}$, the added arrows do not meet the copied source $(i,v)$.
Thus,
\[
 \deg_{\widetilde{\mathbb Q}_M^\rho}(i,v)
 =\deg_{Q(\mathbb T_i,\kappa_i)}(v).
\]
Claim~\ref{c:reflected-blocks-sources} of Lemma~\ref{lem:reflected-blocks} gives
\[
 d_M^\rho(i,v)=1\quad\Longrightarrow\quad
 \deg_{\widetilde{\mathbb Q}_M^\rho}(i,v)=2.
\]
Consequently,
\[
 C(X)=\{\phi_0(x_t):t\in[c]_{\Z}\}.
\]

Let us identify the components after deleting $C(X)$.
Each $Q(\mathbb T_i,\kappa_i)$ is connected because it is obtained from $T_i$ by deleting its sink leaves.
By Definition~\ref{def:attachments}, deleting the new sources $x_t$ removes every added arrow and leaves exactly the disjoint copied quivers.
Thus, for each $i\in J_S$, there is a unique $K_i\in\mathcal K(X)$ satisfying
\[
 \operatorname{Vert}(K_i)=\phi_0\bigl(\iota_{i,0}(\operatorname{Vert}(Q(\mathbb T_i,\kappa_i)))\bigr).
\]
Define $\lambda_0(x_t)=\phi_0(x_t)$ for $t\in[c]_{\Z}$ and $\lambda_0(i)=K_i$ for $i\in J_S$.
These maps are bijections onto $C(X)$ and $\mathcal K(X)$, respectively.

Let us verify incidence.
The added arrows in Definition~\ref{def:attachments} and Definition~\ref{def:marked-chain-data}(\ref{d:recovered-skeleton}) give
\[
 \llbracket x_t,i\rrbracket\in\operatorname{Edge}(S)
 \quad\Longleftrightarrow\quad
 \llbracket\phi_0(x_t),K_i\rrbracket\in\operatorname{Edge}(S_X).
\]
Hence, $\lambda_0$ induces the required isomorphism $\lambda:S\to S_X$.
\end{proof}

\begin{lemma}[Marked chains and their supports]\label{lem:forget-markings}
Let $a\in\Z_{>0}$ and $k\in\Z_{>0}$ satisfy $k\ge2$ and $a\ge100000k$.
Fix a family of representatives $\rho=(\rho_M)_{M\in\mathcal H_{a,k}}$.
Let $c=k-1$ and $S=\operatorname{LSTrivCh}(c)$. Consider the map $\pi_\rho\colon\mathcal H_{a,k}\to\mathcal S_{a,k}^\rho$.
Let
\[
 L=J_S\cap\operatorname{Leaf}(S),\qquad I=\{i\in J_S:\deg_S(i)=2\}.
\]
Then, we have the following.
\begin{enumerate}
\item\label{c:forget-markings-eight} For $H\in\mathcal S_{a,k}^\rho$,
we have $\#\pi_\rho^{-1}(\{H\})\le8$.
\item\label{c:forget-markings-six} If $c=1$, then
$\#\pi_\rho^{-1}(\{H\})\le6$ for $H\in\mathcal S_{a,k}^\rho$.
\item\label{c:marked-color-count} There are $6^c$ families $\tau=(\tau_t)_{t\in[c]_{\Z}}$ such that each $\tau_t:[3]_{\Z}\to[0,2]_{\Z}$ is a bijection.
\item\label{c:marked-block-count} Let $f\in\mathcal F_{a,k}$.
Let $\tau=(\tau_t)_{t\in[c]_{\Z}}$ be a family of bijections $\tau_t:[3]_{\Z}\to[0,2]_{\Z}$.
Let $a_i=a_i(f)$ for $i\in J_S$.
The number of choices of $(B_i)_{i\in L\sqcup I}$ is
\[
 \prod_{i\in L}P_{a_i}
 \prod_{i\in I_=(\tau)}Q_{a_i}^{=}
 \prod_{i\in I_{\ne}(\tau)}Q_{a_i}^{\ne}.
\]
For every $(P_j,Q_j^{=},256)$-admissible pair $(A,B)$, this number is at least
\[
 \beta^{a-c}A^{c+2}B^{c-1}
 \prod_{i\in L}a_i^{-\frac32}
 \prod_{i\in I}a_i^{-\frac12}.
\]
\end{enumerate}
\end{lemma}
\begin{proof}
\prf{Claim~\ref{c:forget-markings-eight}} Let us bound the number of marked chains over one support.
Choose a representative $X=(\widetilde{\mathbb G},d_G)$ of $H$.
Let us count the bipartition-preserving isomorphisms
\[
 \mathcal I_X=\{\lambda:S\xrightarrow{\sim}S_X:
 \lambda_0(I_S)=I_{S_X},\ \lambda_0(J_S)=J_{S_X}\}.
\]

Let us associate an element of $\mathcal I_X$ to each $M\in\pi_\rho^{-1}(\{H\})$.
Choose an isomorphism
\[
 \phi_M:(\widetilde{\mathbb Q}_M^\rho,d_M^\rho)\longrightarrow X.
\]
Claim~\ref{c:marked-recovery-skeleton} of Lemma~\ref{lem:marked-chain-recovery} gives a bipartition-preserving isomorphism $\lambda_M:S\to S_X$ satisfying
\[
 \begin{aligned}
 (\lambda_M)_0(x_t)&=(\phi_M)_0(x_t)&& (t\in[c]_{\Z}),\\
 \operatorname{Vert}((\lambda_M)_0(i))
 &=(\phi_M)_0\bigl(\iota_{i,0}^\rho(\operatorname{Vert}(Q(\mathbb T_i,\kappa_i)))\bigr)
 &&(i\in J_S).
 \end{aligned}
\]
Thus, $\lambda_M\in\mathcal I_X$.
By Claim~\ref{c:marked-recovery-unique} of Lemma~\ref{lem:marked-chain-recovery}, each $\lambda\in\mathcal I_X$ determines at most one such $M$.
Hence,
\[
 \#\pi_\rho^{-1}(\{H\})\le\#\mathcal I_X.
\]

First, suppose $c\ge2$. Deleting the sink leaves of $S$ gives a path with endpoints $x_1,x_c$ and vertex set
\[
 \{x_t:t\in[c]_{\Z}\}\sqcup\{z_t:t\in[c-1]_{\Z}\}.
\]
An isomorphism sends $x_1$ to one of the two endpoints of the corresponding path in $S_X$.
For $\lambda\in\mathcal I_X$, choosing $\lambda_0(x_1)$ determines
$\lambda_0(x_t)$ for $t\in[c]_{\Z}$ and $\lambda_0(z_t)$ for $t\in[c-1]_{\Z}$.
The two sink leaves at $x_1$ are $z_0,y_1$, and those at $x_c$ are $y_c,z_c$.
Each pair has $2!$ possible permutations.
For $1<t<c$, the only sink leaf at $x_t$ is $y_t$, whose image is then determined.
Thus,
\[
 \#\pi_\rho^{-1}(\{H\})\le\#\mathcal I_X
 \le2\cdot2!\cdot2!=8.
\]

Second, suppose $c=1$. The only source is $x_1$, and its three sink neighbors are $z_0,y_1,z_1$.
An isomorphism is determined by a permutation of these sinks. Therefore,
\[
 \#\pi_\rho^{-1}(\{H\})\le\#\mathcal I_X\le3!=6\le8.
\]

\prf{Claim~\ref{c:forget-markings-six}}
The case $c=1$ in the proof of Claim~\ref{c:forget-markings-eight} gives
\[
 \#\pi_\rho^{-1}(\{H\})\le6.
\]

\prf{Claim~\ref{c:marked-color-count}} Let us count the color bijections independently at each source.
Each $\tau_t$ is a bijection between
sets of three elements. The choices are independent. Thus, their number is
$(3!)^c=6^c$.

\prf{Claim~\ref{c:marked-block-count}} Let us count the choices of zero-sum representations. Fix $f$ and $\tau$.
Definition~\ref{def:marked-block-counts}(1) gives $P_{a_i}$ choices for $i\in L$.
Parts~(2) and~(3) of that definition give $Q_{a_i}^{=}$ choices for $i\in I_=(\tau)$
and $Q_{a_i}^{\ne}$ choices for $i\in I_{\ne}(\tau)$.
The choices of zero-sum representations are independent. Thus, their counts multiply.

Let us check that the counts depend only on equality of the marked colors.
For colors $d,d'\in[0,2]_{\Z}$, choose a permutation $\sigma$ of $[0,2]_{\Z}$ such that $\sigma(d)=d'$.
For ordered pairs $(d,e)$ and $(d',e')$ satisfying
\[
 d=e\quad\Longleftrightarrow\quad d'=e',
\]
choose $\sigma$ such that $\sigma(d)=d'$ and $\sigma(e)=e'$.
Replacing $\kappa$ by $\sigma\circ\kappa$ gives a bijection of the corresponding marked classes.
Its inverse uses $\sigma^{-1}$. Thus, the counts apply to every family $\tau=(\tau_t)_{t\in[c]_{\Z}}$.

Let us obtain a uniform lower bound for the product of marked counts.
Let $(A,B)$ be a $(P_j,Q_j^{=},256)$-admissible pair.
Claim~\ref{c:slots-distinct} of Lemma~\ref{slots} gives $Q_j^{\ne}\ge Q_j^=$.
Claims~\ref{c:flow-choices-sum} and~\ref{c:flow-choices-range} of Lemma~\ref{lem:flow-choices}
give $\sum_i a_i=a-c$ and $a_i>256$.
The $(P_j,Q_j^{=},256)$-admissibility of $(A,B)$ gives the lower bound since $\#L=c+2$ and $\#I=c-1$.
\end{proof}

\begin{proposition}[Counting stable dimensioned supports]\label{prop:cover-count}
Let $k\ge2$ and $a>0$ be integers. Let $c=k-1$, $x=\frac ak$,
$N=256$, $\gamma=1-\frac Nx$, and $\lambda=\frac1{100}$.
Let $S=\operatorname{LSTrivCh}(c)$.
Assume the following.
\begin{hypotheses}
\item\label{hyp:cover-dimensions} $a\ge100000k$.
\item\label{hyp:cover-slot-bounds} Let $(A,B)$ be a $(P_j,Q_j^{=},256)$-admissible pair.
\end{hypotheses}
Fix a family of representatives $\rho=(\rho_M)_{M\in\mathcal H_{a,k}}$.
Then,
\begin{equation}
\#\mathcal S_{a,k}^\rho\ge\beta^a x^{-\frac{5}{2}}\frac{A^2}{8B\lambda N^2}
 \left(\frac{6AB\lambda\gamma^2}{\beta}\right)^c.\label{count}
\end{equation}
\end{proposition}
\begin{proof}
Let us count the marked chains for a fixed color family.
Assumption~\ref{hyp:cover-dimensions} gives $x\ge100000$.
Fix a family $\tau=(\tau_t)_{t\in[c]_{\Z}}$ of bijections
$\tau_t:[3]_{\Z}\to[0,2]_{\Z}$.
Claim~\ref{c:flow-choices-count} of Lemma~\ref{lem:flow-choices} gives
\[
 \#\mathcal F_{a,k}\ge\left(\frac{\gamma x}{N}\right)^{2c}.
\]
Claim~\ref{c:flow-choices-distinct} of Lemma~\ref{lem:flow-choices}
gives distinct labeled size lists for distinct $f\in\mathcal F_{a,k}$.
Apply Claim~\ref{c:marked-block-count} of Lemma~\ref{lem:forget-markings}
and Claim~\ref{c:flow-choices-weight} of Lemma~\ref{lem:flow-choices}.
The $c-1$ pairs in~\eqref{weight} and the leaves of sizes
$x-l_1$, $x-r_c$, and $l_c+r_c-1$ give
\[
\begin{aligned}
 &\#\{(S,f,\tau,(B_i)_i)\in\mathcal H_{a,k}\}\\
 &\quad\ge\left(\frac{\gamma x}{N}\right)^{2c}
   \beta^{a-c}A^{c+2}B^{c-1}
   (\lambda N^2x^{-2})^{c-1}x^{-\frac92}\\
 &\quad=\beta^{a-c}A^{c+2}B^{c-1}\gamma^{2c}x^{-\frac52}N^{-2}\lambda^{c-1}.
\end{aligned}
\]
For $c=1$, the product over pairs is empty and $B^{c-1}\lambda^{c-1}=1$.

Let us sum over color families and pass to isomorphism classes.
Claim~\ref{c:marked-color-count} of Lemma~\ref{lem:forget-markings}
gives $6^c$ choices of $\tau$.
Claim~\ref{c:forget-markings-eight} of Lemma~\ref{lem:forget-markings} gives
\[
\begin{aligned}
 \#\mathcal S_{a,k}^\rho
 &\ge\frac18\#\mathcal H_{a,k}\\
 &\ge\frac{6^c}{8}\beta^{a-c}A^{c+2}B^{c-1}
 \gamma^{2c}x^{-\frac52}N^{-2}\lambda^{c-1}.
\end{aligned}
\]
Substitute $\beta^{a-c}=\beta^a\beta^{-c}$ to obtain~\eqref{count}.
\end{proof}

\begin{lemma}[Traversals without immediate reversal in a tree]\label{lem:tree-traversal-distinct}
Let $G$ be a tree and let $n\in\Z_{\ge0}$.
Let $v_i\in\operatorname{Vert}(G)$ for $i\in[0,n]_{\Z}$ and suppose
\[
 e_i=\llbracket v_{i-1},v_i\rrbracket\in\operatorname{Edge}(G)
 \quad(i\in[n]_{\Z}),\qquad
 e_i\ne e_{i+1}\quad(i\in[n-1]_{\Z}).
\]
Then, $v_p\ne v_q$ for $0\le p<q\le n$.
\end{lemma}
\begin{proof}
Let us suppose $v_p=v_q$ for some $0\le p<q\le n$.
For $i\in[p,q]_{\Z}$, let $h_i$ be the length of the unique simple path from $v_p$ to $v_i$.
Then,
\[
 h_p=h_q=0,\qquad h_{p+1}=1.
\]
In particular, $q\ge p+2$.
Choose $j\in[p+1,q-1]_{\Z}$ satisfying
\[
 h_j=\max_{i\in[p,q]_{\Z}}h_i.
\]
The unique simple paths from $v_p$ to adjacent vertices differ by one edge.
Thus, $|h_{i+1}-h_i|=1$ for $p\le i<q$.
Maximality of $h_j$ gives
\[
 h_{j-1}=h_j-1=h_{j+1}.
\]
Both $v_{j-1}$ and $v_{j+1}$ are the predecessor of $v_j$ on the unique simple path from $v_p$ to $v_j$.
Consequently,
\[
 v_{j-1}=v_{j+1},\qquad
 e_j=\llbracket v_{j-1},v_j\rrbracket
 =\llbracket v_j,v_{j+1}\rrbracket=e_{j+1}.
\]
This contradicts $e_j\ne e_{j+1}$.
\end{proof}

\begin{lemma}[Embedding a colored tree]\label{lem:cover-tree-embedding}
Consider the universal covering
\[
 \varpi:\widetilde K_3\longrightarrow K(3)
\]
with arrow coloring $\operatorname{CovCol}_{\varpi}=c_3\circ\varpi_1$.
Assume the following.
\begin{hypotheses}
\item\label{hyp:embedding-tree} $(T,(I_T,J_T))$ is a bipartite quiver such that $T$ is finite and $\operatorname{UnderG}(T)$ is a tree.
The map $\kappa:\operatorname{Edge}(T)\to[0,2]_{\Z}$ is a proper coloring.
\item\label{hyp:embedding-root} $v_0\in\operatorname{Vert}(T)$ and $\widehat v_0\in\operatorname{Vert}(\widetilde K_3)$ satisfy
\[
 v_0\in I_T\quad\Longleftrightarrow\quad\widehat v_0\in I_{\widetilde K_3}.
\]
\end{hypotheses}
Then, there is a unique bipartition-preserving embedding $\iota:T\to\widetilde K_3$ such that
\[
 \iota_0(v_0)=\widehat v_0,\qquad
 \operatorname{CovCol}_{\varpi}(\iota_1(\alpha))=\kappa(\alpha)
 \quad(\alpha\in\operatorname{Edge}(T)).
\]
\end{lemma}
\begin{proof}
Let us define the vertex and arrow maps.
Define $q:T\to K(3)$ by $q_0(v)=1$ for $v\in I_T$ and $q_0(v)=2$ for $v\in J_T$, where $1$ and $2$ denote the source and sink vertices of $K(3)$, respectively. Let
\[
 q_1(\alpha)=c_3^{-1}(\kappa(\alpha)).
\]
For $v\in\operatorname{Vert}(T)$, let $\gamma_v$ be the unique simple path from $v_0$ to $v$, viewed as an arrow traversal.
Its projection $q\gamma_v$ starts at $\varpi_0(\widehat v_0)$ by assumption~\ref{hyp:embedding-root}.
Define $\iota_0(v)$ to be the final vertex of $\operatorname{Lift}_{\varpi}(q\gamma_v;\widehat v_0)$.
For $\alpha:u\to v$, define $\iota_1(\alpha)$ to be the unique lift of $q_1(\alpha)$ starting at $\iota_0(u)$.
The simple paths to adjacent vertices differ by one edge.
Uniqueness of their lifts gives
\[
 s_{\widetilde K_3}(\iota_1(\alpha))=\iota_0(u),\qquad
 t_{\widetilde K_3}(\iota_1(\alpha))=\iota_0(v).
\]
The equality $\varpi_1(\iota_1(\alpha))=q_1(\alpha)$ gives $\operatorname{CovCol}_{\varpi}(\iota_1(\alpha))=\kappa(\alpha)$.

Let us prove injectivity.
Let $u\in\operatorname{Vert}(T)$ and $v\in\operatorname{Vert}(T)$ satisfy $u\ne v$.
Let
\[
 (z_0,\bar\alpha_1,z_1,\ldots,\bar\alpha_n,z_n),\qquad
 z_0=u,\quad z_n=v,\quad n\ge1
\]
be the unique simple path from $u$ to $v$.
For $i\in[0,n]_{\Z}$, let $w_i=\iota_0(z_i)$.
For $i\in[n]_{\Z}$, let $\widehat\alpha_i=\iota_1(\alpha_i)$.
Assumption~\ref{hyp:embedding-tree} and the color identity give
\[
 \operatorname{CovCol}_{\varpi}(\widehat\alpha_i)
 \ne\operatorname{CovCol}_{\varpi}(\widehat\alpha_{i+1})
 \qquad(i\in[n-1]_{\Z}).
\]
Let $H=\operatorname{UnderG}(\widetilde K_3)$ and
\[
 e_i=\llbracket w_{i-1},w_i\rrbracket\in\operatorname{Edge}(H)
 \qquad(i\in[n]_{\Z}).
\]
The distinct colors give $\widehat\alpha_i\ne\widehat\alpha_{i+1}$ for $i\in[n-1]_{\Z}$.
Since $H$ has no parallel edges, $e_i\ne e_{i+1}$ for $i\in[n-1]_{\Z}$.
Lemma~\ref{lem:tree-traversal-distinct} gives
\[
 \iota_0(u)=w_0\ne w_n=\iota_0(v).
\]
If $\iota_1(\alpha)=\iota_1(\beta)$, their two endpoints agree by injectivity of $\iota_0$.
The tree $\operatorname{UnderG}(T)$ has no parallel edges by assumption~\ref{hyp:embedding-tree}.
Hence, $\alpha=\beta$.

Let us prove uniqueness.
For any embedding $\jmath:T\to\widetilde K_3$ satisfying $\jmath_0(v_0)=\widehat v_0$ and $\operatorname{CovCol}_{\varpi}\circ\jmath_1=\kappa$, the traversal $\jmath(\gamma_v)$ is the lift of $q\gamma_v$ from $\widehat v_0$.
Lemma~\ref{lem:path-lift} determines this lift and hence every vertex and arrow image.
\end{proof}

\Needspace{8\baselineskip}
\begin{lemma}[Deck transformations with a fixed vertex image]\label{lem:cover-deck-transitivity}
Consider the universal covering $\varpi:\widetilde K_3\to K(3)$.
Assume the following.
\begin{hypotheses}
\item\label{hyp:deck-vertices} $u\in\operatorname{Vert}(\widetilde K_3)$ and $v\in\operatorname{Vert}(\widetilde K_3)$.
\item\label{hyp:deck-fiber} $\varpi_0(u)=\varpi_0(v)$.
\end{hypotheses}
Then, there is a unique $\delta\in\operatorname{Deck}(\varpi)$ such that $\delta_0(u)=v$.
\end{lemma}
\begin{proof}
Let us construct $\delta$ by lifting paths.
For $w\in\operatorname{Vert}(\widetilde K_3)$, let $\gamma_w$ be the unique simple path from $u$ to $w$.
By assumption~\ref{hyp:deck-fiber}, the vertex $v$ lies over the initial vertex of $\varpi\gamma_w$. Lemma~\ref{lem:path-lift} gives its unique lift starting at $v$.
Define $\delta_0(w)$ to be its final vertex.
For an arrow $\alpha:w\to w'$, define $\delta_1(\alpha)$ to be the unique lift of $\varpi_1(\alpha)$ starting at $\delta_0(w)$.
Uniqueness in Lemma~\ref{lem:path-lift} gives its endpoint $\delta_0(w')$ and
\[
 \varpi_0\circ\delta_0=\varpi_0,\qquad
 \varpi_1\circ\delta_1=\varpi_1.
\]
Interchanging $u$ and $v$ gives a quiver homomorphism $\epsilon$ over $K(3)$.
The composites $\epsilon\circ\delta$ and $\delta\circ\epsilon$ fix $u$ and $v$, respectively.
Applied to the paths from these vertices, uniqueness of lifting gives
\[
 \epsilon\circ\delta=1_{\widetilde K_3},\qquad
 \delta\circ\epsilon=1_{\widetilde K_3}.
\]
Hence, $\delta\in\operatorname{Deck}(\varpi)$.
If $\delta'\in\operatorname{Deck}(\varpi)$ satisfies $\delta'_0(u)=v$, then $\delta'\gamma_w$ and $\delta\gamma_w$ both lift $\varpi\gamma_w$ from $v$ for every $w\in\operatorname{Vert}(\widetilde K_3)$.
Lemma~\ref{lem:path-lift} gives $\delta'_0(w)=\delta_0(w)$ for every $w$ and $\delta'_1(\alpha)=\delta_1(\alpha)$ for every arrow $\alpha$.
Thus, $\delta'=\delta$.
\end{proof}

\begin{lemma}[Dimension vectors from colored trees]\label{lem:cover-dimension-equivalence}
Consider the universal covering
\[
 \varpi:\widetilde K_3\longrightarrow K(3)
\]
with arrow coloring $\operatorname{CovCol}_{\varpi}=c_3\circ\varpi_1$.
Assume the following.
\begin{hypotheses}
\item\label{hyp:dimensioned-trees} $((T,\kappa),(I_T,J_T))$ and $((T',\kappa'),(I_{T'},J_{T'}))$ are properly colored bipartite quivers such that $T$ and $T'$ are finite and $\operatorname{UnderG}(T)$ and $\operatorname{UnderG}(T')$ are trees.
\item\label{hyp:dimensioned-embeddings} $\iota:T\hookrightarrow\widetilde K_3$ and $\iota':T'\hookrightarrow\widetilde K_3$ preserve the bipartitions and satisfy
\[
 \operatorname{CovCol}_{\varpi}\circ\iota_1=\kappa,\qquad
 \operatorname{CovCol}_{\varpi}\circ\iota'_1=\kappa'.
\]
\item\label{hyp:positive-tree-dimensions} $\rho:\operatorname{Vert}(T)\to\Z_{>0}$ and $\rho':\operatorname{Vert}(T')\to\Z_{>0}$.
\end{hypotheses}
Define finitely supported dimension vectors $d$ and $d'$ by
\[
\begin{aligned}
 d_{\iota_0(v)}&=\rho(v)&& (v\in\operatorname{Vert}(T)),&
 d_w&=0&& (w\notin\iota_0(\operatorname{Vert}(T))),\\
 d'_{\iota'_0(v')}&=\rho'(v')&& (v'\in\operatorname{Vert}(T')),&
 d'_w&=0&& (w\notin\iota'_0(\operatorname{Vert}(T'))).
\end{aligned}
\]
Then, $[d]_{\varpi}=[d']_{\varpi}$ if and only if there is a bipartition-preserving quiver isomorphism $\psi:T\to T'$ satisfying
\[
 \kappa'\circ\psi_1=\kappa,\qquad
 \rho'\circ\psi_0=\rho.
\]
\end{lemma}
\begin{proof}
Let us restrict a deck transformation to the positive supports.
Suppose $d'=\delta\cdot d$ for $\delta\in\operatorname{Deck}(\varpi)$.
Assumption~\ref{hyp:positive-tree-dimensions} gives
\[
 \operatorname{VecSupp}(d)=\iota_0(\operatorname{Vert}(T)),\qquad
 \operatorname{VecSupp}(d')=\iota'_0(\operatorname{Vert}(T')).
\]
The images of $\iota$ and $\iota'$ are full subquivers of $\widetilde K_3$: an additional edge between image vertices would form a cycle together with their path in the image tree.
Thus, $\delta$ restricts to an isomorphism from the image of $\iota$ to the image of $\iota'$.
The maps $\iota$ and $\iota'$ in assumption~\ref{hyp:dimensioned-embeddings} give an isomorphism $\psi:T\to T'$ such that
\[
 \iota'\circ\psi=\delta\circ\iota.
\]
Deck transformations preserve $\operatorname{CovCol}_{\varpi}$, $I_{\widetilde K_3}$, and $J_{\widetilde K_3}$.
The equality $d'=\delta\cdot d$ gives $\rho'(\psi_0(v))=\rho(v)$ for every $v$.

Let us prove the converse by extending the colored tree isomorphism to the covering.
Conversely, let $\psi:T\to T'$ be a bipartition-preserving quiver isomorphism satisfying $\kappa'\circ\psi_1=\kappa$ and $\rho'\circ\psi_0=\rho$.
By assumption~\ref{hyp:dimensioned-trees}, choose $v_0\in\operatorname{Vert}(T)$.
Bipartition preservation gives
\[
 \varpi_0(\iota_0(v_0))=\varpi_0(\iota'_0(\psi_0(v_0))).
\]
Lemma~\ref{lem:cover-deck-transitivity} gives $\delta\in\operatorname{Deck}(\varpi)$ satisfying $\delta_0(\iota_0(v_0))=\iota'_0(\psi_0(v_0))$.
Their arrow colors satisfy
\[
 \operatorname{CovCol}_{\varpi}\circ\delta_1\circ\iota_1
 =\kappa=\kappa'\circ\psi_1
 =\operatorname{CovCol}_{\varpi}\circ\iota'_1\circ\psi_1.
\]
Uniqueness in Lemma~\ref{lem:cover-tree-embedding} gives $\delta\circ\iota=\iota'\circ\psi$.
The identity $\rho'\circ\psi_0=\rho$ then gives $d'=\delta\cdot d$.
\end{proof}

\begin{lemma}[Finiteness of connected support classes]\label{lem:cover-trees-finite}
Consider the universal covering $\varpi:\widetilde K_3\to K(3)$.
Let $a\in\Z_{>0}$ and $b\in\Z_{>0}$.
There are finitely many classes in
\[
 \left\{[d]_{\varpi}:d\in\operatorname{CovDim}_{\varpi}(a,b),\
 \operatorname{SubQuiv}(\widetilde K_3,\operatorname{VecSupp}(d))
 \text{ is connected}\right\}.
\]
\end{lemma}
\begin{proof}
Let us bound each connected support after a deck transformation.
Let $d\in\operatorname{CovDim}_{\varpi}(a,b)$ and assume its positive support $D=\operatorname{VecSupp}(d)$ is connected.
Since every $d_v$ for $v\in D$ is a positive integer,
\[
 \#D\le\sum_{v\in D}d_v=a+b.
\]
The inequality $a>0$ gives a source in $D$.
Fix a source $o$ of $\widetilde K_3$.
Lemma~\ref{lem:cover-deck-transitivity} gives a representative of $[d]_{\varpi}$ with $o\in D$.
For a vertex $v$, define $\operatorname{dist}(o,v)$ as the number of edges in the unique path from $o$ to $v$.
Since $D$ is connected, its path from $o$ to each $v\in D$ has at most $\#D-1$ edges. Hence,
\[
 D\subseteq\{v:\operatorname{dist}(o,v)\le a+b-1\},\qquad
 d_v\in[a+b]_{\Z}\quad(v\in D).
\]
Every vertex of $\widetilde K_3$ has degree three.
Thus, $\{v\in\operatorname{Vert}(\widetilde K_3):\operatorname{dist}(o,v)\le a+b-1\}$ is finite.
There are finitely many sets $D\subseteq\{v:\operatorname{dist}(o,v)\le a+b-1\}$ and finitely many functions $D\to[a+b]_{\Z}$.
These finite choices contain a representative of every class in the statement.
\end{proof}

\begin{lemma}[Cover stability and localization]\label{lem:cover-localization}
Let $a$ and $k$ be integers. Assume the following.
\begin{hypotheses}
\item\label{hyp:localization-dimensions} $k\ge2$ and $a\ge100000k$.
\item\label{hyp:localization-coprimality} $\gcd(a,k)=1$.
\end{hypotheses}
Fix a family of representatives $\rho=(\rho_M)_{M\in\mathcal H_{a,k}}$.
Let $I=I_{\widetilde K_3}$ and $J=J_{\widetilde K_3}$.
Then, the following hold.
\begin{enumerate}
\item\label{c:cover-stability-weight} Let $R$ be a representation of dimension $\widetilde d\in\operatorname{CovDim}_{\varpi}(a,a+k)$.
For $W\in\operatorname{SubRep}(R)$, let
\[
 u=\sum_{i\in I}\dim W_i,\qquad v=\sum_{j\in J}\dim W_j.
\]
Then,
\[
 \theta^{\varpi}_{a,a+k}(\dim R)=0,\qquad
 \theta^{\varpi}_{a,a+k}(\dim W)=av-(a+k)u.
\]
Thus, the stability condition~\eqref{eq:stability} with $b=a+k$ is equivalent to
$\theta^{\varpi}_{a,a+k}(\dim W)>0$ for every proper nonzero subrepresentation $W$.
\item\label{c:cover-stability-connected} Let $R$ be a representation of dimension $\widetilde d\in\operatorname{CovDim}_{\varpi}(a,a+k)$.
If $R$ satisfies the stability condition in Claim~\ref{c:cover-stability-weight},
then the quiver $\operatorname{SubQuiv}(\widetilde K_3,\operatorname{VecSupp}(\widetilde d))$ is connected.
\item\label{c:cover-localization-bound} We have
\[
 E_k(a)\ge\#\mathcal S_{a,k}^\rho.
\]
\end{enumerate}
\end{lemma}
\begin{proof}
\prf{Claim~\ref{c:cover-stability-weight}} Let us check the dimension-vector and stability conditions for the support.
Definition~\ref{def:cover-dimension-classes}(1)--(2) gives
\[
 \theta^{\varpi}_{a,a+k}(\dim R)=a(a+k)-(a+k)a=0,
 \qquad\theta^{\varpi}_{a,a+k}(\dim W)=av-(a+k)u.
\]
Substituting $b=a+k$ in~\eqref{eq:stability} proves the equivalence.

\prf{Claim~\ref{c:cover-stability-connected}} Let us prove connectedness of the support by contradiction.
Suppose that $\operatorname{VecSupp}(\widetilde d)$ is disconnected.
Then, $R=R_1\oplus R_2$ for two nonzero subrepresentations.
Claim~\ref{c:cover-stability-weight} gives
\[
 0<\theta^{\varpi}_{a,a+k}(\dim R_1)+\theta^{\varpi}_{a,a+k}(\dim R_2)
 =\theta^{\varpi}_{a,a+k}(\dim R)=0,
\]
a contradiction.

\prf{Claim~\ref{c:cover-localization-bound}} Let us obtain a sum formula for $E_k(a)$ over cover dimension vectors.

Let us transfer stability to each chosen representative.
We use the stability condition in Claim~\ref{c:cover-stability-weight}
for $\cM^s_{\widetilde d}(\widetilde K_3)$.
On $\widetilde K_3$, the equivalence classes of dimension vectors in~\cite[Corollary 3.16]{Weist} are the $\operatorname{Deck}(\varpi)$-orbits.
By~\cite[Corollary 3.16]{Weist}, we have
\begin{equation}
 E_k(a)=\chi\bigl(\cM^s_{(a,a+k)}(K(3))\bigr)
 =\sum_{[\widetilde d]_{\varpi}\in\mathcal D_{a,k}}
 \chi\bigl(\cM^s_{\widetilde d}(\widetilde K_3)\bigr).
\label{eq:cover-sum}
\end{equation}

Let us restrict the sum in~\eqref{eq:cover-sum} to finitely many nonempty moduli spaces.
Claim~\ref{c:cover-stability-connected} gives connected positive support for each nonempty summand.
Lemma~\ref{lem:cover-trees-finite}
with $b=a+k$ gives only finitely many such classes.
For each such class $[\widetilde d]_{\varpi}$, let $S_{\widetilde d}=\operatorname{VecSupp}(\widetilde d)$.

Let us restrict stable representations to $S_{\widetilde d}$.
Let $\operatorname{SubQuiv}(\widetilde K_3,S_{\widetilde d})$ be the full subquiver on $S_{\widetilde d}$.
Claims~\ref{c:restriction-zero-left}--\ref{c:restriction-zero-right} of Lemma~\ref{lem:rep-restriction-zero} give inverse correspondences
\[
 [R]\longmapsto[\operatorname{Res}_{\widetilde K_3,S_{\widetilde d}}R],\qquad
 [N]\longmapsto[\operatorname{ExtZero}_{\widetilde K_3,S_{\widetilde d}}N].
\]
Claim~\ref{c:restriction-zero-subrep} of Lemma~\ref{lem:rep-restriction-zero} preserves subrepresentations and their dimensions at every vertex of $S_{\widetilde d}$, and hence their $\theta^{\varpi}_{a,a+k}$-weights.
The restriction and zero-extension formulas in Definition~\ref{def:rep-restriction-zero} apply to families of representations and induce
\begin{equation}\label{eq:cover-finite-support-moduli}
 \cM^s_{\widetilde d}(\widetilde K_3)
 \cong\cM^s_{\widetilde d|_{S_{\widetilde d}}}(\operatorname{SubQuiv}(\widetilde K_3,S_{\widetilde d})).
\end{equation}
The quiver $\operatorname{SubQuiv}(\widetilde K_3,S_{\widetilde d})$ is finite and acyclic.

Let us identify the classes in $\mathcal S_{a,k}^\rho$ with distinct cover dimension classes $[d]_{\varpi}$.
Let $M\in\mathcal H_{a,k}$.
Claim~\ref{c:flow-choices-stable} of Lemma~\ref{lem:flow-choices} gives stability of $R_M^\rho$.
Claims~\ref{c:flow-stability-support-tree}, \ref{c:flow-stability-support-colors}, and~\ref{c:flow-stability-total-dimension} of Lemma~\ref{lem:flow-stability} give a properly colored tree support and total dimension $(a,a+k)$.
Lemma~\ref{lem:cover-tree-embedding} gives a color-preserving embedding
\[
 \iota:\widetilde Q_M^\rho\longrightarrow\widetilde K_3.
\]

Let us extend the embedded representation by zero.
Let $U_M=\iota_0(\operatorname{Vert}(\widetilde Q_M^\rho))$.
An additional arrow of $\widetilde K_3$ between vertices of $U_M$ would form a cycle with the path between its endpoints in the embedded tree.
Thus, the image is the full subquiver $\operatorname{SubQuiv}(\widetilde K_3,U_M)$.
Transport $R_M^\rho$ to a representation $N_M$ of this subquiver and let
\[
 \widehat R_M^\rho=\operatorname{ExtZero}_{\widetilde K_3,U_M}N_M,\qquad
 \widehat d_M^\rho(v)=\dim_{\C}(\widehat R_M^\rho)_v
 \quad(v\in\operatorname{Vert}(\widetilde K_3)).
\]
Claim~\ref{c:restriction-zero-subrep} of Lemma~\ref{lem:rep-restriction-zero} preserves subrepresentations and their dimension vectors.
Hence, $\widehat R_M^\rho$ is stable and $[\widehat d_M^\rho]_{\varpi}\in\mathcal D^{\mathrm{st}}_{a,k}$.

Let us pass to dimension classes.
Lemma~\ref{lem:cover-dimension-equivalence} gives a well-defined injective map
\[
 \mathcal S_{a,k}^\rho\longrightarrow\mathcal D^{\mathrm{st}}_{a,k},\qquad
 [(\widetilde{\mathbb Q}_M^\rho,d_M^\rho)]\longmapsto[\widehat d_M^\rho]_{\varpi}.
\]

Let us prove that a subdimension vector of weight zero is zero or the full dimension vector.
For $\widetilde d\in\operatorname{CovDim}_{\varpi}(a,a+k)$, let $0\le e\le\widetilde d$ be a subdimension vector.
Let
\[
 u=\sum_{i\in I}e_i,\qquad v=\sum_{j\in J}e_j,
 \qquad u\in[0,a]_{\Z},\quad v\in[0,a+k]_{\Z}.
\]
By assumption~\ref{hyp:localization-coprimality}, we have $\gcd(a,a+k)=\gcd(a,k)=1$.
Definition~\ref{def:cover-dimension-classes}(2) gives
\[
 \theta^{\varpi}_{a,a+k}(e)=0
 \ \Longrightarrow\ av=(a+k)u
 \ \Longrightarrow\ (u,v)\in\{(0,0),(a,a+k)\}.
\]
If $(u,v)=(0,0)$, nonnegativity gives $e=0$.
If $(u,v)=(a,a+k)$, then
\[
 \sum_x(\widetilde d_x-e_x)=0,
 \qquad \widetilde d_x-e_x\ge0\quad(x\in\operatorname{Vert}(\widetilde K_3)).
\]
Hence, $e=\widetilde d$.
Thus, no proper nonzero subdimension vector has weight zero.

Let us bound the Euler characteristic of each nonempty summand from below.
Let $X=\cM^s_{\widetilde d}(\widetilde K_3)$.
Suppose that $X\ne\varnothing$.
Every common divisor of the entries of $\widetilde d|_{S_{\widetilde d}}$ divides $a$ and $a+k$.
Thus, $\gcd\{\widetilde d_x:x\in S_{\widetilde d}\}=1$.
Since $\theta^{\varpi}_{a,a+k}(e)\ne0$ for every subdimension vector $e$ satisfying $0\le e\le\widetilde d$ and $e\notin\{0,\widetilde d\}$, $\theta^{\varpi}_{a,a+k}$-semistability and $\theta^{\varpi}_{a,a+k}$-stability coincide for dimension vector $\widetilde d$.
By~\eqref{eq:cover-finite-support-moduli} and~\cite[Propositions 4.3 and 5.3, Remark 5.4]{King},
$X$ is a smooth projective fine moduli space for the finite acyclic support quiver.

Let us compute $\chi(X)$.
King--Walter's theorem~\cite[Theorem 3(c)]{KingWalter} gives
\[
 H^{2j+1}(X;\Q)=0\qquad(j\in\Z_{\ge0}).
\]
Hence,
\[
 \chi(X)=\sum_{j=0}^{\dim_{\C}X}\dim_{\Q}H^{2j}(X;\Q)
 \ge\dim_{\Q}H^0(X;\Q)\ge1.
\]
This proves
\begin{equation}
 \cM^s_{\widetilde d}(\widetilde K_3)\ne\varnothing
 \quad\Longrightarrow\quad
 \chi\bigl(\cM^s_{\widetilde d}(\widetilde K_3)\bigr)\ge1.
\label{eq:positive-summands}
\end{equation}
Therefore, equations~\eqref{eq:cover-sum} and~\eqref{eq:positive-summands} give
\[
 E_k(a)\ge
 \#\mathcal D^{\mathrm{st}}_{a,k},
\]
and $\#\mathcal D^{\mathrm{st}}_{a,k}\ge\#\mathcal S_{a,k}^\rho$. Proposition~\ref{prop:cover-count} bounds $\#\mathcal S_{a,k}^\rho$ from below.
\end{proof}

\Needspace{6\baselineskip}
\begin{theorem}\label{thm:central-lower}
We have the following.
\begin{enumerate}
\item\label{c:central-lower-finite} There are constants $c_*\in(0,1]$ and $q_*\in(0,1]$ with the following property.
For each $a\in\Z$ and $k\in\Z$ satisfying
\begin{hypotheses}
\item\label{hyp:lower-dimensions} $k\ge2$ and $a\ge100000k$,
\item\label{hyp:lower-coprimality} $\gcd(a,k)=1$,
\end{hypotheses}
we have
\begin{equation}\label{finite}
 E_k(a)\ge c_*\beta^a\left(\frac{k}{a}\right)^{\frac{5}{2}}q_*^{k-1}.
\end{equation}
\item\label{c:central-lower-limit} Let $((a_n,k_n))_{n\ge1}$ be a sequence in $\Z_{>0}^2$ satisfying
\[
 \begin{gathered}
 k_n\ge2,\qquad \gcd(a_n,k_n)=1\qquad(n\ge1),\\
 \lim_{n\to\infty}a_n=\infty,\qquad
 \lim_{n\to\infty}\frac{k_n}{a_n}=0.
 \end{gathered}
\]
Then, $E_{k_n}(a_n)>0$ for all sufficiently large $n$, and
\begin{equation}\label{entropy}
 \liminf_{n\to\infty}\frac{\log E_{k_n}(a_n)}{a_n}\ge\log\beta.
\end{equation}
\end{enumerate}
\end{theorem}
\begin{proof}
\prf{Claim~\ref{c:central-lower-finite}} Let us obtain the finite lower bound with constants independent of the dimensions.
For each $a\in\Z_{>0}$ and $k\in\Z_{>0}$ satisfying $k\ge2$ and $a\ge100000k$, fix a family of representatives $\rho=(\rho_M)_{M\in\mathcal H_{a,k}}$.

Let us apply the counting and localization bounds.
Choose a $(P_j,Q_j^{=},256)$-admissible pair $(A,B)$ by Lemma~\ref{lem:slot-bounds}.
The inequality $a\ge100000k$ verifies the dimension hypothesis of Proposition~\ref{prop:cover-count}.
The additional equality $\gcd(a,k)=1$ verifies the coprimality hypothesis of
Claim~\ref{c:cover-localization-bound} of Lemma~\ref{lem:cover-localization}.
Set $N=256$, $\lambda=\frac1{100}$, $x=\frac ak$, and $\gamma=1-\frac Nx$.
These results give
\[
 E_k(a)\ge\#\mathcal S_{a,k}^{\rho}
 \ge\beta^a x^{-\frac52}\frac{A^2}{8B\lambda N^2}
 \left(\frac{6AB\lambda\gamma^2}{\beta}\right)^{k-1}.
\]
Let
\[
 \gamma_0=1-\frac{N}{100000},\qquad
 c_*=\min\left\{1,\frac{A^2}{8B\lambda N^2}\right\},\qquad
 q_*=\min\left\{1,\frac{6AB\lambda\gamma_0^2}{\beta}\right\}.
\]
These constants are positive. Assumption~\ref{hyp:lower-dimensions}
gives $x=\frac ak\ge100000$. Hence, $\gamma\ge\gamma_0$. Substituting $x=\frac{a}{k}$ and $c=k-1$ in~\eqref{count}
proves~\eqref{finite}.

\prf{Claim~\ref{c:central-lower-limit}} Let us apply the finite bound to $(a_n,k_n)$.
Since $\lim_{n\to\infty}\frac{k_n}{a_n}=0$, choose $n_0\in\Z_{>0}$ such that
\[
 a_n\ge100000k_n\qquad(n\ge n_0).
\]
Together with $k_n\ge2$ and $\gcd(a_n,k_n)=1$, this verifies the hypotheses of Claim~\ref{c:central-lower-finite}.
Taking logarithms in~\eqref{finite} and dividing by $a_n$ gives
\begin{equation}\label{eq:central-lower-log-bound}
 \frac{\log E_{k_n}(a_n)}{a_n}\ge\log\beta+
 \frac{\log c_*-\frac52\log\left(\frac{a_n}{k_n}\right)+(k_n-1)\log q_*}{a_n}
 \qquad(n\ge n_0).
\end{equation}

Let us prove that each error term tends to zero as $n\to\infty$.
The constants $c_*$ and $q_*$ do not depend on $n$.
For $n\ge n_0$, we have
\[
 \begin{aligned}
 \lim_{n\to\infty}\frac{\log c_*}{a_n}&=0,\\
 0\le\frac{\log(a_n/k_n)}{a_n}&\le\frac{\log a_n}{a_n}\longrightarrow0,\\
 \left|\frac{(k_n-1)\log q_*}{a_n}\right|
 &\le\frac{k_n}{a_n}|\log q_*|\longrightarrow0.
 \end{aligned}
\]
Taking the lower limit as $n\to\infty$ in~\eqref{eq:central-lower-log-bound} proves~\eqref{entropy}.
\end{proof}

